\sectiontitle{Section III:}

                                                       Section III:
                                             Sparking and Slipping




Pattern, Poetry, and Power in the Music of Frédéric Chopin      171

                                                             Section 111:
                                                   Sparking and Slipping
        The concern of the following five chapters is creativity: its wellsprings and its
mechanizability. One of the most common metaphors for creativity is that of "spark": an
electric leap of thought from one place to a remote one, without any apparent justification
beforehand, but with all the justification in the world after the fact. Besides being used as
a noun, "spark" is also used as a verb: one idea sparks another. Creative mental activity
becomes, in this imagery, a set of sparks flying around in a space of concepts. Just how
different is this metaphor for the mind from the reality of computers? They are filled with
electricity rushing from one place to another at the most unimaginable speeds. Isn't that
enough to turn the mechanical into the fluid? Or do computers still lack something
ineffable? Are their mechanical attempts at thinking still too rigid, too dry? Is something
liquid and slippery missing? My word for the elusive aspect of human thought still
lacking in synthetic imitations is "slippability". Human thoughts have a way of slipping
easily along certain conceptual dimensions into other thoughts, and resisting such
slippage along other dimensions. A given idea has slightly different slippabilities-
predispositions to slip-in each different human mind that, it comes to live in. Yet some
minds' slippabilities seem to give rise to what we consider genuine creativity, while
others' do not. What is this precious gift? Is there a formula to the creative act? Can spark
and slippability be canned and bottled? In fact, isn't that just what a human brain is-an
encapsulated creativity machine? Or is there more to creativity and mind than can ever be
encapsulated in any finite physical object or mathematical model?




Pattern, Poetry, and Power in the Music of Frédéric Chopin                               172

                      Pattern, Poetry, and Power
                   in the Music of Frédéric Chopin
                                       April, 1982



THE abstract visual pattern in Figure 9-1 is a graphical representation of the opening
of one of the most difficult and lyrical pieces for piano eve: composed, namely the
eleventh etude in Frederic Chopin's Opus 25, writer in about 1832, when he was in his
early twenties. As a boy, I heard the Chopin etudes many times over on my parents'
phonograph, and I quickly grew to love them. They became as familiar to me as the faces
of my friends. Indeed, I cannot imagine who I would be if I did not know these pieces.
          A few years later, as a teen-ager who enjoyed playing piano, I wanted to learn to
play some of these old friends. I went to the local music store and found a complete
volume of them. I will never forget my reaction on opening the book and looking for my
friends. They were nowhere to be found! I saw nothing but masses of black notes and
chords: complex, awesome visual patterns that I had never imagined. It was as if,
expecting to meet old friends, I had instead found their skeletons grinning at me. It was
terrifying. I closed the book and left, somewhat in shock.
          I remember going back several times to that music store, each time pulled by the
same curiosity tinged with fear. One day I worked up my courage and actually bought
that book of etudes. I suppose I hoped that if I simply sat down at the piano and tried
playing the notes I saw, I would hear my old friends, albeit a little slowly. Unfortunately,
nothing of the kind happened. 11, general, I could not even play the two hands together
comfortably, let alone recreate the sounds I knew so well. This left me disheartened and a
little frightened at the realization of the awesome complexities I had taken for granted.
You can look at it two ways. One way is to be amazed at how human perception can
integrate a huge set of independent elements and "hear" only a single quality: the other is
to be amazed at the incredible skill of a pianist who can play so many notes so quickly
that they all blur into one shimmering mass, a "co-hear-ent" totality.
          At first it was bewildering to see that “friends” had anatomies of such




Pattern, Poetry, and Power in the Music of Frédéric Chopin                              173

overwhelming complexity. But looking back, I don't know what I expected. Did I expect
that a few simple chords could work the magic that I felt? No; if I had thought it over, I
would have realized this was impossible. The only possible source of that magic was in
some kind of complexity-patterned complexity, to be sure. And I think this experience
taught me a lifelong lesson: that phenomena perceived to be magical are always the
outcome of complex patterns of nonmagical activities taking place at a level below
perception. More succinctly: The magic behind magic is pattern. The magic of life itself
is a perfect example, emerging as it does out of patterned but lifeless activities at the
molecular level. The magic of music emerges from complex, nonmagical-or should I say
metamagical?-patterns of notes.

                                         *    *    *

         Having bought this volume, I felt drawn to it, wanted to explore it somehow. I
decided that, hard work though it might be, I would learn an etude. I chose the one that
was my current favorite-the one pictured in Figure 9-1-and set about memorizing the
finger pattern in the right hand, together with the patterns that follow it, making up the
first two pages or so. I played the pattern literally thousands of times, and gradually it
became natural to my fingers, although never as natural as it had always sounded to my
ears-or rather, to my mind.
         It was then that I first observed the amazing subtlety of the lightning flash of the
right hand, how it is composed of two alternating and utterly different components: the
odd-numbered notes (in red) trace out a perfect descending chromatic scale for four
octaves, while the even-numbered notes (in black), wedged between them like pickets
between the spaces in a picket fence, dictate an arpeggio with repeated notes. To execute
this alternating pattern, the right hand flutters down the keyboard, tilting from side to side
like a swift in flight, its wings beating alternately.
         A word of explanation. On a piano there are twelve notes (some black, some
white) from any note to the corresponding note one octave away. Playing them all in
order creates a chromatic scale, as contrasted with the more familiar diatonic scales
(usually major or minor). These latter involve only seven notes apiece (the eighth note
being the octave itself). The seven intervals between the successive notes of a diatonic
scale are not all equal. Some are twice as large as others, yet to the ear there is a perfect
intuitive logic to it. Rather paradoxically, in fact, most people can sing a major scale
without any trouble, uneven intervals notwithstanding, but few can sing a chromatic scale
accurately, even though it "ought" to be much more straightforward-or so it would seem,
since all its intervals are exactly the same size. The chromatic scale is so called because
the extra notes it introduces to fill up the gaps in a diatonic scale have a special kind of
"bite" or sharpness to them that adds color or piquancy to a piece. For that reason, a piece
filled with notes other than the seven notes belonging to the key it is in is said to be
chromatic.




Pattern, Poetry, and Power in the Music of Frédéric Chopin                                174

FIGURE 9-2. The strikingly different visual textures of six Chopin Etudes. On top, Op.
10, No. 11, in E-flat major; Op. 25, No. 1, in A flat major, and Op. 25, No. 2, in F minor.
Below, Op. 25, No. 3, in F major; Op. 25, No. 6, in G-sharp minor; and Op. 25, No. 12,
in C minor: [From the G. Schirmer (Friedheim) edition.]




Pattern, Poetry, and Power in the Music of Frédéric Chopin                             175

An arpeggio is a broken chord played one or more times in a row, moving up or down the
keyboard. Thus it bears a resemblance to a spread-out scale, a little like someone
bounding up a staircase three or four steps at a time. Chopin's music is filled with both
arpeggios and chromatic passages, but the intricate fusion of these two opposite structural
elements in the eleventh etude struck me as a masterpiece of ingenuity. And what is
amazing is how it is perceived when the piece moves quickly. The chromatic scale comes
through loud and clear, forming a smooth "envelope" of the pattern (your eye picks it out
too), but the arpeggio blurs into a kind of harmonic fog that deeply affects one's
perception, if only subliminally, or so it seems at least to the untrained ear.
        Each etude in that book I bought has a characteristic appearance, a visual texture
(see Figure 9-2). This was one of the most striking things about the book at first. I was
not at all accustomed to the idea of written music as texture; the simple pieces I had
played up to that time were slow, so that every note was distinctly heard. In other words,
the pieces in my playing experience were coarse-grained compared with the fine grain of
a Chopin etude, where notes often go by in a blur and are merely parts of an auditory
gestalt. Conversion of this kind of auditory experience to notated music sheets often
yields quite stunning textures and patterns. Each composer has a characteristic set of
patterns the eye becomes familiar with, and these etudes provided for me a stunning
realization of that fact.

                                            *     *     *

        Sadly, I was forced to abandon etude Op. 25, No. 11, after having learned only a
little more than a page-it was simply too hard for me. James Huneker, an American critic
and one of Chopin's earliest English-language biographers, wrote of this study: "Small-
souled men, no matter how agile their fingers, should not attempt it." Well, whatever the
size of my soul, my fingers were not agile enough. For a while, that discouraged me from
attacking any more Chopin etudes at all. A few years later, though, when I was working
more earnestly on improving my modest piano skills, I came across an isolated Chopin
etude in a book of medium-difficult selections. It turned out to be one of three etudes he
had composed later in life, none of which had been on my parents' records. This was a
real find! Luckily its texture looked less prickly, its pace less forbidding. Somewhat
gingerly, I played through it very slowly and discovered that it was astonishingly
beautiful and not as inaccessible as the others I'd tried.
        Like all the rest of Chopin's studies, this one is centered on a particular technical
point, although to think of the etudes primarily in that way is like thinking of the fantastic
gymnastic performances of Nadia Comaneci as merely fancy fitness exercises. Louis
Ehlert, a nineteenth-century musicologist, wrote of one of the most beautiful etudes in
Opus 25 (the sixth one, in G-sharp minor): "Chopin not only versifies an exercise in
thirds; he transforms it into such a work of art that in studying it one could sooner




Pattern, Poetry, and Power in the Music of Frédéric Chopin                                176

fancy oneself on Parnassus than at a lesson. He deprives every passage of all mechanical
appearance by promoting it to become the embodiment of a beautiful thought, which in
turn finds graceful expression in its motion." Similar words apply to this easier,
posthumously published etude in A-flat major, whose chief technical concern is the
concept of three against two, a special case of the general concept of polyrhythm.
        Mathematically, the concept is simple enough: play two musical lines
simultaneously, one of them sounding three notes to the other's two. Usually the triplet
and doublet are aligned so that they start at the same instant. When they are both plotted
on a unit interval (see Figure 9-3a), you can see that the doublet's second note is struck
halfway between the triplet's second and third notes. Of course, this is simply a pictorial
representation of the fact that 1/2 is the arithmetic mean of 1/3 and 2/3.
        In theory, two voices playing a three-against-two pattern need not be perfectly
aligned. If you shift the upper voice by, say, 1/12 to the right, you get a different picture
(see Figure 9-3b). Here the triplet's third note starts halfway through the doublet's second.
As you can see, the triplet extends beyond the end of the interval, presumably to join onto
another identical

FIGURE 9-3. The 3-against-2 phenomenon. In (a), as it is usually heard, with both voices
"in phase'' . In (b), one voice is shifted by 1112 with respect to the other, producing a
quite unusual pattern of beats. In (c), it is shown how in principle the relative staggering
of the two voices could be adjusted continuously by a knob arrangement.




Pattern, Poetry, and Power in the Music of Frédéric Chopin                               177

    pattern. We can fold the pattern around and represent its periodicity in a circle, as is
  shown in Figure 9-3c. By rotating either of the concentric circles like a knob, we get all
    possible ways of hearing three beats against two. In Chopin and most other Western
   music, however, the only possibility that I have seen explored is where the triplet and
                               doublet are perfectly "in phase".
At first I found the three-against-two rhythm hard to perform exactly. One has to learn
how to hear the voices separately, to hear the roundish lilt of the three-rhythm weaving
itself into the square mesh of the two-rhythm. Of course, it's easy to hear when someone
else is playing; the trick is to hear it in one's own playing! In principle the task is not
hard, but it is one of coordination, and requires practice. I found that once I had mastered
the problem of playing the two rhythms evenly and independently, I could play the whole
etude. To play it-or to hear it-is like smiling through tears, it is so beautiful and sad at the
same time.
It is impossible to pinpoint the source of the beauty, needless to say, but it is certainly due
in part to the way the chords in the right hand flow into one another. (See Figure 9-4.)
Almost all the way through the piece, the




FIGURE 9-4. The opening two measures of the posthumous etude in Afat major, showing
its typical 3-against-2 pattern with slowly shifting chords in the right hand. [Music
printed by Donald Byrd's SMUT program at Indiana University.]

right hand plays three-note chords (six to a measure) against single notes by the left hand
(four to a measure). The delicacy of the piece comes from the fact that very often, when
one chord flows into the next one, only a single note changes. And to add to the subtlety
of this slowly shifting sound-pattern, usually the steps taken by the shifting voice are
single scale-steps rather than wide jumps. These "rules" do not hold all the way, of
course; there are numerous exceptions. Nevertheless, there is a uniform aural texture to
the piece that imbues it with its soft melancholy, known in Polish as tgsknota.

                         *   *          *    *               *       *

It is interesting to speculate about the extent to which such formal considerations
occurred to Chopin while he was composing. It is well known that Chopin revered Bach's
music. "Always play Bach" was his advice to a




Pattern, Poetry, and Power in the Music of Frédéric Chopin                                 178

FIGURE 9-5. Chopin's Etude in C major from Opus 10, his first etude, computer-printer
as to reproduce as closely as possible the stunning usual pattern that Chopin himself
caret produced in his manuscript. Aside from the beautiful alignment of crests and
troughs, Chop manuscript features whole notes centered in their measures (in the bass
clef). [.'Music printer Donald Bird's SMUT program at Indiana University'.]




Pattern, Poetry, and Power in the Music of Frédéric Chopin                       179

pupil, and he was particularly devoted to the Well-Tempered Clavier, a paragon of
elegant formal structures. Chopin confided to his friend Eugene Delacroix, the painter,
that "The fugue is like pure logic in music ... To know the fugue deeply is to be
acquainted with the element of all reason and consistency in music." Clearly, Chopin
loved pattern.
A stunning demonstration of Chopin's extreme awareness of the visual appeal of the
textures in his etudes is provided by the appearance of the manuscript of his etude Op. 10,
No. 1, in C major, one about which James Huneker wrote, in his inimitable prose:

   The irregular black ascending and descending staircases of notes strike the
   neophyte with terror. Like Piranesi's marvellous aerial architectural dreams, these
   dizzy acclivities and descents of Chopin exercise a charm, hypnotic, if you will,
   for eye as well as ear. Here is the new technique in all its nakedness, new in the
   sense of figure, design, pattern, web, new in a harmonic way. The old order was
   horrified at the modulatory harshness; the young sprigs of the new, fascinated and
   a little frightened. A man who could thus explode a mine that assailed the stars
   must be reckoned with.

         That "terror-stricken neophyte" might well have been me. Huneker's words form
an amusing contrast with what the nineteen-year-old Chopin himself wrote of this, his
first etude, in a letter to his friend Tytus Woyciechowski in 1829: "I have written a large
exercise in form, in my own personal style; when we get together, I'll show it to you." A
finished copy, believed to be in Chopin's hand, is now in the Museum of the Frederic
Chopin Society in Warsaw. With the present turmoil in Poland, it would be difficult to
gain permission to reproduce it directly. Fortunately, a long-standing research project of
my friend Donald Byrd at Indiana University has been to develop a computer program
that can print out music according to specification, and at professional standards. With
some help from our friend Adrienne Gnidec, Don and I coaxed his marvelous program
into printing the music in a very strange and visually striking way (see Figure 9-5). This
figure reproduces quite accurately the large-scale visual patterns of Chopin's own
manuscript, in which Chopin took great care to align all the crests of the massive waves.
When this piece is played at the proper speed, each sweep up and down the keyboard is
heard as one powerful surge, like the stroke of an eagle's wing, with the notes of each
crest sparkling brilliantly like wingtips flashing in the sun.
Another interesting feature of Chopin's notation, here copied, is his positioning of the
doubled whole notes in the bass. Instead of placing them at the very start of each
measure, aligned with the sixteenth-note rests, Chopin centered each one in its own
measure, thereby creating an elegant visual balance, though losing some notational
clarity. Musically, such centering has no effect. Since a whole note lasts for the duration
of an entire 4/4 measure, it must be struck at the start of the measure, otherwise it would
overflow into the next measure, and that is impossible. (Or rather, it would




Pattern, Poetry, and Power in the Music of Frédéric Chopin                             180

Pattern, Poetry, and Power in the Music of Frédéric Chopin   181

Pattern, Poetry, and Power in the Music of Frédéric Chopin   182

 violate a much more rigid convention of music notation-namely, that n note can
designate a sound that overflows the boundaries of its measure. Hence the only possible
interpretation is that the whole note is to be struck at the outset. In other words, the
centering is simply a charming artistic touch with a quaint nineteenth-century flavor, like
the ornaments on Victorian house. The modern music-reading eye is used to more
function notation; in particular, it expects the staff to be in essence a graph of the sound,
in which the horizontal axis is time. Thus notes struck simultaneously are expected to line
up vertically.
         But let us return to the matter of Chopin's preoccupation with form an structure.
Few composers of the romantic era have penned such visual] patterned pages, have spun
a whole cloth out of a single textural idea. Wit Chopin, though, preoccupation with strict
pattern never took precedent over the expression of heartfelt emotions. One must
distinguish, it seers to me, between "head pattern" and "heart pattern", or, in moi
objective-sounding terms, between syntactic pattern and semantic patters The notion of a
syntactic pattern in music corresponds to the form; structural devices used in poetry:
alliteration, rhyme, meter, repetition < sounds, and so on. The notion of a semantic
pattern is analogous to the pattern or logic that underlies a poem and gives it reason to
exist: the inspiration, in short.
         That there are such semantic patterns in music is as undeniable as the; there are
courses in the theory of harmony. Yet harmony theory has no more succeeded in
explaining such patterns than any set of rules has yet succeeded in capturing the essence
of artistic creativity. To be sure, they are words to describe well-formed patterns and
progressions, but no theory yet invented has even come close to creating a semantic sieve
so fine as t let all bad compositions fall through and to retain all good ones. Theories of
musical quality are still descriptive and not generative; to some extent they can explain in
hindsight why a piece seems good, but they are not sufficient to allow someone to create
new pieces of quality and interest. is nonetheless fascinating, if not downright
compelling, to try to find certain earmarks of greatness, to try to understand why it is that
one composer music can reach in and touch your innermost core while another composer
music leaves you cold and unmoved. It is a mystery.

                                           *     *     *

After learning the posthumous A-flat etude, I felt encouraged to tack some of the others.
One of the ones I had loved the most was Op. 25, N 2, in F minor. To me, it was a soft,
rushing whisper of notes, a fluttering like the leaves of a quaking aspen in a gentle
breeze. Yet it was not just scene of nature; it expressed a human longing, a melancholy
infused with strange and wild yearnings for something unknown and remote-tgsknota
again. I knew this melody inside out from many years of hearing it, and looked forward
to transferring it to my fingers.




Pattern, Poetry, and Power in the Music of Frédéric Chopin                               183

After a couple of months' practice, my fingers had built up enough stamina to play the
piece fairly evenly and softly. This was very satisfying to me until one day, an
acquaintance for whom I was playing it commented, "But you're playing it in twos-it's
supposed to be in threes!" What she meant by this was that I was stressing every second
note, rather than every third. Bewildered, I looked at the score, and of course, as she had
pointed out, the melody was written in triplets. But surely Chopin had not meant it to be
played in threes. After all, I knew the melody perfectly! Or did I? I tried playing it in
threes. It sounded strange and unfamiliar, a perceptual distortion the like of which I had
never experienced.
        I went home and took out my parents' old Remington LP of the Chopin etudes
Opus 25 (played by a wonderful but hardly remembered pianist named Alexander
Jenner). I put on the F minor etude and tried to hear which way he played it. I found I
could hear it either way. Jenner had played it so smoothly, so free of accent (as they say
Chopin did, by the way), that one really could not tell which way to hear it. All of a
sudden I saw that I really knew two melodies composed of the exact same sequence of
notes! I felt myself to be very fortunate, because now I could experience this familiar old
melody in a fresh new way. It was like falling in love with the same person twice.
        I had to practice hard to undo the bad habits of "biplicity" and to replace them
with the indicated "triplicity", but it was a delight. The hardest part, however, was
combining the two hands. With duplets in the right hand, this had presented no problem;
all the accented notes fell in coincidence with notes in the left hand, moving at exactly
half the speed of the right hand in a pattern of wide arpeggios. But if I were to spread my
accents thinner, so that I accented only every third note of the right hand, then many of
the notes in the left hand would be struck simultaneously with weak notes in the right.
This may sound simple enough, but I found it very tricky. The difference is shown in
Figure 9-6 (which, like most of the others in this article, was created by Dori Byrd's
program).

FIGURE 9-6. The opening of Etude Op. 25, No. 2, printed in two ways. In (a), as Chopin
penned it, and as it is usually conceived: in threes. In (b), as I first heard it and first
learned to play it: in twos. [Music printed by Donald Byrd's SMUT program at Indiana
University. ]




Pattern, Poetry, and Power in the Music of Frédéric Chopin                             184

        Even after mastering the right-hand solo in triplets, I found that 1 put the parts
together, it was at first nearly impossible to keep from accenting the melodic notes
coinciding with the bass. It was a fearsome of coordination, yet I enjoyed it greatly. After
a while something "snapped into place", and I found I was doing it. It was not something
could consciously control or explain; I simply was playing it right, a sudden. Huneker, in
his commentary on this etude, quotes Theodor Kullak another Chopin specialist, about
the "algebraic character o tone-language" and then adds his own image: "At times so
delicate design that it recalls the faint fantastic tracery made by frost on gla
        Chopin's music is filled to the brim with such "algebraic" tri( cross-rhythm. He
seemed to revel in them in a way that no pre composer ever had. A famous example is his
iconoclastic waltz, Opus A-flat major, written in 1840. In this waltz, the bass line follows
the "oom-pah-pah" convention, but the melody of the first section comp counters this
three-ness; its six eighth-notes, instead of being broken into three pairs aligned with the
left hand's bounces, form two triplets, as in minor etude just discussed (see Figure 9-7).
Here, though, in contrast I nearly accentless shimmering desired in that etude, the initial
not successive triplets are to be clearly emphasized and prolonged, thus creating a higher-
level melody (shown in red) abstracted out of the quietly rip right hand. This melody is
composed of two notes per measure, be regularly against the three notes of the waltzing
bass. It is a marvelous trompe-l'oreille effect, one that Chopin exploited again in his E
major sch Opus 54, written in 1842, when he was 32.

                                       *     *     *

        In that same year, Chopin wrote what some admirers consider to b greatest work:
the fourth Ballade, in F minor. This piece is filled noteworthy passages, but one in
particular had a profound effect or One day, long after I knew the piece intimately from
recordings, a fi told me that he had been practicing it and wanted to show me "a bit of t
polyrhythm" that was particularly interesting. I was actually not interested in hearing
about polyrhythm at the moment, and so I didn’t' pay much attention when he sat down at
the keyboard. Then he started to He played just two measures, but by the time they were
over, I felt someone had reached into the very center of my skull and caused some to
explode deep down inside. This "bit of tricky polyrhythm" had un, me completely. What
in the world was going on?
        Of course, it was much more than just polyrhythm, but that is part As you can see
in our three-color plot of the two measures concerned (Figure 9-8), the left hand forms
large, rumbling waves of sound, like ocean waves on which a ship is sailing. Each wave
consists of six n forming a rising and falling arpeggio (in blue). High above these billows
of




Pattern, Poetry, and Power in the Music of Frédéric Chopin                              185

sound, a lyrical melody (in red) soars and floats, emerging out of a blur of notes swirling
around it like a halo (in black). This high melody and its halo are actually fused together
in the right hand's eighteen notes per measure. They are written as six groups of three, so
that in each half-measure, nine high notes beat against the six-note ocean wave below-
already a clear problem in three-against-two. But look: on top of those flying triplets,
there are eight-note flags placed on every fourth note! Thus there is a flag on the first
note of the first triplet, on the second note of the second triplet, on the third note of the
third triplet, on the fourth note of the fourth triplet ... Well, that cannot be. In fact, the
fourth triplet has no flag at all; the flag goes to the first note of the fifth. triplet, and the
pattern resumes. Flags waving in wind, high on the masts of a sea-borne sailing ship.
         This wonderfully subtle rhythmic construction might just might-have been
invented by anyone, say by a rhythm specialist with no feeling for melody. And yet it was
not. It was invented by a composer with a supreme gift for melody and harmony as well
as for rhythm, and this can be no coincidence. A mere "rhythms hacker" would not have
the sense to know what to do with this particular rhythm any more than with any other
rhythmic structure. There is something about this passage that shows true genius, but
words alone cannot define it. You have to hear it. It is a burning lyricism, having a power
and intensity that defy description.
         One must wonder about the soul of a man who at age 32 could write such
possessed music-a man who at the tender age of nineteen could write such perfectly
controlled and poetic outbursts as the etudes of Opus 10. Where could this rare
combination of power, poetry, and pattern, this musical self-confidence and maturity,
have come from?

                                              *    *     *

In search of an answer, one must look to Chopin's roots, both his family roots and his
roots in his native land, Poland. Chopin was born in a small and peaceful country village
30 miles west of Warsaw called Zelazowa Wola, which means Iron Will. His father,
Nicolas (Mikolaj) Chopin, was French by birth but emigrated to Poland and became an
ardent Polish patriot (so ardent, in fact, that he participated in the celebrated but ill-fated
insurrection led by the national hero Jan Kilinski in 1794 against the Russian occupation
of Warsaw). Chopin's mother, Justyna Krzyzanowska, was a distant relative of the rich
and aristocratic Skarbek family, who lived in Zelazowa Wola. She lived with them as a
family member and took care of various domestic matters. When Mikolaj Chopin came to
be the tutor of the Skarbek children, he and Justyna met and married. In addition to being
a gentle and loving mother, she was as fervent a Polish patriot as her husband, and had a
romantic and dreamy streak. They had four children, of whom Frederic, born in 1810,
was the second. The other three children were girls, one of whom died young, of
tuberculosis-a disease that in the end would




Pattern, Poetry, and Power in the Music of Frédéric Chopin                                  186

claim Frederic as well, at age 39. The four children doted on one another It was a close-
knit family, and all in all, Chopin had a very happy childhood(
    The family moved to Warsaw when Frederic was very young, and they he was
exposed to culture of all kinds, since his father was a teacher an knew university people
of all disciplines. Frederic was a fun-loving an spirited boy. The summer he was fourteen
he spent away from home in lilac-filled village called Szafarnia. He wrote home a series
of letters gleefull ,hocking the style of the Warsaw Courier, a gossipy provincial paper of
th times. One item from his "Szafarnia Courier" ran as follows (in full):

   The Esteemed Mr. Pichon [an anagram of "Chopin"] was in Golub on the 26th of
   the current month. Among other foreign wonders and oddities, he came:across a
   foreign Pig, which Pig quite specially attracted the attention of this most
   distinguished Voyageur.

   Chopin's musical talent, something he shared with his mother, emerge( very early and
was nurtured by two excellent piano teachers, first by a gentle and good-humored old
Czech named Wojciech Zywny, and later by the director of the Warsaw Conservatory,
Jozef Eisner.
   Chopin grew up in the capital city of the "Grand Duchy of Warsaw"what little
remained of Poland after it had been decimated, in three successive "partitions" in the late
eighteenth century, by its greed) neighbors: Russia, Prussia, and Austria. The turn of the
century was marked by a mounting nationalistic fervor; in Warsaw and Cracow, the two
main Polish cities, there occurred a series of rebellions against the foreign occupiers, but
to no avail. A number of ardent Polish nationalists went abroad and formed "Polish
Legions" whose purpose was to fight for the liberation of all oppressed peoples and to
eventually return to Poland and reclaim it from the occupying powers. When Napoleon
invaded Russia in 1806, a Polish state was established for a brief shining instant; then all
was lost again. The Polish nation's flame flickered and nearly went out totally, but as the
words to the Polish national anthem proclaim, "Jeszcze Polska nit zginela, poki my
zyjemy." It is a curious sentence, built out of past and present tenses, and literally
translated it runs: "Poland has not yet perished, as long as we live." The first clause
sounds so fatalistic, as if to admit that Poland surely will someday perish, but not quite
yet! Some Poles tell me that the connotations are not that despairing, that a better overall
translation would be, "Poland will not perish, as long as we live." Others, though, tell me
that the construction is subtly ambiguous, that its meaning floats somewhere between
grim fatalism and ardent determination.

                                         *    *     *

The Poles are a people who have learned to distinguish sharply between two conceptions
of Poland: Poland the abstract social entity, at whose core




Pattern, Poetry, and Power in the Music of Frédéric Chopin                              187

are the Polish language and culture, and Poland the concrete geographical entity, the land
that Poles live in. Narod polski-the "Polish nation"represents a spirit rather than a piece
of territory, although of course the nation came into existence because of the bonds
between people who lived in a certain region. It is the fragility of this flickering flame,
and the determination to keep it alive, that Chopin's music reflects so purely and
poignantly. There is a certain fusion of bitterness, anger, and sadness called zal that is
uniquely Polish. One hears it, to be sure, in the famous mazurkas and polonaises, pieces
that Chopin composed in the form of national dances. The mazurkas are mostly smaller
pieces based on folk-like tunes with a lilting 3/4 rhythm; the polonaises are grand, heroic,
and martial in spirit. But one hears this burning flame of Poland just as much in many of
Chopin's other pieces-for example, in the slow middle sections of such pieces as the
waltzes in A minor (Op. 34, No. 2) and A-flat major (Op. 64, No. 3), the pathos-filled
Prelude in F-sharp major (Op. 28, No. 13), and particularly in the middle part of the F-
sharp minor Polonaise (Opus 44), where a ray of hope bursts through dark visions like a
gleam in the gloom. One hears zal in the angry, buzzing harmonies of the etude in C-
sharp minor (Op. 10, No. 4) and in the passion of the etude in E major (Op. 10, No. 3). In
fact, Chopin is said to have cried out once, on hearing this piece played in his presence,
"O ma patrie!" ("0 my homeland!").
        But aside from the fervent patriotism of Chopin's music there is in it that different
and softer kind of Polish nostalgia: tcsknota. It is his yearning for home-for his childhood
home, for his family, for a dream-Poland that at age twenty he had left forever. In 1830,
at the height of the turmoil in Warsaw, Chopin set out for France. He had a premonition
that he would never return. Traveling by way of Vienna, he made slow progress. When
things boiled over in late 1831-when, in September 1831, the Russians finally crushed the
desperate Warsaw insurrection-Chopin was in Stuttgart. On hearing the news, he was
overwhelmed with agitation and grief, partly out of fear for the fate of his family, partly
out of love for his stricken homeland. He wavered about going back to Poland and
fighting for his nation, but the idea eventually receded from his mind.
It was at about this time that he composed the twelfth and final etude of his Opus 10. Of
this etude, Chopin's Polish biographer Maurycy Karasowski wrote:

   Grief, anxiety, and despair over the fate of his relatives and his dearly beloved
   father filled the measure of his sufferings. Under the influence of this mood he
   wrote the C minor etude, called by many the "Revolutionary Etude". Out of the
   mad and tempestuous storm of passages for the left hand the melody rises aloft,
   now passionate and anon proudly majestic, until thrills of awe stream over the
   listener, and the image is evoked of Zeus hurling thunderbolts at the world.




Pattern, Poetry, and Power in the Music of Frédéric Chopin                               188

This is pretty strong language. Huneker echoes these sentiments, as doe the French
pianist Alfred Cortot, who in his famous Student's Edition of th etudes refers to the piece
as "an exalted outcry of revolt .... wherein the emotions of a whole race of people are
alive and throbbing." I myself hay never found this etude as overwhelming as these
authors do, although it i unquestionably a powerful outburst of emotion. If someone had
told m that one of the etudes had come to be known as the "Revolutionary Etude and had
asked me to guess which one, I would certainly have picked one c the last two of Opus
25, either No. 11 in A minor, the one pictured at th beginning of this article, with its
tumultuous cascades of notes in the right hand against the surging, heroic melody in the
left hand, or else No. 12 i C minor, which sounds to me like a glowing inferno seen at
night from fa away, flaring up unpredictably and awesomely. As for the actual
"Revolutionary Etude", I have always found its ending enigmatic fluctuating as it does
between major and minor, between the keys of F an C, like an indecisive thunderclap.
        Still, this piece, like the martial A-flat major Polonaise (Opus 53), ha become a
symbol of the tragic yet heroic Polish fate. Wherever an whenever it is played, it is
special to Poles; their hearts beat faster, and their spirits cannot fail to be deeply moved. I
will never forget how I heard i nightly as the clarion call of Poland, when, from a small
town in German in 1975, I would try to tune in Radio Warsaw. Two measures of shrill
rousing chords above a roaring left hand, like a call to arms, were repeate4 over and over
again as the call signal, preceding a nightly broadcast c Chopin's music. Nor will I ever
forget how that feeble signal of Radio Warsaw faded in and out, symbolizing to me the
flickering flame of Poland's spirit.

                                             *    *     *

However one chooses to describe it-whether in terms of zal and tesknota or patriotyzm
and polyrhythm, or chromaticism and arpeggios-Chopin' music has had a deep influence
on the composers of succeeding generations. It is perhaps most visible in the piano music
of Alexander Scriabin, Sergei Rachmaninoff, Gabriel Faure, Felix Mendelssohn, Robert
and Clara Schumann, Johannes Brahms, Maurice Ravel, and Claude Debussy, but
Chopin's influence is far more pervasive than even that would suggest. It has become one
of the central pillars of Western music, and a such it has its effect on the music perceived
and created by everyone in the Western world.
In one way, Chopin's music is purely Polish, and that Polishness polsknose-extends even
to foreign-inspired pieces such as his Bolero Tarantella, Barcarolle, and so on. In another
way, though, Chopin's music




Pattern, Poetry, and Power in the Music of Frédéric Chopin                                 189

is universal, so that even his most deeply Polish pieces-the mazurkas and polonaises-
speak to a common set of emotions in everyone. But what are these emotions? How are
they so deeply evoked by mere pattern? What is the secret magic of Chopin? I know of
no more burning question.


Post Scriptum.
This column is a unique one, in that it expresses certain kinds of emotions that are not
expressed as directly in my other published writings. But the part of me represented by it
is no smaller and no less important than the part of me from which my other writings
flow. It was provoked, of course,
by the worsening crisis in Poland in late 1981, just at the time of the takeover by the
military and the tragic collapse of Solidarity. In fact, it was almost exactly 150 years after
the tragic takeover of Warsaw by the Russians that triggered the Revolutionary Etude. I
guess Poland has not yet perishedbut it is certainly going through terrible tribulations,
once again.
I received some heart-warming correspondence in response to this column. One letter,
from Andrzej Krasinski, a Pole living in West Germany, ran this way:

   I just read your nice article about Chopin's music in the April issue of Scientific
   American in which you have shown so much sympathy and understanding for a
   Polish soul, and so much care for the Polish language. I enjoyed it a lot, although I
   am no expert in music. However, by my birth, I happen to be an expert in the Polish
   language, and I wish to point out a minor error you have made. The name of the
   village where Chopin was born, Zelazowa Wola, does not mean "Iron Will",
   although you might have picked such a meaning by looking for the two words in a
   dictionary separately. The word wola, which means "will" alone, when applied as a
   part of a village's name means that the village was founded by somebody's will, and
   then the other part of the village's name usually stems from a person's name. There
   are numerous examples of such names in Poland, and normally they are attached to
   small hamlets. Consequently, Wola as a village's name has a second meaning in
   Polish, and that is simply "small village". The word Zelazowa does not seem to
   stem from a person's name (although I have no literature here to answer that
   question with certainty). It suggests that the founding of the village had something
   to do either with iron ore being found somewhere in the neighborhood or with iron
   being processed there. So the best translation of Zelazowa Wola would be "Iron
   Village" or "Iron-Ore Village". "Iron will" in Polish would be Zelazna Wola, and
   the name of Chopin's village does quite certainly not mean that.

I stand corrected!




Pattern, Poetry, and Power in the Music of Frédéric Chopin                                190

        Jakub Tatarkiewicz, a physicist writing from Warsaw, very gently pointed out that
I had somehow managed to invent a new Polish word: polsknosc. I was quite surprised to
learn that I had invented it, since I was sure I had sees it somewhere, but as it tu rns out,
what I had actually seen was polskose (will no n'). Tatarkiewicz complimented me,
however, for my talent in coming up with a good neologism, for, he said, my word has
poignant overtones o such loaded words as tfsknota and Solidarnosc. As he put it: "I can
only doubt if you really meant all those connotations-or is it just Chopin's music that
played in your soul?!" I don't know. I guess I'd chalk it up to serendipity
        Great art has a way of evoking continual commentary; it is a bottomless source of
inspiration to others. I have my blind spots in terms o understanding music, that's for
sure; but Chopin hits some kind of bull's-eyl in my soul. If I could meet any one person
from the past, it would be Chopin without any doubt. What saddens me enormously is his
relatively small output. He died at age 39, with his expressive powers clearly as strong a;
ever. What ever would he have produced, had he lived to the age of, say, 65 as Bach did?
Unbelievable firegems, I am sure. Indeed, I cannot imagine who I would be if I knew
those pieces.




Pattern, Poetry, and Power in the Music of Frédéric Chopin                               191


                      Parquet Deformations:
                    A Subtle, Intricate Art Form
                                         July, 1983


WHATS the difference between music and visual art? If I were asked this, I would
have no hesitation in replying. To me, the major difference is clearly temporality. Works
of music intrinsically involve time; works of art do not. More precisely, pieces of music
consist of sounds intended to be played and heard in a specific order and at a specific
speed. Music is thus fundamentally one-dimensional; it is tied to the rhythms of our
existence. Works of visual art, by contrast, are generally two-dimensional or three-
dimensional. Paintings and sculptures seldom have any intrinsic "scanning order" built
into them that the eye must follow. Mobiles and other pieces of kinetic art may change
over time, but often without any specific initial state or final state or intermediate stages.
You are free to come and go as you please.
        There are exceptions to this generalization, of course. European art has its grand
friezes and historic cycloramas, and Oriental art has intricate pastoral scrolls of up to
hundreds of feet in length. These types of visual art impose a temporal order and speed
on the scanning eye. There is a starting point and a final point. Usually, as in stories,
these points represent states of relative calm-especially the end. In between them, various
types of tension are built up and resolved in an idiosyncratic but pleasing visual rhythm.
The calmer end states are usually orderly and visually simple, while the tenser
intermediate states are usually more chaotic and visually confusing. If you replace
"visual" by "aural", virtually the same could be said of music.
        I have been fascinated for many years by the idea of trying to capture the essence
of the musical experience in visual form. I have my own ideas as to how this can be done;
in fact, I spent. several years working out a form of visual music. It is perhaps the most
original and creative thing I have ever done. However, by no means do I feel that there is
a unique or best way to carry out this task of "translation", and indeed I have often
wondered how




Parquet Deformations: A Subtle, Intricate Art Form                                        191

others might attempt to do it. I have seen a few such attempts, but most of them,
unfortunately, did not grab me. One striking counterexample is the set of "parquet
deformations" meta-composed by William Huff, a professor of architectural design at the
State University of New York at Buffalo.
        I say "meta-composed" for a very good reason. Huff himself has never executed a
single parquet deformation. He has elicited hundreds of them, however, from his
students, and in so doing has brought this form of art to a high degree of refinement. Huff
might be likened to the conductor of a fine orchestra, who of course makes no sound
whatsoever during a performance. And yet we tend to give the conductor most of the
credit for the quality of the sound. We can only guess how much preparation and
coaching went into this performance. And what about the selection of the pieces and
tempos and styles-not to mention the many-year process of culling the performers
themselves?
        So it is with William Huff. For 23 years, his students at Carnegie-Mellon and
SUNY at Buffalo have been prodded into flights of artistic inspiration, and it is thanks to
Huff's vision of what constitutes quality that some very beautiful results have emerged.
Not only has he elicited outstanding work from students, he has also carefully selected
what he feels to be the best pieces and these he is preserving in archives. For these
reasons, I shall at times refer to Huff's "creations", but it is always in this more indirect
sense of "meta-creations" that I shall mean it.
        Not to take credit from the students who executed the individual pieces, there is a
larger sense of the term "credit" that goes exclusively to Huff, the person who has shaped
this whole art form himself. Let me use an analogy. Gazelles are marvelous beasts, yet it
is not they themselves but the selective pressures of evolution that are responsible for
their species' unique and wondrous qualities. Huff's judgments and comments have here
played the role of those impersonal evolutionary selective pressures, and out of them has
been molded a living and dynamic tradition, a "species" of art exemplified and extended
by each new instance.

                                         *   *   *

       All that remains to be said by way of introduction is the meaning of the term
parquet deformation. It is nearly self-explanatory, actually: traditionally, a parquet is a
regular mosaic made out of inlaid wood, on the floor of an elegant room; and a
deformation-well, it's somewhere in between a distortion and a transformation. Huff's
parquets are more abstract: they are regular tessellations (or tilings) of the plane, ideally
drawn with zero-thickness line segments and curves. The deformations are not arbitrary
but must satisfy two basic requirements:

       (1) There shall be change only in one dimension, so that one can see a temporal
progression in which one tessellation gradually becomes another;




Parquet Deformations: A Subtle, Intricate Art Form                                       192

         (2) At each stage, the pattern must constitute a regular tessellation of the plane
(i.e., there must be a unit cell that could combine with itself so as to cover an infinite
plane exactly).
         (Actually, the second requirement is not usually adhered to strictly. It would be
more accurate to say that the unit cell at any stage of a parquet deformation can be easily
modified so as to allow it to tile the plane perfectly.)
         From this very simple idea emerge some stunningly beautiful creations. Huff
explains that he was originally inspired, back in 1960, by the woodcut "Day and Night"
of M. C. Escher. In that work, forms of birds tiling the plane are gradually distorted (as
the eye scans downwards) until they become diamond-shaped, looking like the
checkerboard pattern of cultivated fields seen from the air. Escher is now famous for his
tessellations, both pure and distorted, as well as for other hauntingly strange visual games
he played with art and reality.
         Whereas Escher's tessellations almost always involve animals, Huff decided to
limit his scope to purely geometric forms. In a way, this is like a decision by a composer
to use austere musical patterns and to totally eschew anything that might conjure up a
"program" (that is, some sort of image or story behind the sounds). An effect of this
decision is that the beauty and visual interest must come entirely from the complexity and
the subtlety of the interplay of abstract forms. There is nothing to "charm" the eye, as
with pictures of animals. There is only the uninterpreted, unembellished perceptual
experience.
         Because of the linearity of this form of art, Huff has likened it to visual music. He
writes:

   Though I am spectacularly ignorant of music, tone deaf, and hated those piano
   lessons (yet can be enthralled by Bach, Vivaldi, or Debussy), I have the students
   'read' their designs as I suppose a musician might scan a work: the themes, the
   events, the intervals, the number of steps from one event to another, the rhythms,
   the repetitions (which can be destructive, if not totally controlled, as well as
   reinforcing). These are principally temporal, not spatial, compositions (though all
   predominantly temporal compositions have, of necessity, an element of the spatial
   and vice versa-e.g., the single-frame picture is the basic element of the moving
   picture).

                                      *       *       *

         What are the basic elements of a parquet deformation? First of all, there is the
class of allowed parquets. On this, Huff writes the following:
         We play a different (or rather, tighter) gate than does Escher. We work with only
A tiles (i.e., congruent tiles of the same handedness). We do not use, as he does, A and A'
tiles (i.e., congruent tiles of both handednesses). Finally, we don't use A and B tiles (i.e.,
two different interlocking tiles), since two such tiles can always be seen as subdivisions
of a single larger tile.




Parquet Deformations: A Subtle, Intricate Art Form                                        193

                                               FIGURE: 10-2 Crossover by Richard
FIGURE: 10-1 Fylfor Flipflop by Fred           Long. Created in the studio of William
Watts. Created in the studio of William        Huff (1963).
Huff (1963).




Parquet Deformations: A Subtle, Intricate Art Form                                 194

The other basic element is the repertoire of standard deforming devices. Typical devices
include:

    • lengthening or shortening a line;
    • rotating a line;
    • introducing a "hinge" somewhere inside a line segment so that it can "flex";
    • introducing a "bump" or "pimple" or "tooth" (a small intrusion or extrusion
      having a simple shape) in the middle of a line or at a vertex;
    • shifting, rotating, expanding, or contracting a group of lines that form a natural
      subunit;

and variations on these themes. To understand these descriptions, you must realize that a
reference to "a line" or "a vertex" is actually a reference to a line or vertex inside a unit
cell, and therefore, when one such line or vertex is altered, all the corresponding lines or
vertices that play the same role in the copies of that cell undergo the same change. Since
some of those copies may be at 90 degrees (or other angles) with respect to the master
cell, one locally innocent-looking change may induce changes at corresponding spots,
resulting in unexpected interactions whose visual consequences may be quite exciting.

*   *   *

        Without further ado, let us proceed to examine some specific pieces. Look at the
one called "Fylfot Flipflop" (Figure 10-1). It is an early one, executed in 1963 by Fred
Watts at Carnegie-Mellon. If you simply let your eye skim across the topmost line, you
will get the distinct sensation of scanning a tiny mountain range. At either edge, you
begin with a perfectly flat plain, and then you move into gently rolling hills, which
become taller and steeper, eventually turning into jagged peaks; then past the centerpoint,
these start to soften into lower foothills, which gradually tail off into the plain again. This
much is obvious even upon a casual glance. Subtler to see is the line just below, whose
zigging and zagging is 180 degrees out of phase with the top line. Thus notice that in the
very center, that line is completely at rest: a perfectly horizontal stretch flanked on either
side by increasingly toothy regions. Below it there are seven more horizontal lines. Thus
if one completely filtered out the vertical lines, one would see nine horizontal lines
stacked above one another, the odd-numbered ones jagged in the center, the even-
numbered ones smooth in the center.
        Now what about the vertical lines? Both the lefthand and righthand borderlines
are perfectly straight vertical lines. However, their immediate neighbors are as jagged as
possible, consisting of repeated 90-degree bends, back and forth. Then the next vertical
line nearer the center is practically straight up and down again. Then there is a wavy one
again, and so on. As




Parquet Deformations: A Subtle, Intricate Art Form                                         192

you move across the picture, you see that the jagged ones gradually get less jagged and
the straight ones get increasingly jagged, so that in the middle the roles are completely
reversed. Then the process continues, so that by the time you've reached the other side,
the lines are back to normal again. If you could filter out the horizontal lines, you would
see a simple pattern of quite jaggy lines alternating with less jaggy lines.
         When these two extremely simple independent patterns-the horizontal and the
vertical-are superimposed, what emerges is an unexpectedly rich perceptual feast. At the
far left and right, the eye picks out fylfots-that is, swastikas-of either handedness
contained inside perfect squares. In the center, the eye immediately sees that the central
fylfots are all gone, replaced by perfect crosses inside pinwheels.
         And then a queer perceptual reversal takes place. If you just shift your focus of
attention diagonally by half a pinwheel, you will notice that there is a fylfot right there
before your eyes! In fact, suddenly they appear all over the central section where before
you'd been seeing only crosses inside pinwheels! And conversely, of course, now when
you look at either end, you'll see pinwheels everywhere with crosses inside them. No
fylfots! It is an astonishingly simple design, yet this effect catches nearly everyone really
off guard.
         This is a simple example of the ubiquitous visual phenomenon called regrouping,
in which the boundary line of the unit cell shifts so that structures jump out at the eye that
before were completely submerged and invisible -while conversely, of course, structures
that a moment ago were totally obvious have now become invisible, having been split
into separate conceptual pieces by the act of regrouping, or shift of perceptual
boundaries. It is both a perceptual and conceptual phenomenon, a delight to that subtle
mixture of eye and mind that is most sensitive to pattern.
         For another example of regrouping, take a look at "Crossover" (Figure 10-2), also
executed at Carnegie-Mellon in 1963 by Richard Lane. Something really amazing
happens in the middle, but I won't tell you what. Just find it yourself by careful looking.
         By the way, there are still features left to be explained in "Fylfot Flipflop". At
first it appears to be mirror-symmetric. For instance, all the fvlfots at the left end are
spinning counterclockwise, while all the ones at the right end are spinning clockwise. So
far, so symmetric. But in the middle, all the fylfots go counterclockwise. This surely
violates the symmetry. Furthermore, the one-quarter-way and three-quarter-way stages of
this deformation, which ought to be mirror images of each other, bear no resemblance at
all to each other. Can you figure out the logic behind this subtle asymmetry between the
left and right sides?
         This piece also illustrates one more way in which parquet deformations resemble
music. A unit cell-or rather, a vertical cross-section consisting of a stack of unit cells-is
analogous to a measure •in music. The regular pulse of a piece of music is given by the
repetition of unit cells across the page.




Parquet Deformations: A Subtle, Intricate Art Form                                        193

And the flow of a melodic line across measure boundaries is modeled by the flow of a
visual line-such as the mountain range lines-across many unit cells.

                                         *   *   *

        Bach's music is always called up in discussions of the relationship of
mathematical patterns to music, and this occasion is no exception. I am reminded
especially of some of his texturally more uniform pieces, such as certain preludes from
the Well-Tempered Clavier, in which in each measure there is a certain pattern executed
once or twice, possibly more times. From measure to measure this pattern undergoes a
slow metamorphosis, meandering over the course of many measures from one region of
harmonic space to far distant regions and then slowly returning via some circuitous route.
For specific examples, you might listen to (or look at the scores of): Book I, numbers 1,
2; Book II, numbers 3, 15. Many of the other preludes have this feature in places, though
not for their entirety.
        Bach seldom deliberately set out to play with the perceptual systems of his
listeners. Artists of his century, although they occasionally played perceptual games,
were considerably less sophisticated about, and less fascinated with, issues that we now
deem part of perceptual psychology. Such phenomena as regrouping would undoubtedly
have intrigued Bach, and I for one sometimes wish that he had known of and been able to
try out certain effects-but then I remind myself that whatever time Bach might have spent
playing with new-fangled ideas would have had to be subtracted from his time to produce
the masterpieces that we know and love, so why tamper with something that precious?
        On the other hand, I don't find that argument 100 percent compelling. Who says
that if you're going to imagine playing with the past, you have to hold the lifetimes of
famous people constant in length? If we can imagine telling Bach about perceptual
psychology, why can't we also imagine adding a few extra years to his lifetime to let him
explore it? After all, the only divinely imposed (that is, absolutely unslippable) constraint
on Bach's years is that they and Mozart's years add up to 100, no? So if we award Bach
five extra ones, then we merely take five years away from Mozart. It's painful, to be sure,
but not all that bad. We could even let Bach live to 100 that way! (Mozart would never
have existed.) It starts to get a little questionable if we go much beyond that point,
however, since it is not altogether clear what it means to live a negative number of years.
        Although it is difficult to imagine and impossible to know what Bach's music
would have been like had he lived in the twentieth century, it is certainly not impossible
to know what Steve Reich's music would have been like, had he lived in this century. In
fact, I'm listening to a record of it right now (or at least I would have been if I hadn't
gotten distracted by this radio program). Now Reich's is music that really is conscious of
perceptual psychology. All the way through, he plays with perceptual shifts and




Parquet Deformations: A Subtle, Intricate Art Form                                       194

ambiguities, pivoting from one rhythm to another, from one harmonic origin to another,
constantly keeping the listener on edge and tingling with nervous energy. Imagine a piece
like Ravel's "Bolero", only with a much finer grain size, so that instead of roughly a one-
minute unit cell, it has a three-second unit cell. Its changes are tiny enough that
sometimes you barely can tell it is changing at all, while other times the changes jump
out at you. What Reich piece am I listening to (or rather, would I be listening to if I
weren't still listening to this radio program)? Well, it hardly matters, since most of them
satisfy this characterization, but for the sake of specificity you might try "Music for a
Large Ensemble", "Octet", "Violin Phase", "Vermont Counterpoint", or his recent choral
work "Tehillim".

                                         *   *   *

        Let us now return to parquet deformations. "Dizzy Bee" (Figure 10-3), executed
by Richard Mesnik at Carnegie-Mellon in 1964, involves --perceptual tricks of another
sort. The left side looks like a perfect honeycomb or-somewhat less poetically-a perfect
bathroom floor. However, as we move rightward, its perfection seems cast in doubt as the
rigidity of the lattice gives way to rounder-seeming shapes. Then we notice that three of
them have combined to form one larger shape: a super hexagon made up of three rather
squashed pentagons. The curious thing is that if we now sweep our eyes right to left, back
to the beginning, we can no longer

FIGURE 10-3. Dizzy Bee, by Richard Mesnik. Created in the studio of William Huff
(1964).




Parquet Deformations: A Subtle, Intricate Art Form                                     195

FIGURE 10-4. Consternation, by Scott Grady. Created in the studio of William Huff
(1977).

see the left side in quite the way we saw it before. The small hexagons now are constantly
grouping themselves into threes, although the grouping changes quickly. We experience
"flickering clusters" in our minds, in which groups form for an instant and then disband,
their components immediately regrouping in new combinations, and so on. The poetic
term "flickering clusters" comes from a famous theory of how water molecules behave,
the bonding in that case coming from hydrogen bonds rather than mental ones. (See the
P.S. to Chapter 26.)
        Even more dizzying, perhaps, than "Dizzy Bee" is "Consternation" (Figure 10-4),
executed by Scott Grady of SUNY at Buffalo in 1977. This is another parquet
deformation in which hexagons. ,and cubes vie for perceptual supremacy. This one is so
complex and agitated in appearance that I scarcely dare to attempt an analysis. In its
intermediate regions, I find the same extremely exciting kind of visual pseudo-chaos as in
Escher's best deformations.
        Perhaps irrelevantly, but I suspect not, the names of many of these studies remind
me of pieces by Zez Confrey, a composer most famous during the twenties for his
novelty piano solos such as "Dizzy Fingers", "Kitten on the Keys", and-my favorite-
"Flutter by, Butterfly". Confrey specialized in pushing rag music to its limits without
losing musical charm, and some of the results seem to me to have a saucy, dazzling
appeal not unlike the jazzy appearance of this parquet deformation, and others.
        The next parquet deformation, "Oddity out of Old Oriental Ornament"




Parquet Deformations: A Subtle, Intricate Art Form                                    196

FIGURE 10-5. Oddity out of Old Oriental Ornament, by Francis O'Donnell. Created in
the studio of William Huff (1966).

(Figure 10-5), executed by Francis O'Donnell at Carnegie-Mellon in 1966, is based on an
extremely simple principle: the insertion of a "hinge" in one single line segment, and
subsequent flexing of the segment at that hinge! The reason for the stunningly rich results
is that the unit cell that creates the tessellation occurs both vertically and horizontally, so
that flexing it one way induces a crosswise flexing as well, and the two flexings combine
to yield this curious and unexpected pattern.
         Another one that shows the amazing results of an extremely simple but carefully
chosen tranformation principle is "Y Knot" (Figure 10-6),




          FIGURE 10-6. Y Knot, hi Leland Chen. Created in the studio of William Huff
(1977).

executed by Leland Chen at SUNY at Buffalo in 1977. If you look at it with full
attention, you will see that its unit cell is in the shape of a three-bladed propeller, and that
unit cell never changes whatsoever in shape. All that does change is the 'Y' lodged tightly
inside that unit cell. And the only way that 'Y' changes is by rotating clockwise very
slowly! Admittedly, in the final stages of rotation, this forces some previously constant
line segments to extend themselves a little bit, but this does not change the outline of the
unit cell whatsoever. What well-chosen simplicity can do!




Parquet Deformations: A Subtle, Intricate Art Form                                         197

          FIGURE 10-7. Crazy Cogs, by Arne Larson. Created in the studio of William Huff
(1963).


Three of my favorites are "Crazy Cogs" (Figure 10-7, done by Arne Larson, Carnegie-
Mellon, 1963), "Trifoliolate" (Figure 10-8, done by Glen Paris, Carnegie-Mellon, 1966),
and "Arabesque" (Figure 10-9, done by Joel Napach, SUNY at Buffalo, 1979). They all
share the feature of getting more and more intricate as you move rightward. Most of the
earlier ones we've seen don't have this extreme quality of irreversibility-that is, the
ratcheted quality that signals that an evolutionary process is taking place. I can't help
wondering if the designers didn't feel that they'd painted themselves into a corner,
especially in the case of "Arabesque". Is there any way you can back out of that super-
tangle except by retrograde motion-that is, retracing your steps? I suspect there is, but I
wouldn't care to try to discover it.
        To contrast with this, consider "Razor Blades", an extended study in relative
calmness (Figure 10-10). It was done at Carnegie-Mellon in 1966, but unfortunately it is
unsigned. Like the first one we discussed, this one can be broken up into very long
waving horizontal lines and vertical structures crossing them. It's a little easier to see
them if you start at the right side. For instance, you can see that just below the top, there
is a long snaky line

FIGURE 10-8. Trifoliolate, by Glen Paris. Created n( //), hub,-I\_ William Huff (1966).




Parquet Deformations: A Subtle, Intricate Art Form                                       198

FIGURE 10-9: Arabesque by Joel Napach Created in the studio of William Huff (1979).




Parquet Deformations: A Subtle, Intricate Art Form                               199

FIGURE 10-10: Razaor Blades (unsigned) Created in the studio of William Huff (1966).




Parquet Deformations: A Subtle, Intricate Art Form                               192

with numerous little "nicks" in it, undulating its way leftwards and in so doing shedding
some of those nicks, so that at the very left edge it has degenerated into a perfect "square
wave", as such a periodic wave form is called in Fourier analysis. Complementing this
horizontal structure is a similar vertical structure that is harder to describe. The thought
that comes to my mind is that of two very ornate, rather rectangular hourglasses with
ringed necks, one on top of the other. But you can see for yourself.
        As with "Fylfot Flipflop" (Figure 10-1), each of these patterns by itself is
intriguing, but of course the real excitement comes from the daring act of superimposing
them. Incidentally, I know of no piece of visual art that better captures the feeling of
beauty and intricacy of a Steve Reich piece, created by slow "adiabatic" changes floating
on top of the chaos and dynamism of the lower-level frenzy. Looking back, I see I began
by describing this parquet deformation as "calm". Well, what do you know? Maybe I
would be a good candidate for inclusion in The New Yorker's occasional notes titled
"Our Forgetful Authors".
        More seriously, there is a reason for this inconsistency. One's emotional response
to a given work of art, whether visual or musical, is not static and unchanging. There is
no way to know how you will respond, the next time you hear or see one of your favorite
pieces. It may leave you unmoved, or it may thrill you to the bones. It depends on your
mood, what has recently happened, what chances to strike you, and many other subtle
intangibles. One's reaction can even change in the course of a few minutes. So I won't
apologize for this seeming lapse.
        Let us now look at "Cucaracha" (Figure 10-11), executed in 1977 by Jorge
Gutierrez at SUNY at Buffalo. It moves from the utmost geometricity-a lattice of perfect
diamonds-through a sequence of gradually more arbitrary modifications until it reaches
some kind of near-freedom, a dance of strange, angular, quasi-organic forms. This
fascinates me. Is entropy increasing or decreasing in this rightward flow toward freedom?
        A gracefully spiky deformation is the one wittily titled "Beecombing Blossoms"
(Figure 10-12), executed this year by Laird Pylkas at SUNY at Buffalo. Huff told me that
Pylkas struggled for weeks with this one, and at the end, when she had satisfactorily
resolved her difficulties, she mused, "Why is it that the obvious ideas always take so long
to discover?"

                                         *   *   *

As our last study, let us take "Clearing the Thicket" (Figure 10-13), executed in 1979 by
Vincent Marlowe at SUNY at Buffalo, which involves a mixture of straight lines and
curves, right angles and cusps, explicit squarish swastikoids and implicit circular holes.
Rather than demonstrate my inability to analyze the ferocious complexity of this design, I
would like to use it as the jumping-off point for a discussion of computers and creativity -
one of my favorite hobbyhorses.




Parquet Deformations: A Subtle, Intricate Art Form                                      193

                                               FIGURE 10-12: Becoming Blossoms, by
                                               Laird Pylhas. Created in the studio of
                                               Wiolliam Huff (1983).




FIGURE 10-11: cucuracha by Jorge
Gutierrez Created in the studio of
William Huff (1977).




Parquet Deformations: A Subtle, Intricate Art Form                               194

FIGURE 10-13. Clearing the Thicket, by Vincent Marlowe. Created in the studio of
William Huff (1979).

         Some totally new things are going on in this parquet deformation-things that have
not appeared in any previous one. Notice the hollow circles on the left side that shrink as
you move rightward; notice also that on the right side there are hollow "anticircles"
(concave shapes made from four circular arcs turned inside out) that shrink as you move
leftward. Now, according to Huff, such an idea had never appeared in any previously
created deformations. This means that something unusual happened here-something
genuinely creative, something unexpected, unpredictable, surprising, intriguing-and not
least, inspiring to future creators.
         So the question naturally arises: Would a computer have been able to invent this
parquet deformation? Well, put this way it is a naive and ill-posed question, but we can
try to make some sense of it. The first thing to point out is that, of course, the phrase "a
computer" refers to nothing more than an inert hunk of metal and semiconductors. To go
along with this bare computer, this hardware, we need some software and some energy.
The former is a specific pattern inserted into the matter binding it with constraints yet
imbuing it with goals; the latter is what breathes "life" into it, making it act according to
those goals and constraints.
         The next point is that the software is what really controls what the machine does;
the hardware simply obeys the software's dictates, step by step. And yet, the software
could exist in a number of different "instantiations"-that is, realizations in different
computer languages. What really counts about the software is not its literal aspect, but a
more abstract, general, overall "architecture", which is best described in a nonformal
language, such as English. We might say that the plan, the sketch, the central idea of a
program is what we are talking about here-not its final realization in some specific formal
language or dialect. That is something we can leave to apprentices to carry out, after we
have presented them with our informal sketch.
         So the question actually becomes less mundane-sounding, more theoretical and
philosophical: Is there an architecture to creativity? Is there a




Parquet Deformations: A Subtle, Intricate Art Form                                       192

plan, a scheme, a set of principles that, if elucidated clearly, could account for all the
creativity embodied in the collection of all parquet deformations, past, present, and
future?

                                         *   *   *

        Note that we are asking about the collection of parquet deformations, not about
some specific work. It is a truism that any specific work of art can be recreated, even
recreated in various slightly novel ways, by a programmed computer.
        For example, the Dutch artist Piet Mondrian evolved a highly idiosyncratic,
somewhat cryptic style of painting over a period of many years.. You can see, if you trace
his development over the course of time, exactly where he came from and where he was
headed. But if you focus in on just a single Mondrian work, you cannot sense this stylistic
momentum -this quality of dynamic, evolving style that any great artist has. Looking at
just one work in isolation is like taking a snapshot of something in motion: you capture
its instantaneous position but not its momentum. Of course, the snapshot might be
blurred, in which case you get a sense of the momentum but lose information about the
position. But when you are looking at just a single work of art, there is no mental blurring
of its style with that of recent works or soon-to-come works; you have exact position
information ("What is the style now ?"), but no momentum information ("Where was it
and where is it going?").
        Some years ago, the mathematician and computer artist A. Michael Noll took a
single Mondrian painting-an abstract, geometric study with seemingly random elements-
and from it extracted some statistics concerning the patterns. Given these statistics, he
then programmed a computer to generate numerous "pseudo-Mondrian paintings" having
the same or different values of these randomness-governing parameters. (See Figure 10-
14.) Then he showed the results to naive viewers. The reactions were interesting, in that
more people preferred one of the pseudo-Mondrians to the genuine Mondrian!
        This is quite amusing, even provocative, but it also is a warning. It proves that a
computer can certainly be programmed, after the fact, to imitate-and well-mathematically
capturable stylistic aspects of a given work. But it also warns us: Beware of cheap
imitations!
        Consider the case of parquet deformations. There is no doubt that a computer
could be programmed to do any specific parquet deformation-or minor variations on it-
without too much trouble. There just aren't that many parameters to any given one. But
the essence of any artistic act resides not in selecting particular values for certain
parameters, but far deeper: it's in the balancing of a myriad intangible and mostly
unconscious mental forces, a judgmental act that results in many conceptual choices that
eventually add up to a tangible, perceptible, measurable work of art.




Parquet Deformations: A Subtle, Intricate Art Form                                      193

FIGURE 10-14. One genuine Mondrian plus three computer imitations. Can you spot the
Mondrian? If you rotate the figure so that east becomes south, it will be the one in the
northwest corner. The Mondrian, done in 1917, is titled Composition with Lines; the
three others, done in 1965, comprise a work called Computer Composition with Lines,
and were created by a computer at Bell Telephone Laboratories at the behest of
computer tamer A. Michael Noll. The subjectively "best "picture was found through
surveys; it is the one diagonally opposite thegenuine Mondrian!

        Once the finished work exists, scholars looking at it may seize upon certain
qualities of it that lend themselves easily to being parametrized. Anyone can do statistics
on a work of art once it is there for the scrutiny, but the ease of doing so can obscure the
fact that no one could have said, a priori, what kinds of mathematical observables would
turn out to be relevant to the capturing of stylistic aspects of the as-yet-unseen work of
art.
        Huff's own view on this question of mechanizing the art of parquet deformations
closely parallels mine. He-believes that some basic principles could be formulated at the
present time enabling a computer to come up




Parquet Deformations: A Subtle, Intricate Art Form                                      194

with relatively stereotyped yet novel creations of its own. But, he stresses, his students
occasionally come up with rule-breaking ideas that noneth°less enchant the eye for
deeper reasons than he has so far been able to verbalize. And so, this way, the set of
explicit rules gets gradually increased.
        Comparing the creativity that goes into parquet deformations with the creativity
of a great musician, Huff has written:

           I don't know about the consistency of the genius of Bach, but I did work
   with the great American architect Louis Kahn (1901- 1974) and suppose that it
   must have been somewhat the same with Bach. That is, Kahn, out of moral,
   spiritual, and philosophical considerations, formulated ways he would and ways he
   would not do a thing in architecture. Students came to know many of his ways, and
   some of the best could imitate him rather well (though not perfectly). But as Kahn
   himself developed, he constantly brought in new principles that brought new
   transformations to his work; and he even occasionally discarded an old rule.
   Consequently, he was always several steps ahead of his imitators who knew what
   was but couldn't imagine what will be. So it is that computer-generated `original'
   Bach is an interesting exercise. But it isn't Bach -that unwritten work that Bach
   never got to, the day after he died.

        The real question is: What kind of architecture is responsible for all of these
ideas? Or is there any one architecture that could come up with them all? I would say that
the ability to design good parquet deformations is probably deceptive, in the same way as
the ability to play good chess is: it looks more mathematical than it really is.
        A brilliant chess move, once the game is' over and can be viewed in retrospect,
can be seen as logical-as "the correct thing to do in that situation". But brilliant moves do
not originate from the kind of logical analysis that occurs after the game; there is no time
during the game to check out all the logical consequences of a move. Good chess moves
spring from the organization of a good chess mind: a set of perceptions arranged in such a
way that certain kinds of ideas leap to mind when certain subtle patterns or cues are
present. This way that perceptions have of triggering old and buried memories underlies
skill in any type of human activity, not only chess. It's just that in chess the skill is
particularly deceptive, because after the fact, it can all be justified by a logical analysis, a
fact that seems to hint that the original idea came from logic.
        Writing lovely melodies is another one of those deceptive arts. To the
mathematically inclined, notes seem like numbers and melodies like number patterns.
Therefore all the beauty of a melody seems as if it ought to be describable in some simple
mathematical way. But so far, no formula has produced even a single good melody. Of
course, you can look back at any melody and write a formula that will produce it and
variations on it. But that is retrospective, not prospective. Lovely chess moves and lovely
melodies (and lovely theorems in mathematics, etc.) have this in common: every one has
idiosyncratic nuances that seem logical a posteriori but that are not easy to




Parquet Deformations: A Subtle, Intricate Art Form                                         195

anticipate a priori. To the mathematical mind, chess-playing skill and melody-writing
skill and theorem-discovering skill seem obviously formalizable, but the truth turns out to
be more tantalizingly complex than that. Too many subtle balances are involved.
        So it is with parquet deformations, I reckon. Each one taken alone is in some
sense mathematical. However, taken as a class, they are not mathematical. This is what's
tricky about them. Don't let the apparently mathematical nature of an individual one fool
you, for the architecture of a program that could create all these parquet deformations and
more good ones would have to incorporate computerized versions of concepts and
judgments-and those are much more elusive and complex things than are numbers. In a
way, parquet deformations are an ideal case with which to make this point about the
subtlety of art, for the very reason that each one on its own appears so simple and rule-
bound.
        At this point, many critics of computers and artificial intelligence, eager to find
something that "computers can't do" (and never will be able to do) often jump too far:
they jump to the conclusion that art and, more generally, creativity, are fundamentally
uncomputerizable. This is hardly the implied conclusion! The implied conclusion is just
this: that for computers to act human, we will have to wait until we have good computer
models of such human things as perception, memory, mental categories, learning, and so
on. We are a long way from that. But there is no reason to assume that those goals are in
principle unattainable, even if they remain far off for a long time.

                                         *   *   *

         I have been playing with the double meaning, in this column, of the term
"architecture": it means both the design of a habitat and the abstract essence of a grand
structure of any sort. The former has to do with hardware and the latter with software. In
a certain sense, William Huff is a professor of both brands of architecture. Obviously his
professional training is in the design of "hardware": genuine habitats for humans, and he
is in a school where that is what they do. But he is also in the business of forming, in the
minds of his students, a softer type of architecture: the mental architecture that underlies
the skill to create beauty. Fortunately for him, he can take for granted the whole
complexity of a human brain as his starting point upon which to build this architecture.
But even so, there is a great art to instilling a sensitivity for beauty and novelty.
         When I first met William Huff and saw how abstract and seemingly impractical
were the marvelous works produced in his design studio ranging from parquet
deformations to strange ways of slicing a cube to gestalt studies using thousands of dots
to eye-boggling color patterns-I at first wondered why this man was a professor of
architecture. But after conversing with him and his colleagues, my horizons were
extended about the nature of their discipline.




Parquet Deformations: A Subtle, Intricate Art Form                                      196

       The architect Louis Kahn had great respect for the work of William Huff, and it is
with his words that I would like to conclude:

           What Huff teaches is not merely what he has learned from someone else,
   but what.is drawn from his natural gifts and belief in their truth and value. In my
   belief what he teaches is the introduction to discipline underlying shapes and
   rhythms, which touches the arts of sight, the arts of sound, and the arts of structure.
   It teaches students of drawing to search for the abstract and not the representational.
   This is so good as a reminder of order for the instructors/architectural sketchers
   (like me), and so good especially for the student sketchers without background. It is
   the introduction to exactitudes of the kind that instill the religion of the ordered
   path.


Post Scriptum.
        "The religion of the ordered path"-a lovely phrase. I did not know at the time this
column was written that it would be my last full column (the one reporting on the results
of the Luring Lottery, here Chapter 31, was only a half-column). Both William Huff and I
were pleased with my bowing out this way, and I was especially pleased with the phrase
with which I bowed out. Though ambiguous, it captures much of the spirit that I
attempted to get across in all my columns: dedicated questing after patterned beauty, and
particularly after the reasons that certain particular patterns are beautiful.
        In this column, I repeatedly claimed that it is relatively easy to make a computer
program that creates attractive art within a formula, but not at all easy to make a
computer program that constantly comes up with novelty. Some people familiar with the
computer art produced in the last couple of decades might pick a fight with me over this.
They might point to complex patterns produced by simple algorithms, and then add that
there are certain simple algorithms which, when you change merely a few parameters,
come up with astonishingly different patterns that no human would be likely to recognize
as being each other's near kin. An example is a very simple program I know, which fills a
screen with rapidly changing sixfoldsymmetric dot-patterns that look like magnified
snowflakes; in just a few seconds, any given pattern will dissolve and be replaced by an
unbelievably different sixfold-symmetric pattern. I have stood transfixed at a screen
watching these patterns unfold one after another, unable to anticipate in the slightest what
will happen next-and yet knowing that the program itself is only a few lines long! I have
seen small changes in mathematical formulas produce enormous visual changes in what
those formulas represent, graphically.
        The trouble is, these parameter-based changes-knob-twiddlings, as they are called
in Chapters 12 and 13-are of a different nature than the kinds




Parquet Deformations: A Subtle, Intricate Art Form                                       197

of novel ideas people' come up with when they vary a given idea. For a machine to make
simple variants of a given design, it must possess an algorithm for making that design
which has explicit parameters; those parameters are then modifiable, as with the pseudo-
Mondrian paintings. But the way people make variations is quite different. They look at
some creation by an artist (or computer), and then they abstract from it some quality that
they observe in the creation itself (not in some algorithm behind it). This newly
abstracted quality may never have been thought of explicitly by the artist (or programmer
or computer), yet it is there for the seeing by an acute observer. This perceptual act gets
you more than half the way to genuine creativity; the remainder involves treating this
new quality as if it

FIGURE 10-15. 1 at the Center, by David Oleson. Created in the studio of William Huff




Parquet Deformations: A Subtle, Intricate Art Form                                     198

were an explicit knob: "twiddling" it as if it were a parameter that had all along been in
the program that made the creation.
        That way, the perceptual process is intimately linked up with the generative
process: a loop is closed in which perceptions spark new potentials and experimentation
with new potentials opens up the way for new perceptions. The element lacking in
current computer art is the interaction of perception with generation. Computers do not
watch what they do; they simply do it. (See Chapter 23 for more on the idea of self-
watching computers.) When programs are able to look at what they've done and perceive
it in ways that they never anticipated, then you'll start to get close to the kinds of insight-
giving disciplined exercises that Louis Kahn was speaking of when he wrote of the
"religion of the ordered path".

                                          *   *   *

        One of my favorite parquet deformations is called "I at the Center" (Figure 10-
15), and was done by David Oleson at Carnegie-Mellon in 1964. This one violates the
premise with which I began my article: one-dimensionality. It develops its central theme-
the uppercase letter `I' -along two perpendicular dimensions at once. The result is one of
the most lyrical and graceful compositions that I have seen in this form.
        I am also pleased by the metaphorical quality it has. At the very center of a mesh
is an I-an ego; touching it are other things-other I's-very much like the central I, but not
quite the same and not quite as simple; then as one goes further and further out, the
variety of I's multiplies. To me this symbolizes a web of human interconnections. Each of
us is at the very center of our own personal web, and each one of us thinks, "I am the
most normal, sensible, comprehensible individual." And our identity-our "shape" in
personality space-springs largely from the way we are embedded in that network-which is
to say, from the identities (shapes) of the people we are closest to. This means that we
help to define others' identities even as they help to define our own. And very simply but
effectively, this parquet deformation conveys all that, and more, to me.




Parquet Deformations: A Subtle, Intricate Art Form                                         199


                             Stuff and Nonsense
                                    December, 1982

                                                    Buz, quoth the blue fly,
                                                    Hum, quoth the bee,
                                                     Buz and hum they cry,
                                                    And so do we:
                                                    In his ear, in his nose, thus, do you
                                                    see?
                                                    He ate the dormouse, else it was he.
                                                                              -Ben Jonson


EH? What does this mean? What is its point? This little nonsense poem, written
around 1600, begins with an image of insects, slides into an image of someone's face, and
concludes with an uncertain reference to the devourer of a certain rodent. Although it
makes little sense, it is still somehow enjoyable. It reminds us of a nursery rhyme. It is
comfortable, cute, droll.
       Nonsense has been around for a long time. Its style and tone have changed over
the centuries, however. The path of development of nonsense is interesting to trace. What
marks something off as being nonsense? When does nonsense spill over into sense, or
vice versa? Where are the borderlines between nonsense and poetry? These are issues to
be explored in this column.
       A century and a half after Jonson wrote his poem, an English actor named Charles
Macklin became notorious for boasting that he could memorize any passage on one
hearing. To challenge Macklin, his friend the dramatist Samuel Foote wrote the following
odd passage:

  So she went into the garden to cut a cabbage-leaf to make an apple-pie; and, at the
  same time, a great she-bear coming up the street pops its head into the shop-What!
  no soap? So he died; and, she very imprudently married the barber: and there were
  present the Picninnies, and the Joblilies, and the Garyulies, and the great
  Panjandrum himself, with the little round button at





   top. And they all fell to playing the game of `catch as catch can', till the gunpowder
   ran out at the heels of their boots.

Full of non sequiturs and awkward, choppy sentences, this must have been an excellent
challenge for Macklin. Unfortunately, we have no record as to how he fared on first
hearing it, but we do know that he enjoyed the passage immensely, and went around
reciting it with great gusto for years thereafter.
        In the nineteenth century, the reigning monarchs of nonsense were Lewis Carroll
and Edward Lear. Everyone knows Carroll's "Jabberwocky", "Tweedledum and
Tweedledee", and "The Walrus and the Carpenter"; most people have heard of Lear's
"The Owl and the Pussycat". Fewer have heard of Lear's "The Pobble Who Has No Toes"
or "The Dong with the Luminous Nose". Carroll and Lear both enjoyed inventing strange
words and using them innocently, as if they were commonplace. Their nonsense was
expressed largely in poems, where they indulged in much alliteration, many internal
rhymes, catchy rhythms, and off beat imagery. Rather than exhibit works of those two
authors, I have instead chosen, to represent their era, an anonymous poem with some of
the same charming qualities:

                                     INDIFFERENCE

                               In loopy links the canker crawls,
                               Tads twiddle in their 'polian glee,
                               Yet sinks my heart as water falls.
                               The loon that laughs, the babe that bawls,
                               The wedding wear, the funeral palls,
                               Are neither here nor there to me.
                               Of life the mingled wine and brine
                               I sit and sip pipslipsily.

       Many of Carroll's nonsense poems were parodies of popular songs or ditties of his
day. Ironically, the parodies are remembered, and the things that triggered them are
mostly completely forgotten. Carroll loved to poke fun, in his gentle manner, at the stuffy
mores and hypocritical mannerisms of society. One of the characteristics of "genteel"
poetry of the nineteenth century was its precious use of classical literary allusions. Carroll
seldom parodies this quality, but Charles Battell Loomis, a little-known writer, admirably
caught the style in this poem.

                                     A CLASSIC ODE

                       Oh, limpid stream of Tyrus, now I hear
                       The pulsing wings of Armageddon's host,
                       Clear as a colcothar and yet more clear
                       (Twin orbs, like those of which the Parsees boast);





                       Down In thy pebbled deeps in early spring
                       The dimpled naiads sport, as in the time
                       When Ocidelus with untiring wing
                       Drave teams of prancing tigers, 'mid the chime

                       Of all the bells of Phicol. Scarcely one
                       Peristome veils its beauties now, but then
                       Like nascent diamonds, sparkling in the sun,
                       Or sainfoin, circinate, or moss in marshy fen.

                       Loud as the blasts of Tubal, loud and strong,
                       Sweet as the songs of Sappho, aye more sweet;
                       Long as the spear of Arnon, twice as long,
                       What time he hurled it at King Pharaoh's feet.

This poem has the curious quality that when you read it, you feel that surely it makes
sense-perhaps another reading will reveal it to you. And then you read it again and find
that same head-scratching feeling comes back to you. This is a problem with much
modern poetry: It is very hard to be certain that you're not simply being taken for a ride
by the poet, sucked in by some practical joker who has actually nothing in mind except
tricking readers into thinking there is profound meaning where there is none.
        The limerick is a form of poetry often featured in nonsense anthologies, probably
because it is a playful form. However, very few limericks make no sense. They may
involve mild impossibilities, such as a young woman who travels faster than the speed of
light, or other more off-color feats, but in actuality, limericks are seldom nonsensical.
One limerick that, in its own way, is pure nonsense is the following gem, by W. S.
Gilbert (of Gilbert and Sullivan):

                       There was an old man of St. Bees,
                       Who was stung in the arm by a wasp.
                       When asked, "Does it hurt?"
                       He replied, "No, it doesn't
                       I'm so glad it wasn't a hornet."

Why do I call this "nonsense"? Well, if it were a prose sentence, nothing about it would
attract much attention, except perhaps the name of the town. The nonsense is certainly
not in the content, but in the way it utterly violates every standard set up for the limerick
form. It doesn't rhyme, its meter is a little bumpy, and it has absolutely nothing funny in
it-which is what makes it funny. And that makes it qualify as nonsense.





        Is nonsense always funny? Up until the twentieth century, it certainly seemed that
way. In fact, nonsense and humor have traditionally been so closely allied that
anthologies of nonsense seem to be composed largely of
humorous passages of any sort whatsoever, irrespective of how sensible they are. But
nonsense and humor took widely divergent paths in the early
twentieth century. Perhaps the greatest nonsense writer who ever lived was Gertrude
Stein, although she is seldom mentioned in this connection. Entire
collections of nonsense have been published without featuring a single piece of her work.
Her most audacious piece in this genre is a volume of nearly
400 pages, modestly titled How to Write. Here is a sample taken from the Chapter called
"Arthur a Grammar".

Arthur a grammar.
Questionnaire in question.
What is a question.
Twenty questions.
A grammar is an astrakhan coat in black and other colors it is an obliging management of
their requesting in indulgence made mainly as if in predicament as in occasion made
plainly as if in serviceable does it shine.
A question and answer.
How do you like it.
Grammar can be contained on account of their providing medaling in a ground of
allowing with or without meant because which made coupled become blanketed with a
candidly increased just as if in predicting example of which without meant and coupled
inclined as much without meant to be thought as if it were as ably rested too.
Considerable as it counted heavily in part.
What is grammar when they make it round and round. As round as they are called.
Did they guess whether they wished. A politely definitely detailed blame' of when they
go.
What is a grammar ordinarily. A grammar is question and answer answer undoubted
however how and about.
What is Arthur a grammar.
Arthur is a grammar.
Arthur a grammar.
What can there be in a difficulty.
Seriously in grammar.
Thinking that a little baby can sigh.
That is so much.
Sayn can say only he is dead that he is interested in what is said. That is another in
consequence.
Better and flutter must and man can beam.
Now think of seams.
Embroidery consists in remembering that it is but what she meant. There an instance of
grammar.
Suppose embroidery is two and two. There can be reflected that it is as if it were having
red about.
This is an instance of having settled it.


Grammar uses twenty in a predicament. Include hyacinths and mosses which
grow to abundance.
Grammar. In picking hyacinths quickly they suit admirably this makes
grammar a preparation. Grammar unites parts and praises. In just this way. Grammar
untiringly.
Grammar perhaps grammar.

         It is quite perplexing. It is simply an absurd string of non sequiturs, often totally
lacking grammar, meandering randomly from "topic" to "topic". It
is frustrating because there is nothing to grab onto. It is like trying to climb a mountain
made of sand.
         Stein's experiments in absurdity parallel the Dadaist and Surrealist movements of
roughly the same period, and they mark the trend away from
exuberant and laughable nonsense toward troubling and, later, macabre nonsense.
However, her work still has a freshness and silliness that makes it amusing and light
rather than disturbing and heavy.

                                          *   *   *

        As we move further into the twentieth century, we encounter the philosophy of
existentialism and the master expositor of existential malaise,
Trish-born playwright Samuel Beckett. In Beckett's most famous play, Waiting for
Godot, written in the early 1950's, the pathetic character ironically
called "Lucky" has exactly one speech, coming in about the middle of the play. He has
been being taunted by the other characters with cries of
..ink, pig!" and with sharp tugs on the rope around his neck, by which they are holding
him. Eventually he is driven beyond the breaking point, and out pours an incoherent,
wild, tormented piece of absolute confusion, resembling regurgitated academic
coursework crossed with stock phrases and garbled memorized lists of one sort and
another. Here is Lucky's famous speech:

Given the existence as uttered forth in the public works of Puncher and Wattmann of a
personal God quaquaquaqua with white beard quaquaquaqua outside time without
extension who from the heights of divine apathia divine athambia divine aphasia loves us
dearly with some exceptions for reasons unknown but time will tell and suffers like the
divine Miranda with those who for reasons unknown but time will tell are plunged in
torment plunged in fire whose fire flames if that continues and who can doubt it will fire
the firmament that is to say blast hell to heaven so blue still and calm so calm with a calm
which even though intermittent is better than nothing but not so fast and considering what
is more that as a result of the labors left unfinished crowned by the Acacacacademy of
Anthropopopometry of Essy-in-Possy of Testew and Cunard it is established beyond all
doubt all other doubt than that which clings to the labors of men that as a result of the
labors unfinished of Testew and Cunard it is established as hereinafter but not so fast for
reasons unknown that as a result of the public works of Puncher and Wattmann it is
established beyond




all doubt that in view of the labors of Fartov and Belcher left unfinished for reasons
unknown of Testew and Cunard left unfinished it is established what many deny that man
in Possy of Testew and Cunard that man in Essy that man in short that man in brief in
spite of the strides of alimentation and defecation wastes and pines wastes and pines and
concurrently simultaneously what is more for reasons unknown in spite of the strides of
physical culture the practice of sports such as tennis football running cycling gliding
conating camogie skating tennis of all kinds dying flying sports of all sorts autumn
summer winter winter tennis of all kinds hockey of all sorts penicilline and succedanea in
a word I resume flying gliding golf over nine and eighteen holes tennis of all sorts in a
word for reasons unknown in Feckham Peckham Fulham Clapham namely concurrently
simultaneously what is more for reasons unknown but time will tell fades away I resume
Fulham Clapham in a word the dead loss per head since the death of Bishop Berkeley
being to the tune of one inch four ounce per head approximately by and large more or
less to the nearest decimal good measure round figures stark naked in the stockinged feet
in Connemara in a word for reasons unknown no matter what matter the facts are there
and considering what is more much more grave that in the light of the labors lost of
Steinweg and Peterman it appears what is more much more grave that in the light the
light the light of the labors lost of Steinweg and Peterman that in the plains in the
mountains by the seas by the rivers running water running fire the air is the same and
then the earth namely the air and then the earth in the great cold the great dark the air and
the earth abode of stones in the great cold alas alas in the year of their Lord six hundred
and something the air the earth the sea the earth abode of stones in the great deeps the
great cold on sea on land and in the air I resume for reasons unknown in spite of the
tennis the facts are there but time will tell I resume alas alas on on in short in fine on on
abode of stones who can doubt it I resume but not so fast I resume the skull fading fading
fading and concurrently simultaneously what is more for reasons unknown in spite of the
tennis on on the beard the flames the tears the stones so blue so calm alas alas on on the
skull the skull the skull the skull in Connemara in spite of the tennis the labors abandoned
left unfinished graver still abode of stones in a word I resume alas alas abandoned
unfinished the skull the skull in Connemara in spite of the tennis the skull alas the stones
Cunard tennis ... the stones ... so calm ... Cunard ... unfinished .. .

        Around the same time as Beckett was writing this play, or perhaps a few years
earlier, the Welsh poet Dylan Thomas, intoxicated with the sounds of the English
language, was creating poems that are remarkably opaque. Consider the opening two
stanzas (there are five altogether) of his poem "How Soon the Servant Sun":

               How soon the servant sun,
               (Sir morrow mark),
               Can time unriddle, and the cupboard stone,
               (Fog has a bone
               He'll trumpet into meat),
               Unshelve that all my gristles have a gown
               And the naked egg stand straight,





               Sir morrow at his sponge,
               (The wound records),
               The nurse of giants by the cut sea basin,
               (Fog by his spring
               Soaks up the sewing tides),
               Tells you and you, my masters, as his strange
               Man morrow blows through food.

Poems like this make me want to cry out that this emperor has no clothes. As far as I can
discern, close to no meaning can be pulled from these lines. But how can I be sure? I
cannot. All I can say is that it would probably take such a great effort to "decode" these
lines that I suspect very, very few people would be willing to make it.

                                         *   *   *

         It is perhaps not so well known that the American singer Bob Dylan (whose name
was inspired by that of Dylan Thomas) is also an author of inspired nonsense. Some of
his nonsense written during the 1960's was collected and published in a book called
Tarantula. Its tone is often bitter and it exudes the confused mood of those difficult years.
Most of the pieces in the book consist of an outburst of free associations followed by a
letter from some strangely-named personage or other. The following sample is called "On
Busting the Sound Barrier":

   the neon dobro's F hole twang \& climax from disappointing lyrics of upstreet
   outlaw mattress while pawing visiting trophies \& prop up drifter with the bag on
   head in bed next of kin to the naked shade-a tattletale heart \& wolf of silver drizzle
   inevitable threatening a womb with the opening of rusty puddle, bottomless, a rude
   awakening \& gone frozen with dreams of birthday fog/ in a boxspring of sadly
   without candle sitting \& depending on a blemished guide, you do not feel so gross
   important/ success, her nostrils whimper. the elder fables \& slain kings \& inhale
   manners of furious proportion, exhale them against a glassy mud ... to dread misery
   of watery bandwagons, grotesque \& vomiting into the flowers of additional help to
   future treason \& telling horrid stories of yesterday's influence/ may these voices
   join with agony \& the bells \& melt their thousand sonnets now ... while the moth
   ball woman, white, so sweet, shrinks on her radiator, far away \& watches in with
   her telescope/ you will sit sick with coldness \& in an unenchanted closet ... being
   relieved only by your dark jamaican friend-you will draw a mouth on the lightbulb
   so it can laugh more freely

   forget about where youre bound youre bound for a three octave fantastic hexagram.
   you'll see it. dont worry. you are Not bound to pick wildwood flowers
    .... like i said, youre bound for a three octave titanic tantagram

                                                                        your little squirrel,
                                                                      Pety, the Wheatstraw




       Dylan is not the only popular singer of the sixties to have had a literary bent. John
Lennon, when he was in his early twenties, reveled in the nonsensical, and published two
short books called In His Own Write and A Spaniard in the Works. The books contain
mostly nonsense poetry, although there are also several prose selections. Two of
Lennon's poems will serve to illustrate his idiosyncratic style.

                                       I SAT BELONELY

                              I sat belonely down a tree,
                                humbled fat and small.
                              A little lady sing to me
                                I couldn't see at all.

                              I'm looking up and at the sky,
                                to find such wondrous voice.
                              Puzzly puzzle, wonder why,
                                I hear but have no choice.

                              "Speak up, come forth, you ravel me",
                                I potty menthol shout.
                              "I know you hiddy by this tree".
                                But still she won't come out.

                              Such softly singing lulled me sleep,
                                an hour or two or so
                              I wakeny slow and took a peep
                                and still no lady show.

                              Then suddy on a little twig
                               I thought I see a sight,
                              A tiny little tiny pig,
                               that sing with all its might.

                              "I thought you were a lady",
                                I giggle,-well I may,
                              To my suprise the lady,
                                got up-and flew away.





             THE FAULTY BAGNOSE

             Softly softly, treads the Mungle
             Thinner thorn behaviour street.
             Whorg canteell whorth bee asbin?
             Cam we so all complete,
             With all our faulty bagnose?

             The Mungle pilgriffs far awoy
             Religeorge too thee worled.
             Sam fells on the waysock-side
             And somforbe on a gurled,
             With all her faulty bagnose!

             Our Mungle speaks tonife at eight
             He tell us wop to doo
             And bless us cotten sods again
             Oamnipple to our jew
             (With all their faulty bagnose).

             Bless our gurlished wramfeed
             Me cursed cafe kname
             And bless thee loaf he eating
             With he golden teeth aflame
             Give us OUR faulty bagnose!

             Good Mungle blaith our meathalls
             Woof mebble morn so green the wheel
             Staggaboon undie some grapeload
             To get a little feel
             of my own faulty bagnose.

             Its not OUR faulty bagnose now
             Full lust and dirty hand
             Whitehall the treble Mungle speak
             We might as wealth be band
             Including your faulty bagnose

             Give us thisbe our daily tit
             Good Mungle on yer. travelled
             A goat of many coloureds
             Wiberneth all beneath unravelled
             And not so MUCH OF YER FAULTY BAGNOSE!





The first of these is transparent and charming, while the second is somewhat baffling and
disturbing. What in the world is a bagnose? No clear image comes through. And why are
all these bagnoses faulty? And does "faulty" have its normal meaning here? Hard to tell.

                                            *   *   *

The idea of "normal meanings" is turned on its head in a recent book of poetry by
William Benton called just that: Normal Meanings. One section of the book is titled
"Normal Meanings"; here is an extract from it.

Escape is, escape
       was, once more,

continued.

              Vineyards
       as dusky.


He watches it wrinkle into a school bell.

                It isn't music sometimes, I'm
       happy.


Leaves, practically falling
off and into the air.


                Hills river

       sunset ice-cream

                cone.


                        The buildings. Things
                               build up. It
                        must be so many
                               normal meanings.


The downstairs lights. Probably I doubt.





These and other

stories.
       The loveliness
of houses.

                        Clarissa is the name of the, bug I just sent somewhere.

The falseness it abjures has seemed in statements
                                we are losing.
                       It's hard to say. A note of privilege
which turns up here in their appearances.
                       I drink.

                               The cobweb is becoming a strand of
               lamplight, its black heart

                                 blessed.


                        A nice

                        Elaine

                        by the beer.

Some may find the amorphousness of this type of poetry amusing or engaging; others
may find it tiresome, confusing. I personally find it provocative for a while, but then I
begin to lose interest.

                                            *   *   *

I have somewhat greater interest in the writings of the little-known American rhetorician,
Y. Serm Clacoxia, who, in the past 25 years or so, has sporadically penned various pieces
of nonsense poetry and prose. Clacoxia's prose is marked by a certain degree of
vehemence and fire, although it is sometimes a little hard to figure out exactly what he is
ranting and raving about. Here follows one of his most lyrical tracts, entitled "The
Illusions of Alacrity".

   For millennia it has been less than appreciated how futile are the efforts of those
   who seek to sow sobriety in the furrows of trivia. To those of us who have striven
   to clarify what has been left unclear, it has proven a loss. To others who, whilst
   valiantly straddling the fine line that divides arid piquancy from acrid pungency,
   have struggled to set right the many Undeeds and Unsaids of yore, life has shown
   itself as a beast of many colors, a mountain of many flags, a hole of many anchors.




          Who, in fact, were the Outcasts of Episode, if not the champions of clarity?
  Where, indeed, were the witnesses to litany, when their fortress of fecundity was a-
  being stormed by the Ovaltine Monster, that incubus of frozen cheerios and swollen
  bananas? And dare one wonder, with the bassoon of lunacy so shrilly betoning the
  ruined fiddles of flatulism, how it is that doublethink, narcolepsy, and poseurism
  are unthreading themselves across our land like tall, statuesque, half-uneaten yet
  virtuous whippoorwhills? Can it be that a cornflake-catechism has beguiled us into
  an unsworn acceptance of never-takism?
          What sort of entiments are they, that would uncouth a mulebound lout and
  churlishly swirl his burly figure, unfurl and twirl his curly figure, hurl his whirly
  figure, into the circuline vaults of hysteresis? With a drop of sweat unroasting his
  feverish brow, we decry his fate; with the patience of a juggernaut and the
  telemachy of a dozen opossums, we lament his disparity. And summoning all the
  powers that be, we unbow the jelly of our broken dreams, dashing it with the full
  fury of a pleistocene hurdy-gurdy against the lubrified and bulbous nexus of that
  which, having doomed the dinosaurs, seeks the engulfing of all that moves.
          Thus we act; and perhaps action itself is the Anatole's Curlicue of our era. It
  is high time to recognize that action, and action alone, will be the agent that
  transmutes the flowery barrier of unutterability into an arbitrary but sacred iota of
  purposefulness, which cannot help but penetrate into an otherwise nameless and
  universally spaghettified lack of meaning, which smears and beclouds the crab-lit
  hopes of half-beings begging for deliverance from their own private, yet strangely
  tuberculine maelstroms that begat, and begotten were from, a howling sea of
  ribosomal plagiarism.

This is deliberate nonsense, of course, to be contrasted with the nondeliberate nonsense
of, say, Dylan Thomas, or the nonsense to be found in crackpot letters written to
scientists. Crackpot ideas seem to be an inevitable ingredient of any society in which
serious scientific research is carried out; there is no way to plug all the cracks, so to
speak. There is no way to ensure that only high-quality science will be done. Fortunately,
most journals do not publish absolute nonsense or gobbledygook; it is filtered out at a
very early stage. However, one journal I have come across whose pages are filled with
utter nonsense-meant seriously-is called Art-Language. To show what I mean, here are
two short excerpts from the May, 1975 issue. The first one is taken from the beginning of
an article called "Community Work". It seems, from the table of contents, to have been
written by three people collectively. The second one is taken from an article called
"Vulgar and Popular Opinions", and seems to have a single author.

  Dionysus gets a job. (Re: language has got a hold on U. S.) (It's a Whorfian conspiracy!)

  This is hopeless manqué ontological alienation which is still dealing with ideas
  about 'discovery' as a function of a metaphysics of categories. Only for researchers
  is the failure of a modal logic industry to `catch- my-experience --the birth of
  tragedy.
       Going-on in A-L indexed (somehow) is a thing-in-and-for-(dynamically)




   itself. That we never catch up with the NaturKulturLogik has little to do with the
   'actualizing' sets of the frozen dialogue ... and it's not just a ledger; our problems
   with set-theoretical axiomata are embedded into our praxis as more than just
   historical antecedents ... more than nomological permissibility ... more than
   selective filtration. We still don't recognize ourselves as very fundamental history
   producers.
            The possibility of a defence of a set, as with 'a decision', is an index-margin
   of a prima facie ersatz principle for action (!). (There is no workable distinction
   between oratio recta and oratio obliqua.) All we are left with is a deontic Drang.
   Think of that as a chain strength possibility of what, eventually, comes out as a
   product (epistemic conditions?) and the product is not a Frankfurt-ish packing-it-
   all-in .... A slogan (?) might be thought of as a free-form comprised of multiple
   structural features occurring in a (partially) given, or negotiable, unit relative to
   others. That is, the slogan is a unit in one sense or another. In. going-on
   (ideologically, perhaps), a slogan is a unitary filler-for-and-of that stretch of surf <
   surf which is in a B X S position ... But there is the critical issue of that `filler' as a
   reified function of the pusillanimous tittle-tattle of authenticity in its ellipticality (as
   a Das Volk holism) ... (e.g.) 'the Fox' material, passim, falls into that trap in dealing
   with its cultural space as a wantonly dialectical 'region' approaching the solution to
   `the negation of essence' (of homo sapiens, art or what?).

         I am tempted to quote further, to show how the wild quality of the A-L prose just
goes on and on. But life is short. It is hard for this human being to believe that these
paragraphs were meant to communicate something to anybody, but the journal appears
regularly (at least it used to), and can be found on the shelves of reputable art libraries.
Isn't it time that somebody blew the whistle? The curious thing about Art-Language is
that the collective that writes it appears to consist of people who are deeply concerned
with issues that hold much interest for me: the nature of reference, the relationship of
wholes to parts, the connection of art and reality, the structure of society, the philosophy
of set theory, the questionable existence of mathematical concepts, and so on. What is
amazing is how such concepts can be so obscured by language that it is hard to make out
anything except huge billows of very thick smoke.

                                            *   *   *

        An American poet whose work explores ground midway between nonsense and
sense is Russell Edson. He writes tiny surrealistic vignettes that shed a strange light on
life. Often he performs strange reversals, as of animate and inanimate beings, or humans
and animals. His grammar is also oblique, one of his favorite devices being to refer
repeatedly to something specific with the indefinite article "a", thus disorienting the
reader. A typical sample of Edson's style is the following, drawn from his book The Clam
Theater:





       When Science is in the Country

       When science is in the country a cow meows and the moon jumps from limb to
limb through the trees like a silver ape.
       The cow bow-wows to hear all voice of itself. The grass sinks back into the earth
looking for its mother.

        A farmer dreamed he harvested the universe, and had a barn full of stars, and a
herd of clouds fenced in the pasture.
        The farmer awoke to something screaming in the kitchen, which he identified as
the farmerette.
        Oh my my, cried the farmer, what is to become of what became?
        It's a good piece of bread and a bad farmer man, she cried.
        Oh the devil take the monotony of the field, he screamed.
        Which grows your eating thing, she wailed.
        Which is the hell with me too, he screamed.
        And the farmerette? she screamed.
        And the farmerette, he howled.

       A scientist looked through his magnifying glass in the neighborhood.

        This eerie tale leaves one with a host of unresolved images. That, of course, is
Edson's intent. And in this regard, Edson's work is quite typical. Most of the nonsense of
the twentieth century, it seems, has this deliberately upsetting quality to it, reflecting a
deep malaise. It is utterly different from the nonsense of the preceding centuries. Similar
trends exist in the other arts, particularly in music, where "classical" composers have lost
99 percent of their audience by their experimentation with randomness and cacophony.
However, the spirit of experimentation has also crept into rock music, where electronic
sounds and unusual rhythms are occasionally heard. The surrealistic, nonsensical spirit
also pervades the names of popular groups, such as "Iron Butterfly", "Tangerine Dream",
"Led Zeppelin", "Joy of Cooking", "Human Sexual Response", "Captain Beefheart",
"Brand X", "Jefferson Starship", "Average White Band", and so on.

                                         *   *   *

        Perhaps one of the virtues of nonsense is that it opens our minds to new
possibilities. The mere juxtaposition of a few arbitrary words can send the mind soaring
into imaginary worlds. It is as if sense were too mundane, and we need a breather once in
a while. Perhaps sense is also too confining. Nonsense stresses the incomprehensible face
of the universe, while sense stresses the comprehensible. Clearly both are important. Zen
teachings have striven to impart the path to "enlightenment". Although I don't believe that
such a mystical state exists, I am fascinated by the paths that are offered. Zen itself is
perhaps the archetypal source of utter nonsense. It seems fitting to





close this column with two Zen koans taken from the Mumonkan, or "Gateless Gate"-a
set of koans commented upon by the Zen master Mumon in the thirteenth century.

                         Joshu Examines a Monk in Meditation

Joshu went to a place where a monk had retired to meditate and asked him: "What is, is
what?" The monk raised his fist. Joshu replied: "Ships cannot remain where the water is
too shallow." And he left. A few days later Joshu went again to visit the monk and asked
the same question. The monk answered the same way. Joshu said: "Well given, well
taken, well killed, well saved." And he bowed to the monk.

Mumon's comment.

The raised fist was the same both times. Why is it Joshu did not admit the first and
approved the second one? Where is the fault? Whoever answers this knows that Joshu's
tongue has no bone so he can use it freely. Yet perhaps Joshu is wrong. Or, through that
monk, he may have discovered his mistake. If anyone thinks that the one's insight
exceeds the other's, he has no eyes.

              Mumon's Poem:

              The light of the eyes is as a comet,
              And Zen's activity is as lightning.
              The sword that kills the man
              Is the sword that saves the man.


              Learning is Not the Path

              Nansen said: "Mind is not Buddha. Learning is not the path."

Mumon's comment:

Nansen was getting old and forgot to be ashamed. He spoke out with bad breath and
exposed the scandal of his own home. However, there are few who appreciate his
kindness.

                      Mumon's Poem:

              When the sky is clear the sun appears,
              When the earth is parched rain will fall.
              He opened his heart fully and spoke out,
              But it was useless to talk to pigs and fish.





Post Scriptum:

        I was quite aware that I had omitted some nonsense specialists, such as James
Joyce, when I wrote this column. But there were reasons. I haven't studied Joyce, and I
feel there is a lot of complexity there. To call Joyce's strange concoctions "nonsense" is to
miss the mark.
        Several people wrote in, disappointed that I did not include anything by Walt
Kelly, the creator of "Pogo". I have to agree that Kelly was a unique writer of ingenious
and charming nonsense. In fact, I was lucky enough to grow up knowing "The Pogo Song
Book", a record of some of Kelly's most inspired silly songs, some of them belted out by
Kelly himself. One that gets across the flavor very well is this one:

                        TWIRL, TWIRL

                 Twirl! Twirl! Twinkle between!
                 The tweezers are twist in the twittering twain.
                 Twirl! Twirl! Entwiningly twirl
                 'Twixt twice twenty twigs passing platitudes plain.
                 Plunder the plover and rover rides round.
                 Ride all the rungs on the brassily bound,
                 Billy, Swirl! Swirl! Swingingly swirl!
                 Sweep along swoop along sweetly your swain.

The poem is catchy and rhythmic, and I cannot read it without hearing the song in my
head. Few people know that Kelly was a good composer of catchy melodies. But his
songs, unlike his lyrics, follow very ordinary, "sensible" rules of musical syntax.
        Two other pieces of inspired nonsense that I have run across since writing this
column are Tom Phillips' A Humument, and Luigi Serafini's Codex Seraphinianus. The
former, subtitled "A Treated Victorian Novel", was made, by a sort of literary
cannibalism, from another novel entitled A Human Document, itself written by a little-
known Victorian novelist named William Hurrell Mallock. Phillips "treated" this novel
by colorfully and imaginatively overpainting nearly all its pages, blotting out most of the
text, leaving only a select few words or letters to poke their heads through and make
cameo appearances now and then. This creation (or revelation?) of hidden messages in
someone else's text yields some very strange effects. The first page of A Humument reads
this way (I have slightly modified the two-dimensional placement of the words on the
page):

                                The following sing I a book.
                                        a book of art
                                      of mind and art
                                     that which he hid
                                          reveal I.





                                          *   *   *

        Codex Seraphinianus is a much more elaborate work. In fact, it is a highly
idiosyncratic magnum opus by an Italian architect indulging his sense of fancy to the hilt.
It consists of two volumes in a completely invented language (including the numbering
system, which is itself rather esoteric), penned entirely by the author, accompanied by
thousands of beautifully drawn color pictures of the most fantastic scenes, machines,
beasts, feasts, and so on. It purports to be a vast encyclopedia of a hypothetical land
somewhat like the earth, with many creatures resembling people to various degrees, but
many creatures of unheard-of bizarreness promenading throughout the countryside.
Serafini has sections on physics, chemistry, mineralogy (including many drawings of
elaborate gems), geography, botany, zoology, sociology, linguistics, technology,
architecture, sports (of all sorts), clothing, and so on. The pictures have their own internal
logic, but to our eyes they are filled with utter non sequiturs.
        A typical example depicts an automobile chassis covered with some huge piece of
what appears to be melting gum in the shape of a small mountain range. All over the gum
are small insects, and the wheels of the "car" appear to have melted as well. The
explanation is all there for anyone to read, if only they can decipher Serafinian.
Unfortunately, no one knows that language. Fortunately, on another page there is one
picture of a scholar standing by what is apparently a Rosetta Stone. Unfortunately, the
only language on it, besides Serafinian itself, is an unknown kind of hieroglyphics. Thus
the stone is of no help unless you already know Serafinian. Oh, well ... Many of the
pictures are grotesque and disturbing, but others are extremely beautiful and visionary.
The inventiveness that it took to come up with all these conceptions of a hypothetical
land is staggering.
        Some people with whom I have shared this book find it frightening or disturbing
in some way. It seems to them to glorify entropy, chaos, and incomprehensibility. There
is very little to fasten onto; everything shifts, shimmers, slips. Yet the book has a kind of
unearthly beauty and logic to it, qualities pleasing to a different class of people: people
who are more at ease with free-wheeling fantasy and, in some sense, craziness. I see
some parallels between musical composition and this kind of invention. Both are abstract,
both create a mood, both rely largely on style to convey content.
        Music is, in a way, a kind of nonsense that nobody really understands. It
captivates nearly every human being who can hear and yet, for all that, we still know
amazingly little about how music works its wonders. But if music is a kind of auditory
nonsense, that does not prevent there from arising even more extreme brands of auditory
super-nonsense. The works of Karl-Heinz Stockhausen, Peter Maxwell Davies, Luciano
Berio, and John Cage will provide a wonderful introduction to that genre, in case some
reader does not know what I am talking about. Especially if you like the banging of





             FIGURE 11-1. One page from David Moser´s “Metaculture” (1979)





garbage-can lids or the sound of gangland murders, their "musical offerings" are sure to
be right up your alley.
        David Moser is as fascinated with fringe-language as I am, and has explored
many uncharted regions in that territory. His longest and most adventurous journey
consisted of the writing and drawing of a roughly 40-page booklet called "Metaculture
Comics". Inspired by James Joyce, this volume contains some of the most original and
zany meaningless writings I have ever seen. It is also chock-full of the frame-breaking
and self-referential devices so beloved by modern graphic designers. A one-page sample
is shown in Figure I1-1.

                                         *   *   *

         The purpose of this column was to emphasize the very fine line that separates the
meaningful from the meaningless. It is a boundary line that has a great deal to do with the
nature of human intelligence, because the question of how meaning emerges out of
meaningless constituents when combined in certain patterned ways is still a perplexing
one. Computers are good at producing very simple passages that-to us-seem to have
meaning, and they are excellent at producing passages that are utterly devoid of meaning.
It will be interesting to see if someday a computer can tread the line and produce an
artistic exploration of meaning by producing provocative nonsense in the same way as
these human explorers of the territory have done






                             Variations on a Theme
                            as the Crux of Creativity
                                    October, 1982

                             You see things; and you say "Why?"
                             But I dream things that never were; and say "Why not?"

                                     -George Bernard Shaw in Back to Methuselah


WHEN I first heard this beautiful line it made a deep impression on me. It was in
the spring of 1968, during the presidential campaign, and Robert Kennedy had made
this line his theme. I thought it was wonderfully poetic, and I assumed he himself had
dreamt it up. Only many years later did I find out I was quite wrong: Not only had he
not made it up, but the character who utters it in the Shaw play is the snake in the
Garden of Eden! How disturbing! Why couldn't it have been the way I thought?
         "To dream things that never were"-this is not just a poetic phrase, but a truth
about human nature. Even the dullest of us is endowed with this strange ability1to
come up with counterfactual worlds and to dream. But why do we have this ability-in
fact, this proclivity? What sense does it make? And-how can one "see" what is visibly
not there?
         On my table sits a Rubik's Cube. I look at it and see a 3 X 3 X 3 cube whose
faces turn. I see-so it seems to me-what is there. But some people looked at that cube
and saw things that weren't there. They saw cubes with shaved edges, spherical
"cubes", differently coloured cubes, Magic Dominos, 2 x 2 X 2 cubes, 4 X 4 X 4 and
higher-order cubes, skew-twisting cubes, pyramids, octahedra, dodecahedra,
icosahedra, four-dimensional magic polyhedra. (See figures galore in Chapters 14 and
15.) And the list is not complete yet! Just you wait!
         How did this come about? How is it that, in looking directly at something solid
and real on a table, people can see far beyond that solidity and reality --can see an
"essence", a "core", a "theme" upon which to devise




Variations on a Theme as the Crux of Creativity                                     232

variations? I must stress that the solid cube itself is not the theme (although it is
convenient and easy to speak as if it were). In the mind of each person who perceives
a Rubik's Cube there arises a concept that we could call "Rubik's-Cubicity". It's not
the same concept in each mind, just as not everyone has the same concept of
asparagus or of Beethoven. The variations that are spun off by a given cube-inventor
are variations on that concept. In a discussion of perception and invention, this
distinction between an object and some mind's concept of the object is simple but
crucial.
         Now when Eve Rybody comes up with a new variation-let's say the 4 X4 X 4-
is it as a result of wracking her brain, trying as hard as she can to "go against the
grain", so as to come up with something original? Does she think to herself, "Golly,
that Rubik must have really exerted himself to come up with this totally new idea,
therefore I too must strain my mind to its limits in order to invent something
original."? No, no, no! A thousand times no. Einstein didn't go around wracking his
brain, muttering to himself, "How, oh how, can I come up with a Great Idea?" Like
Einstein (although perhaps on a lesser scale), Eve never needs to ask herself, "Hmm,
let's see, shall I try to figure out some way to spin off a variation on this object sitting
here in front of me?" No; she just does what comes naturally.
         The bottom line is that invention is much more like falling off a log than like
sawing one in two. Despite Thomas Alva Edison's memorable remark, "Genius is 2
percent inspiration and 98 percent perspiration", we're not all going to become
geniuses simply by sweating more or resolving to try harder. A mind follows its path
of least resistance, and it's when it feels easiest that it is most likely being its most
creative. Or, as Mozart used to say, things should "flow like oil"-and Mozart ought to
know! Trying harder is not the name of the game; the trick is getting the right concept
to begin with, so that making variations on it is like taking candy from a baby.
         Uh-oh-now I've given the cat away! So let me boldly state the thesis that I
shall now elaborate: Making variations on a theme is really the crux of creativity.

                                         *   *   *

        On the face of it, this thesis is crazy. How can it possibly be true? Aren't
variations simply derivative notions, never truly original creations? Isn't the notion of
a 4 X 4 X 4 cube simply a result of "twiddling a knob" on the concept of Rubik's-
Cubicity? You merely twist the knob from its "factory setting" of 3 to the new setting
of 4, and presto-you've got it! An inner voice protests:' "That's just too easy. That's
certainly not where Rubik's Cube, the Rite of Spring, relativity, or Romeo and Juliet
came from, is it? Isn't there a `magic spark' that leaps across a gap when a Rubik or a
Stravinsky or an Einstein or a Shakespeare comes up with a great idea, something that
is patently lacking when an Eve Rybody merely twiddles a knob on an already-
existing notion?"
Well, of course, inventing the notion of a 4 X 4 X 4 cube is far less deep




Variations on a Theme as the Crux of Creativity                                        233

than coming up with special or general relativity. I'd be the last to deny that. But that
doesn't mean that the underlying mental processes are necessarily based on totally
different principles. Of course, there is a boring sense in which the underlying mental
processes in your brain, my brain, Eve's brain, and Einstein's brain are all "the same"-
namely, they all depend on neural hardware. But it is not at such a microscopic, such
a biological level that I mean it when I suggest that the underlying mental processes in
different brains are somehow the same. What I mean is that there are mechanisms,
processes, call them what you will, that can be described functionally, without
reference to the neural substrate that enables them to take place in brains.
         Thus, a notion like "twiddling a knob on a concept" bears no relation to the
activities of neurons in the brain-or at least no obvious relation. Well then, is there
any reality to it, or is it just a metaphor? If someday we at last come to understand the
brain, will we then be confident that we're on solid ground when we speak of a brain
literally containing concepts ? Or will such statements forever remain shaky and
metaphorical fawns de parley, compared to such hard-science facts as "At the back of
each human brain there is a cerebellum"? Well, until words like "concept" have
become terms as scientifically legitimate as, say, "neuron" or "cerebellum", we will
not have come anywhere close to understanding the brain-at least not in my book.
         However, it must be admitted that at present, words like "concept" are only
metaphorical. They are protoscientific terms awaiting explication. But this is a very
good reason to try to flesh them out as much as possible, to try to see what the
metaphor of "twiddling knobs on a concept" involves. Pinning down the meaning of
such a metaphor will help us know much more clearly what we would ideally want
from a "hard-science" explanation of the brain.
         This metaphor makes your imagination conjure up a vision of a tangible thing
called a "concept" that literally has a set of knobs on it, just waiting to be twiddled.
What I picture in my mind's eye is something that, instead of being built out of
millions of neurons, is more like a metallic "black box" with a panel on it, containing
a row of plastic knobs with little pointers on them, telling you what each one's setting
is.
         Just to make this image more concrete, let me describe a genuine example of
such a black box with knobs. Back in the old days of player pianos, good pianists
made piano rolls of all sorts of wonderful music. Nowadays, you can buy phonograph
records of those rolls being played back on player pianos -but you can do better than
that. Many of the best rolls made on a special kind of piano called a Vorsetzer have
been converted into digital cassette tapes-not to be played on tape recorders, but on
pianos specially equipped with a device called a "Pianocorder". This "reads" the
magnetic tape and converts it into instructions to the keyboard and pedals, so that
your piano then plays the piece. Each Pianocorder has a black box on the front of
which is a control panel with a row of three knobs (tempo, pianissimo, and fortissimo)




Variations on a Theme as the Crux of Creativity                                      234

and one switch ("soft pedal"). By twisting the tempo knob you can make
Rachmaninoff speed up, by twiddling the pianissimo and fortissimo knobs you can
make Horowitz play more softly or Rubinstein more loudly. It's too bad, there's not a
knob labelled "pianist" so that you can select who plays. After all, it would be
interesting to change Horowitzes in midstream.

                                   *       *      *

        This device takes us one step toward realizing a dream of the unique Canadian
pianist Glenn Gould. Gould is very tuned in to the electronic age, and for years has
been advocating using computers to allow people to control the music they hear. You
begin with an ordinary recording of, say, Glenn Gould himself playing a concerto by
Mozart. But this is merely raw data for you to tamper with. On your space-age record
player, you have a bunch of knobs that allow you to slow the music down or to speed
it up ad libitum, to control the volume of all the separate sections of the orchestra,
even to correct for high notes played too flat by the violinists! In effect, you become
the conductor, with knobs to control every aspect of the performance, dynamically.
The fact that it was originally Glenn Gould at the piano is, by the time you're done
with it, irrelevant. By now you've totally taken over and made it your very own
performance! Presumably, such systems would eventually evolve to the point where
you could start with the mere written score, dispensing entirely with the acoustic
recording stage.
        But why not carry this further, then? If we are allowing ourselves to fantasize,
why not go as far as we can imagine? Why should our "raw data" be limited to the
finite universe of already-composed pieces? Why could there not be a knob to control
the mood of the composition, another to control the composer whose style it is to be
written in? This way, we could get a new piece by our favorite composer in any
desired mood. But really, this is too conservative. Why should we be limited to the
finite universe of already-born composers? Why could there not be a knob to allow us
to interpolate between composers, thus making it possible for us to tune our music-
making machine to an even mixture of Johann Sebastian Bach, Giuseppe Verdi, and
John Philip Sousa (ugh!), or a position halfway between Schubert and the Sex Pistols
(super-ugh!)? And why stop at interpolation? Why not extrapolate beyond a given
composer? For instance, I might want to hear a piece by "the composer who is to
Ravel as Ravel is to Chopin". The machine would merely need to calculate the ratios
of its knob settings 'for Ravel and Chopin, and then multiply the Ravel-settings by
those same ratios to come up with a super-Ravel.
        It's no trickier than solving any old analogy problem-you know, simple
problems like this:

  What is to a triangle as a triangle is to a square?
  What is to a honeycomb as a knight's move is to a city grid?




Variations on a Theme as the Crux of Creativity                                     235

  What is to four dimensions as the "impossible triangle" illusion is to three?
  What is to Greece as the Falkland Islands are to Britain? What is to visual art as
  fugues are to music?
  What is to a waterbed as ice is to water?
  What is to the United States as the Eiffel Tower is to France? What is to
  German as Shakespeare's plays are to English? What is to English as simplified
  characters are to Chinese? What is to 1-2-3-4-4-3-2-1 as 4 is to 1-2-3-4-5-5-4-3-
  2-1? What is to pqc as abc is to aqc?

The truth is, of course, that analogy problems are staunchly resistant to
mechanization. The knobs on most concepts are not so apparent as to allow us to just
read their settings right off. The examples above simply carried a sensible thought to a
ludicrous extreme. However, it is still worthwhile to look seriously at the idea that a
concept can be considered as a "knobbed machine" whose knobs can be twiddled to
produce a bewildering array of variations.

                                       *   *   *

        The Rubik's-Cube concept, with its "order" knob set at 3, produces an ordinary
3 X 3 X 3 cube-and with that knob set at 4, a 4 X 4 X 4. Come to think of it, doesn't
there have to be a separate knob for each dimension, so that you can twiddle each one
independently of the others? After all, not all variations have to be cubical. The Magic
Domino is 3 X 3 X 2. So if we agree that there are three knobs defining the shape,
then in the original cube they all just accidentally happened to have the same setting.
Now given these three knobs, we can use our concept-our knobbed machine-to
generate such mental objects as a 7 X 7 X 7 Rubik's Cube, a 2 X 2 X 8 Magic
Domino, even a 3 X 5 X 9 Rubik's Magic Brick (or, if you'll pardon me, a "Rubrick").
But wait a minute-if there really are just three knobs, then we're locked into three
dimensions! Obviously we don't want that. So let's add a fourth knob to control the
length in the fourth dimension. With this knob, we can now make a four-dimensional
2 X 3 X 5 X 7 Rubrick, as well as any Rubik's Tesseract that we might want. But
needless to say, once we've gone through the gate from three dimensions to four,
certainly we should expect to be able to go further. For any n, we could imagine n-
dimensional Rubik's ot4ects-for example, a 2 X 3 X 4 X 5 X 6 X 7 X 8 Hyper-
Rubrick. But now something peculiar has happened. We must now conceive of our
machine--our concept--as having a potentially unlimited number of knobs on it (one
for each dimension in n-dimensional space). If n is set to 3, there need only be 3 more
knobs. But if n is 100, we need 100 extra knobs!
        'No real machine has a variable number of knobs. Now this may sound like
a somewhat trivial observation. However, it leads into some tricky waters.




Variations on a Theme as the Crux of Creativity                                     236

The point is that, if we wis* to keep on using the metaphor of a concept as a machine
with knobs on it, we have to stretch the very concept of "knob". New knobs must be
able to sprout, depending on the settings of other knobs. Or you can think of it this
way, if you wish: on each concept, there are potentially an infinite number of knobs,
and at any moment, some new knobs may get revealed as a consequence of the
settings of other knobs.
        I´m not sure I like that view, however. It's too cut and dried, too closed and
predetermined for my tastes. I am more in favour of a view that says that the knobs on
any one concept depend on the set of concepts that happen to be awake
simultaneously in the mind of the person. This way, new knobs can spring into
existence seemingly out of nowhere; they don't all have to be present from the outset
in the isolated concept. If we go back to Rubik, this would mean that his concept of
Rubik's Cube didn't (and still doesn't) explicitly-or even implicitly-contain all the
possible variations that people may come up with. Rubik anticipated, and even
designed, many of the objects that have subsequently appeared and that we perceive
as "variations on a theme"-but certainly, his mind did not exhaust that fertile theme.
Once the concept entered the public domain, it started migrating and developing in
ways that Rubik could never have anticipated.

                                       *   *   *

        There is a way that concepts have of "slipping" from one into another,
following a quite unpredictable path. Careful observation and theorizing about such
slippages affords us perhaps our best chance to probe deeply into the hidden murk of
our conceptual networks. An example of such a slip is furnished to us whenever we
make a typo or a grammatical mistake, utter a malapropism ("She's just back from a
one-year stench at Berkeley") or a malaphor (a novel phrase concocted unconsciously
from bits and pieces of other phrases, such as "He's such an easy-go-lucky fellow" or
"Uh-oh, now I've given the cat away"), or confuse two concepts at a deeply semantic
level (e.g., saying "Tuesday" but meaning "February", or saying "midnight" in lieu of
"zero degrees"). These types of slip are totally accidental and come straight out of our
unconscious mind.
        However, sometimes a slippage can be nonaccidental yet still come from the
unconscious mind. By "nonaccidental" here, I do not mean to imply that the slip is
deliberate. It's not that we say to ourselves, "I think .1 shall now slip from one
concept into a variation of it"; indeed, that kind of deliberate, conscious slippage is
most often quite uninspired and infertile. "How to Think" and "How to Be Creative"
books-even very thoughtful ones such as George P51ya's How to Solve It-are, for that
reason, of little use to the would-be genius.
        Strange though it may sound, nondeliberate yet nonaccidental slippage
permeates our mental processes, and is the very crux of fluid thought. That is my
firmly held conviction. This subconscious manufacture of "subjunctive variations on a




Variations on a Theme as the Crux of Creativity                                     237

theme" is something that goes on day and night in each of us, usually without our
slightest awareness of it. It is one of those things that, like air or gravity or three-
dimensionality, tend to elude our perception because they define the very fabric of our
lives.
        To make this concrete, let me contrast an example of "deliberate" slippage
with an example of "nondeliberate but nonaccidental" slippage. Imagine that one
summer evening you and Eve Rybody have just walked into a surprisingly crowded
coffeehouse. Now go ahead and manufacture a few variants on that scene, in whatever
ways you want. What kinds of things do you come up with when you deliberately
"slip" this scene into hypothetical variants of itself?
        If you're like most people, you'll come up with some pretty obvious slippages,
made by moving along what seem to be the most obvious "axes of slippability".
Typical examples are:

   I could have come with Ann Yone instead of Eve Rybody.
   We could have gone to a pancake house instead of a coffeehouse.
   The coffeehouse could have been nearly empty instead of full.
   It could have been a winter's evening instead of a summer's evening.

        Now contrast your variations with one that I overheard one evening this past
summer in a very crowded coffeehouse, when a man walked in with a woman. He
said to her, "I'm sure glad I'm not a waitress here tonight!" This is a perfect example
of a subjunctive variation on the given theme-but unlike yours, this one was made
without external prompting, and it was made for the purposes of communication to
someone. The list above looks positively mundane next to this casually tossed-off
remark. And the remark was not considered to be particularly clever or ingenious by
his companion. She merely agreed with the thought by saying "Yeah." It caught my
attention not so much because I thought it was clever, but mostly because I am always
on the lookout for interesting examples of slippability.
        I found this example not just mildly interesting, but highly provocative. If you
try to analyze it, it would appear at first glance to force you as listener to imagine a
sex-change operation performed in world record time. But when you simply
understand the remark, you see that in actuality, there was no intention in the
speaker's mind of bringing up such a bizarre image. His remark was much more
figurative, much more abstract. It was based on an instantaneous perception of the
situation, a sort of "There-but-for-thegrace-of-God-go-I" feeling, which induces a
quick flash to the effect of "Simply because I am human, I can place myself in the
shoes of that harried waitress-therefore I could have been that waitress." Logical or
not, this is the way our thoughts go.
        So when you look carefully, you see that this particular thought has practically
nothing to do with the speaker, or even with the waitresses he sees. It's just his flip
way of saying, "Hmm, it sure is busy here tonight." And




Variations on a Theme as the Crux of Creativity                                     238

that's of course why nobody really is thrown for a loop by such a remark Yet it was
stated in such a way that it invites you to perform a "light" mapping of him onto a
waitress, just barely noticing (if at all) that there is a sex difference. What an
amazingly subtle thought process is involved here!
        And what is even more amazing (and frustrating) to me is how hard it is to
point out to people how amazing it is! People find it very hard indeed to see what's
amazing about the ordinary behavior of people. They cannot quite imagine how it
might have been otherwise. It is very hard to slip mentally into a world in which
people would not think by slipping mentally into other worlds-very hard to make a
counterfactual world in which counterfactuals were not a key ingredient of thought.
        Another quick example: I was having a conversation with someone who told
me he came from Whiting, Indiana. Since I didn't know where that was, he explained,
"Whiting is very near Chicago-in fact, it would be in Illinois if it weren't for the state
line." Like the earlier one, this remark was dropped casually; it was certainly not an
effort to be witty. He didn't chuckle, nor did I. I simply flashed a quick smile,
signaling my understanding of his meaning, and then we went on. But try to analyze
what this remark means! On a logical level, it is somewhat like a tautology. Of course
Whiting would be in Illinois if the Illinois state line made it be so-but if that's all he
meant, it is an empty remark, because it holds just as well for cities thousands of miles
from Chicago. But clearly, the notion he had in mind was that there is an accidental
quality to where boundary lines fall, a notion that there are counterfactual worlds
"close" to ours, worlds in which the Illinois-Indiana line had gotten placed a couple of
miles further east, and so on. And his remark tacitly assumed that he and I shared such
intuitions about the impermanence and arbitrariness of geographical boundary lines,
intuitions about how state lines could "slip".
        Remarks like this betray the hidden "fault lines of the mind"; they show which
things are solid and which things can slip. And yet, they also reveal that nothing is
reliably unslippable. Context contributes an unexpected quality to the knobs that are
perceived on a given concept. The knobs are not displayed in a nice, neat little control
panel, forevermore unchangeable. Instead, changing the context is like taking a tour
around the concept, and as you get to see it from various angl s, more and more of its
knobs are revealed. Some people get to be good at perceiving fresh new knobs on
concepts where others thought there were none, just as some people get to be good at
perceiving mushrooms in a forest where others see none, even when they stare
mightily.

                                        *   *   *

        It may still be tempting to think that for each well-defined concept, there must
be an "ultimate" or "definitive" set of knobs such that the abstract space traced out by
all possible combinations of the knobs yields all possible




Variations on a Theme as the Crux of Creativity                                       239

instantiations of the concept. A case in point is the concept of the letter `A'. The
typographically naive might think that there are four or five knobs to twiddle here,
and that's all. However, the more you delve into letter forms, the more elusive any
attempt to parametrize them mathematically becomes. .One of the most valiant efforts
at "knobbifying the alphabet" has been the letterform-defining system called
"Metafont", developed at Stanford by the well-known computer scientist Donald
Knuth.
         Knuth's purpose is not to give the ultimate parametrization of the letters of the
alphabet (indeed, I suspect that he would be the first to laugh at the very notion), but
to allow a user to make "knobbed letters"-we could call them letter schemas. This
means that you can choose for yourself what the variable aspects of a letter are, and
then, with Metafont's aid, you can easily construct knobs that allow those aspects to
vary. This includes just about anything you can think of: stroke lengths, widenings or
taperings of strokes, curvatures, the presence or absence of serifs, and so on. The full
power of the computer is then at your disposal; you can twiddle away to your heart's
desire, and the computer will generate all the products your knob-settings define.
         Going further than letters in isolation, Knuth then allowed letters to share
parameters-that is, a single "master knob" can control a feature common to a group of
related letters. This way, although there may be hundreds of knobs when you count
the knobs on all the control panels of all the letters of the alphabet, there will be a far
smaller number of master knobs, and they will have a deeper and more pervasive
influence on the whole alphabet. What happens, in effect, is that by twiddling the
master knobs alone, you have a way of drifting smoothly through a space of
typefaces.
         Perhaps Knuth's greatest virtuoso trick yet with Metafont is what he did with
Psalm 23, which in this version consists of 593 characters. (See Figure 12-1.) Knuth
had defined a full set of letters that shared 28 "master knobs". He began his printed
version of the psalm with all 28 master knobs at their leftmost settings. Then, letter by
letter, he inched his way toward the rightmost settings, turning each knob 1/592 of the
way, so that by the time he had reached the final letter, the extreme opposite end of
the spectrum had been attained. In one sense, every letter in this version of the psalm
is printed in a different typeface! And yet the transition is so smooth as to be locally
undetectable even to a finely trained eye. This example is drawn from Knuth's
inspiring article in Visible Language entitled "The Concept of a Meta-Font".
         One of Knuth's main theses is that with computers, we now are'in the position
of being able to describe not just a thing in itself, but how that thing would vary.
Metafont epitomizes this thesis. In a sense, the computer, rather than simply blindly
reproducing fixed letter shapes, has a crude "understanding" of what it. is drawing,
created by the designer who "knobbified" the letters. And yet, one should be careful
not to fall under the illusion, so easily created by Metafont's extraordinary power, that
these




Variations on a Theme as the Crux of Creativity                                        240

FIGURE 12-1. Psalm 23, printed by Donald Knuth's ME TA FONT program. It starts
out in an old-fashioned, highly serifed typeface and gradually modulates into a
modernistic, sans-serif typeface. Each step, imperceptible on its own, is accomplished
by making a tiny shift in 28 parameters governing the overall appearance of the
computerized alphabet.

28 master knobs-or any finite set of knobs-might actually span the entire space of all
possible typefaces. This is about as far from the truth as would be the claim that the
space of all possible face types (see Figure 12-2) could be captured in a computer
program with 28 knobs.
         Even the space of all versions of the letter `A' is only barely explored when
you twiddle all the knobs in Knuth's representation of `A'-not just the 28 master knobs
it shares with other letters, but the many "private" knobs it has as well. Even a
thousand knobs would not suffice to cover the variety of letter 'A's that people
recognize easily. Some evidence of the richness of the 'A' concept is shown in Figure
12-3. These `A's are all taken from real typefaces in the 1982 Letraset Catalogue. To
illustrate that such richness is not a quirk of our writing system, I have assembled, in
Figure 12-4, a similar collection of variants of the Chinese character meaning "black"




Variations on a Theme as the Crux of Creativity                                     241

Variations on a Theme as the Crux of Creativity   242

Variations on a Theme as the Crux of Creativity   243

(pronounced "hei", rhyming with `a'). I found them in some Chinese language
graphic-design catalogues. This figure is a real eye-opener for people who don't read
Chinese. They usually ask incredulously, "You mean Chinese people can easily tell
that these are all the same character?!" Of course they can, and in a split second just
as we can for the matrix of' A's.
There is a crucial distinction to be drawn here. A machine with one off on switch (the
most trivial kind of knob) for each square in a 500 X 500 grid will certainly define
any of the 'A's shown-but it will not exclude `B's or hei 's or pictures of your
grandmother or of trolley cars. It is another matter altogether to define a set of knobs
whose twiddling covers all the `A's, showing all the interpolations between them (as
well as extrapolations in all possible directions)-yet never leads you out of the space
of recognizable




Variations on a Theme as the Crux of Creativity                                     244

'A's. This is far trickier! Similarly, it is a nearly trivial project to write a computer
program that in theory writes all possible sequences and combinations of tones in all
possible rhythmic patterns-but that is a far cry from writing a program that produces
only pieces in the style of Bach. Putting on the constraints makes the program
unutterably more complex!
        What Metafont gives you, rather than the full space of all typefaces or 'A's, is a
subspace, and such a tightly related subspace that it is perhaps best to call it a family.
Nobody would be able to predict butterflies from having studied ants and wasps and
beetles. Certainly no currently imaginable program would, anyway. Likewise, nobody
would be able to predict the full magnitude of the concept of 'A', from seeing only the
family traced out by the finite number of knobs in any realistic Metafont program for
`A'.
        The next stage beyond Metafont will be a program that, on its own, can extract
a set of knobs from a set of given input letters. This, however, is a program for the
distant future. At present, it takes a highly trained and perceptive typeface designer
months to convert a set of letterforms into Metafont programs with knobs flexible
enough to warrant the trouble taken. It would be relatively easy to do it in some crude
mechanical way, but what one wants is for stylistic unity to be preserved even as the
master knobs are twiddled-and therefore, the task of automating the production of
Metafont programs amounts to automation of artistic perception. It's not just around
the corner.

                                        *   *   *

        There is a curious book called One Book Five Ways, published in 1978 by
William Kaufmann, Inc. It came about this way. As an educational experiment in
comparative publishing procedures, a manuscript on indoor gardening was sent
around to five different university presses, and they all cooperated in coming up with
full publication versions of the book, which turned out to be stunningly different at all
conceivable levels. William Kaufmann had the bright idea of publishing pieces of the
various versions side by side; what resulted was this elegant "metabook". It brings
home the meaning of the old saying that there's more than one way to skin a cat.
        Making this book was an extravagant foray into "possible worlds", the kind of
thing that seems very hard to do. One of Knuth's points, however, is that as computers
become more sophisticated and common, the notion of skinning a cat in nine different
ways will gradually become less extravagant. Once your "cat" has been represented
inside a powerful computer program, it is no longer just one cat; it has become,
instead, a "cat-schema"-a mold for many cats at once, and you can skin them all
differently (or at least until the cat-schema runs out of lives).
        Text formatters and computer typesetting present us easily with many
alternative versions of a piece of text. Metafont shows us how letterforms can glide
into alternative versions of themselves. It is now up to us to




Variations on a Theme as the Crux of Creativity                                       245

FIGURE 12-5. In (a), a stylized implicosphere. In (b) through (d), various degrees of
overlap of two implicospheres are portrayed. Too much overlap (b) leads to mushy,
sloppy thought, while too little overlap (d) leads to sparse, dull thought. The ideal
amount of overlap and autonomy (c) leads to creative, insightful thought. -
       In (e), a related and charming geometrical problem called "Mrs. Miniver's
problem "is shown. The idea is to determine the conditions under which the overlap of
two circles (representing two people) has the same area as each of the two crescents
formed. Mrs. Miniver wishes thereby to symbolize her vision of the ideal romance.
The ideal overlap of course symbolizes how much two lovers ideally have in common.


Variations on a Theme as the Crux of Creativity                                  246

continue this trend of extending our abilities to see further into the space of
possibilities surrounding what is. We should use the power of computers to aid us in
seeing the full concept-the implicit "sphere of hypothetical variations"-surrounding
any static, frozen perception.
        I have concocted a playful name for this imaginary sphere: I call it the
implicosphere, which stands for implicit counterfactual sphere, referring to things that
never were but that we cannot help seeing anyway. (The word can also be taken as
referring to the sphere of implications surrounding any given idea. A visual
representation of an implicosphere is shown in Figure 12-5.) If we wish to enlist
computers as our partners in this venture of inventing variations on a theme, which is
to say, turning implicospheres into "explicospheres", we have to give them the ability
to spot knobs themselves, not just to accept knobs that we humans have spotted. To
do this we will have to look deeply into the nature of "slippability", into the fine-
grained structure of those networks of concepts in human minds.

                                        *   *   *

        One way to imagine how slippability might be realized in the mind is to
suppose that each new concept begins life as a compound of previous concepts, and
that from the slippability of those concepts, it inherits a certain amount of slippability.
That is, since any of its constituents can slip in various ways, this induces modes of
slippage in the whole. Generally, letting a constituent concept slip in its simplest ways
is enough, since when more than one of these is done at a time, that can already create
many unexpected effects. Gradually, as the space of possibilities of the new concept-
the implicosphere-is traced out, the most common and useful of those slippages
become more closely and directly associated with the new concept itself, rather than
having to be derived over and over from its constituents. This way, the new concept's
implicosphere becomes more and more explicitly explored, and eventually the new
concept becomes old and reaches the point where it too can be used as a constituent of
fresh new young concepts.
        Some examples of this sort of thing were presented in my column for
September, 1981 (Chapter 23). Now although September is almost October




Variations on a Theme as the Crux of Creativity                                        247

and 1981 is almost 1982, that doesn't quite mean that you have those examples at your
mind's fingertips, or on the tip of your mind's tongue. So let me present a few more
examples of slippage of a new notion based on slipping some of its parts in their
simplest ways. The notion I have chosen is that of yourself sitting there, reading this
very column at this very moment. Here are some elements of the implicosphere of
that concept:

       You are almost reading the September 1981 issue of Scientific American.
       You are almost reading a piece by Richard Hofstadter, the historian.
       You are almost reading a column by Martin Gardner.
       Your identical twin is almost reading this column.
       You are almost reading this column in French.
       You are almost reading Gödel, Escher, Bach.
       You are almost reading a letter from me.
       You are almost writing this column.
       You are almost hearing my voice.
       I am almost talking to you.
       You are almost ready to throw this copy of Mad magazine out in disgust.

         By now, the original concept is almost lost in a silly sea of "almost"
variations-but it has been enriched by this exploration, and when you come back to it,
it will have been that much more reified as a stand-alone concept, a single entity
rather than a compound entity. After a while, under the proper triggering
circumstances, this very example may be retrieved from memory as naturally and
effortlessly as the concept of "fish" is.
         This is an important idea: the test of whether a concept has really come into its
own, the test of its genuine mental existence, is its retrievability by that process of
unconscious recall. That's what lets you know that it has been firmly planted in the
soil of your mind. It is not whether that concept appears to be "atomic", in the sense
that you have a single word to express it by. That is far too superficial.
         Here is an example to illustrate why. A friend told me recently that the
Encyclopaedia Britannica's first edition (1768-71) consisted of three volumes:
Volume 'I: "A-B"; Volume II: "C-L", and then Volume III: the rest of the alphabet. In
that edition, 511 pages were devoted to topics beginning with `A', while the last
volume had 753 pages altogether! (I guess that in those days there weren't yet many
interesting things around that began with letters between `M' and 'Z'.) Hearing this
amusing fact instantaneously triggered the retrieval of the memory, implanted in me
years and years ago under totally unremembered circumstances, of how records used
to be made, back in the days when there was no magnetic tape and the master disk
was actually cut during the live performance. The performers would be playing along
and all of a sudden the recording engineer would notice that there wasn't much room
left on the plate, so the performers would be given a signal to hurry up, and as a
result, the tempo would be faster and faster




Variations on a Theme as the Crux of Creativity                                       248

the further toward the center the needle came. I think it is obvious why the - one
triggered retrieval of the other. And yet-is it obvious?
        On the surface, these two concepts are completely unrelated. One concerns
printed matter, books, the alphabet, and so on, while the other concerns plastic disks,
sounds, performers, recording techniques, and so on. However, at some deeper
conceptual level, these really are the same idea. There is just one idea here, and this
idea I call a conceptual skeleton. Try to verbalize it. It's certainly not just one word. It
will take you a while. And when you do come up with a phrase, chances are it will be
awkward and stilted-and still not quite right!
        Both of the cited instances of this conceptual skeleton-in itself nameless,
majestically nonverbalizable-are floating about in the implicosphere that surrounds it,
along with numerous other examples that I am unaware of, not yet having twiddled
enough knobs on that concept. I don't yet even know which knobs it has! But I may
eventually find out. The point is that the concept itself has been reed-this much is
proven by the fact that it acts as a point of immediate reference; that my memory
mechanisms are capable of using it as an "address" (a key for retrieval) under the
proper circumstances. The vast majority of our concepts are wordless in this way,
although we can certainly make stabs at verbalizing them when we need to.


                                     *       *       *

         Early in this column, I stated a thesis: that the crux of creativity resides in the
ability to manufacture variations on a theme. I hope now to have sufficiently fleshed
out this thesis that you understand the full richness of what I meant when I said
"variations on a theme". The notion encompasses knobs, parameters, slippability,
counterfactual conditionals, subjunctives, "almost"-situations, implicospheres,
conceptual skeletons, mental reification, memory retrieval-and more.
         The question may persist in your mind: Aren't variations on a theme somehow
trivial, compared to the invention of the theme itself? This leads one back to that
seductive notion that Einstein and other geniuses are "cut from a different cloth" from
ordinary mortals, or at least that certain cognitive acts done by them involve
principles that transcend the everyday ones. This is something I do not believe at all.
If you look at the history of science, for instance, you will see that every idea is built
upon a thousand related ideas. Careful analysis leads one to see that what we choose
to call a new theme is itself always some sort of variation, on a deep level, of previous
themes. The trick is to be able to see the deeply hidden knobs!
         Newton said that if he had seen further than others, it was only by standing on
the shoulders of giants. Too often, however, we simply indulge in wishful thinking
when we imagine that the genesis of a clever or beautiful idea was somehow due to
unanalyzable, magical, transcendent insight rather than to




Variations on a Theme as the Crux of Creativity                                        249

any mechanisms-as if all mechanisms by their very nature were necessarily shallow
and mundane.
        My own mental image of the creative process involves viewing the
organization of a mind as consisting of thousands, perhaps millions, of overlapping
and intermingling implicospheres, at the center of each of which is a conceptual
skeleton. The implicosphere is a flickering, ephemeral thing, a bit like a swarm of
gnats around a gas-station light on a hot summer's night, perhaps more like an
electron cloud, with its quantummechanical elusiveness, about a nucleus, blurring out
and dying off the further removed from the core it is (Figure 12-5). If you have
studied quantum chemistry, you know that the fluid nature of chemical bonds can best
be understood as a direct consequence of the curious quantummechanical overlap of
electronic wave functions in space, wave functions belonging to electrons orbiting
neighboring nuclei. In a metaphorically similar way, it seems to me, the crazy and
unexpected associations that allow creative insights to pop seemingly out of nowhere
may well be consequences of a similar chemistry of concepts with its own special
types of "bonds" that emerge out of an underlying "neuron mechanics".
        Novelist Arthur Koestler has long been a champion of a mystical view of
human creativity, advocating occult views of the mind while at the same time
eloquently and objectively describing its workings. In his book The Act of Creation,
he presents a theory of creativity whose key concept he calls "bisociation"-the
simultaneous activation and interaction of two previously unconnected concepts. This
view emphasizes the comingtogether of two concepts, while bypassing discussion of
the internal structure of a single concept. In Koestler's view, something new can
happen when two concepts "collide" and fuse-something not present in the concepts
themselves. This is in keeping with Koestler's philosophy that wholes are somehow
greater than the sum of their parts.
        By contrast, I have been emphasizing the idea of the internal structure of one
concept. In my view, the way that concepts can bond together and form conceptual
molecules on all levels of complexity is a consequence of their internal structure.
What results from a bond may surprise us, but it will nonetheless always have been
completely determined by the concepts involved in the fusion, if only we could
understand how they are structured. Thus the crux of the matter is the internal
structure of a single concept and how it "reaches out" toward things it is not. The crux
is not some magical, mysterious process that occurs when two indivisible concepts
collide; it is a consequence of the divisibility of concepts into subconceptual elements.
As must be clear from this, I am not one to believe that wholes elude description in
terms of their parts. I believe that if we come to understand the "physics of concepts",
then perhaps we can derive from it a "chemistry of creativity", just as we can derive
the principles of the chemistry of atoms and molecules from those of the physics of
quanta and particles. But as I said earlier, it is not just around the corner. Mental
bonds will probably turn




Variations on a Theme as the Crux of Creativity                                      250

out to be no less subtle than chemical bonds. Alan Turing's words of cautious
enthusiasm about artificial intelligence remain as apt now as they were in 1950, when
he wrote them in concluding his famous article "Computing Machinery and
Intelligence": "We can only see a short distance ahead, but we can see plenty there
that needs to be done."
        Recently I happened to read a headline on the cover of a popular electronics
magazine that blared something about "CHIPS THAT SEE". Bosh! I'll start believing
in "chips that see" as soon as they start seeing things that never were, and asking
"Why not?"

Post Scriptum.
                                                               Knobs, knobs, everywhere
                                                                Just vary a knob to think.

         Some readers objected to the slogan of this column-that making variations on
a theme is the crux of creativity. They felt-and quite rightly -that making variations
(i.e., twisting knobs) is as easy as falling off a log. So how can genius be that easy?
Part of the answer is: For a genius, it is easy to be a genius. Not being a genius would
be excruciatingly hard for a genius. However, this isn't a completely satisfactory
answer for people who pose this objection. They feel that I am unwittingly implying
that it is easy for anybody to be a genius: after all, a crank can crank a knob as deftly
as a genius can. The crux of their objection, then, is that the crux of creativity is not in
twiddling knobs, but in spotting them!
         Well, that is exactly what I meant by my slogan. Making variations is not just
twiddling a knob before you; part of the act is to manufacture the knob yourself.
Where does a knob come from? The question amounts to asking: How do you see a
variable where there is actually a constant ? More specifically: What might vary, and
how might it vary? It's not enough to just have the desire to see something different
from what is there before you. Often the dullest knobs are a result of someone's
straining to be original, and coming up with something weak and ineffective. So
where do good knobs come from? I would say they come from seeing one thing as
something else. Once an abstract connection is set up via some sort of analogy or
reminding-incident, then the gate opens wide for ideas to slosh back and forth
between the two concepts.
         A simple example: A friend and I noticed a fuel-delivery truck pulling into a
driveway, and on it was very conspicuously printed "NSF", standing for "North Shore
Fuel". However, to us those letters meant "National Science Foundation" as surely as
"TNT" means "trinitrotoluene" to Eve Rybody. Now, we could have just let the
coincidence go, but instead we played with it. We envisioned a National Science
Foundation truck pulling up to a




Variations on a Theme as the Crux of Creativity                                        251

research institute. The driver gets out of the cab, drags a thick flexible hose over to a
hole in the wall of a building and inserts it, then starts up a loud motor, and pumps a
truckload of money-presumably in large bills-into the cellar of the building. (Wouldn't
it be nice if grants were delivered that way?) This vision then led us to pondering the
way that money actually does flow between large institutions: usually as abstract,
intangible numbers shot down wires as binary digits, rather than as greenbacks hauled
about in large trucks.
        This very small incident serves well to illustrate how a simple reminding-
incident triggered a series of thoughts that wound up in a region of idea-space that
would have been totally unanticipable moments before. All that was needed was for
an inappropriate meaning of "NSF" to come to mind, and then to be explored a bit.
Such opportunities for being reminded of something remote-such double-entendre
situations-occur all the time, but often they go unobserved. Sometimes the ambiguity
is observed but shrugged off with disinterest. Sometimes it is exploited to the . hilt. In
this example, the result was not earthshaking, but it did cast things in a new light for
both of us, and the image amused us quite a bit. And this way of exploiting
serendipity-that is, exploiting coincidences and unexpected perceived similarities-is
typical of what I consider the crux of the creative process.

                                        *   *   *

        Serendipitous observation and quick exploration of potential are vital elements
in the making of a knob. What goes hand in hand with the willingness to playfully
explore a serendipitous connection is the willingness to censor or curtail an
exploration that seems to be leading nowhere. It is the flip side of the risk-taking
aspect of serendipity. It's fine to be reminded of something, to see an analogy or a
vague connection, and it's fine to try to map one situation or concept onto another in
the hopes of making something novel emerge-but you've also got to be willing and
able to sense when you've lost the gamble, and to cut your losses. One of the
problems with the ever-popular self-help books on how to be creative is that they all
encourage "off-the-wall" thinking (under such slogans as "lateral thinking",
"conceptual blockbusting", "getting whacked on the head", etc.) while glossing over
the fact that most off-the-wall connections are of very little worth and that one could
waste lifetimes just toying with ideas in that way. One needs something much more
reliable than a mere suggestion to "think zany, out-of-the-system thoughts".
        Frantic striving to be original will usually get you nowhere. Far better to relax
and let your perceptual system and your category system work together
unconsciously, occasionally coming up with unbidden connections. At' that point,
you-the lucky owner of the mind in question-can seize the bpportunity and follow out
the proffered hint. This view of creativity has the




Variations on a Theme as the Crux of Creativity                                       252

conscious mind being quite passive, content to sit back and wait for the unconscious
to do its remarkable broodings and brewings.
        The most reliable kinds of genuine insight come not from vague reminding
experiences (as with the letters "NSF"), but from strong analogies in which one
experience can be mapped onto another in a highly pleasing way. The tighter the fit,
the deeper the insight, generally speaking. When two things can both be seen as
instances of one abstract phenomenon, it is a very exciting discovery. Then ideas
about either one can be borrowed in thinking about the other, and that sloshing-about
of activity may greatly illumine both at once. For instance, such a connection (i.e.,
mapping)-between sexism and racism-resulted in my "Person Paper" (Chapter 8).
Another example is Scott Kim's brilliant article "Noneuclidean Harmony", in which
mathematics and music are twisted together in the most amazing ways. It can be
found in The Mathematical Gardner, an anthology dedicated to Martin Gardner,
edited by David Klarner.
        A mapping-recipe that often yields interesting results is projection of oneself
into a situation: "How would it be for me?" This can mean a host of things, depending
on how you choose to inject yourself into the scene, which is in turn determined by
what grabs your attention. The man who focused in on the bustling activity in the
coffeehouse and said, "I'm sure glad I'm not a waitress here tonight!" might instead
have been offended by the sounds reaching his ears and said, "If I were the owner
here, I'd play less Muzak" -or he might have zeroed in on someone purchasing a
brownie and said, "I wish I were that thin." People are remarkably fluid at seeing
themselves in roles that they self-evidently could never fill, and yet the richness of the
insights thus elicited is beyond doubt.

                                        *       *      *

        When I first heard the French saying Plus ça change, plus c'est la meme chose,
it struck me as annoyingly nonsensical: "The more it changes, the samer it gets" (in
my own colloquial translation). I was not amused but nonetheless it stuck in my mind
for years, and finally it dawned on me that it was full of meanings. My favorite way
of interpreting it is this. The more different manifestations you observe of one
phenomenon, the more deeply you understand that phenomenon, and therefore the
more clearly you can see the vein of sameness running through all those different
things. Or put another way, experience with a wide variety of things refines your
category system and allows you to make incisive, abstract connections based on deep
shared qualities. A more cynical way of putting it, and probably more in line with the
intended meaning, would be that superficially different things are often boringly the
same. But the saying need not be taken cynically.
        Seeing clear to the essence of something unfamiliar is often best achieved by
finding one or more known things that you can see it as, then being able to balance
these views. Physicists have long since learned to juggle two views




Variations on a Theme as the Crux of Creativity                                       253

of light: light as waves, light as particles. They know that each contains a grain of the
essence of light, that neither contains it all, and they know when to think of light
which way. Don't be fooled by people who knowingly assure you that physicists don't
depend on crude images or analogies as crutches, that everything they need is
contained in their formulas. The fallacy here
        . is that which formula to apply, how to apply it, and what parts of it to neglect
are all aspects not covered in any formula, which is why doing physics is a great art,
despite the fact that there are formulas all over the place for Eve Rybody and her
brother to use.
        Seeing anything as waves suggests immediate knobs: wavelength, frequency,
amplitude, speed, medium, and a host of other basic notions that define the essence of
undularity. Seeing anything as particles suggests totally different knobs: mass, shape,
radius, rotation, constituents, and a host of other basic notions that define the essence
of corpuscularity. If you choose to see, say, people as waves or as particles, you may
find some of these suggested knobs quite interesting. On the other hand, it may not be
`fruitful to do so. Good analogies usually are not the product of an off-the-wall
suggestion like this, but spring to mind unbidden, from the deep similarity-searching
wells of the unconscious.
        Once you have decided to try out a new way of viewing a phenomenon, you
can let that view suggest a set of knobs to vary. The, act of varying them will lead you
down new pathways, generating new images ripe for perception in their own right.
This sets up a closed loop:

       • fresh situations get unconsciously framed in terms of familiar concepts;
       • those familiar concepts come equipped with standard knobs to twiddle;
       • twiddling those knobs carries you into fresh new conceptual territory.

        A visual image that I always find coming back in this context is that of a
planet orbiting a star, and whose orbit brings it so close to another star that it gets
"captured" and begins orbiting the second star. As it swings around the new star,
perhaps it finds itself coming very close to yet another star, and ficklely changes
allegiance. And thus it do-si-do's its way around the universe.
        The mental analogue of such stellar peregrinations is what the loop above
attempts to convey. You can think of concepts as stars, and knob-twiddling as
carrying you from one point on an orbit to another point. If you twiddle enough, you
may well find yourself deep within the attractive zone of an unexpected but
interesting concept and be captured by it. You may thus migrate from concept to
concept. In short, knob-twiddling is a device that carries you from one concept to
another, taking advantage of their overlapping orbits.
        Of course, all this cannot happen with a trivial model of concepts. We see




Variations on a Theme as the Crux of Creativity                                       254

it happening all the time in minds, but to make it happen in computers or to locate it
physically in brains will require a fleshing-out of what concepts really are. It is fine to
talk of "orbits around concepts" as a metaphor, but developing it into a full scientific
notion that either can be realized in a computer model or can be located inside a brain
is a giant task. This is the task that faces cognitive scientists if they wish to make
"concept" a legitimate scientific term. This goal, suggested at the start of this article,
could be taken to be the central goal of cognitive science, although such things are
often forgotten in the inane hoopla that is surrounding artificial intelligence more and
more these days.
        The cycle shown above spells out what I intend by the phrase "making
variations on a theme", and it is this loop that I am suggesting is the crux of creativity.
The beauty of it is that you let your memory and perceptual mechanisms do all the
hard work for you (pulling concepts from dormancy); all you do is twiddle knobs.
And I'll let you decide what this odd distinction is between something called "you"
and the hard-working mechanisms of "your memory".

                                        *   *   *

         The concept of the "implicosphere" of an idea-the sphere of variations on it
resulting from the twiddling of many knobs a "reasonable" amountis a difficult one,
but it is absolutely central to the meaning of this column. One way of thinking about it
is this. Imagine a single gnat attracted by a bright light. It will buzz about, tracing out
a three-dimensional random walk centered on that light. If you keep a photographic
plate exposed so that you can record its path cumulatively, you will first see a chaotic
broken line, but soon the image will get so dense with criss-crossing lines that it will
gradually turn into a circular smear of slowly increasing radius. At the outer edges of
the smear you might once in a while make out an occasional foray of the lone bug.
For a while, the territory covered expands, but eventually this gnat-o-sphere will reach
a stable size. Its silhouette, instead of being a sharp-edged circle, will be a blurry
circle (see Figure 12-5a) whose approximate radius reveals something about how
gnats are attracted by lights.
         Now if you simply think of this translated into idea-space, you have roughly
the right image. Of course, not all implicospheres have the same radius. Some
people's implicospheres tend to have bigger radii than other people's do, and
consequently their implicospheres overlap more. This can be good but it can be
overdone. Too much overlap (Figure 12-5b) and all you have is a mush of vaguely
associated ideas, an overdone and tasteless mental goulash. Too little overlap (Figure.
12-5d) and you have a very thin, watery mind, one with few big surprises (except for
the meta-level surprise of having so few surprises). There is, in other words, an
optimum amount of overlap for useful creative insight (Figure 12-5c). This is the kind
of thing




Variations on a Theme as the Crux of Creativity                                        255

that cannot be taught, however. It would be like trying to train a gnat to control the
size of the spheres it traces out. Or if you prefer, it would be like trying to train an
entire swarm of gnats to form spheres of a particular size whenever they cluster
around lamps. The problem is, it is already preprogrammed in gnats how much they
are attracted by lights, by each other, and so on.
        In my view, mindpower is a consequence of how implicospheres in idea-space
emerge from the statistical predispositions of neurons to fire in response to each other.
Such deep statistical patterns of each brain cannot be altered, although of course a few
superficial aspects can be altered. You can teach somebody to think of applehood
whenever they think of mother pie, for instance-but adding any number of specific
new associative connections does not have any effect on the underlying statistics of
how their neurons work. So in that sense I am gravely doubtful about courses or
books that promise to improve your thinking style or capabilities. Sure, you can add
new ideas-but that's a far cry from adding pizzazz. The mind's -'perceptual and
category systems are too much at the "subcognitive" level to be reached via cognitive-
level training techniques. If you are old enough to be reading this book, then your
deep mental hardware has been in place for many years, and it is what makes your
thinking-style idiosyncratic and recognizably "you". (If you are not, then what are you
doing reading this book? Put it down immediately!) For more on the ideas of
subcognition and identity, see Chapters 25 and 26.
        When a new idea is implanted in a mind, an implicosphere grows around it.
Since this means, in essence, the linking-up of this new idea with older ideas, I call it
"diffusion in idea-space". My canonical example of this phenomenon, although it is a
rather grim one, has to do with the recent spate of random murders inspired by the
spiking of Tylenol capsules with strichnine. It was the Food and Drug
Administration's response that so intrigued me, because it implicitly revealed a theory
of how this idea would diffuse in the idea-space of a typical potential murderer. The
FDA imposed a set of packaging regulations on manufacturers, with various types of
products being given various deadlines for compliance. The idea was that your
potential murderer could slip from the idea of Tylenol to that of aspirin in a week's
time, but it would take the expanding sphere longer to hit the brilliant idea that it
could be just any over-the-counter drug. Not just the FDA seemed to think this way;
also radio talk-show hosts seemed to love speculating about what drug might be
chosen next-but I never heard them worrying about ordinary food in grocery stores.
Yet why should it give a stochastic killer any less joy to kill by spiking ajar of
mustard than by spiking a drug? In fact, if your goal in life is to see masses of random
people die, there are all sorts of routes you can take that don't involve ingestion at all.
A friend of mine took a train from Washington to New York and en route her train
smashed into a washing machine full of rocks that had been placed on the tracks by
some do-badder. Was this part of the Tylenol-murders




Variations on a Theme as the Crux of Creativity                                        256

implicosphere in the mind of the person who did it? I doubt it, but it is possible.
       In its own gruesome way, the generalization of the Tylenol murders resembles
that of the expanding implicosphere of the Cube-and that of any idea that arises.
Ideas, whether evil or beneficial, have their own dynamics of spreading in and among
minds. Here we are primarily talking about intramind spreading (implicospheres), but
intermind spreading (infectious memes) was discussed in Chapter 3.

                                       *   *   *

         Slippage of thought is a remarkably invisible phenomenon, given its ubiquity.
People simply don't recognize how curiously selective they are in their "choice" of
what is and what is not a hinge point in how they think of an event. It all seems so
natural as to require no explanation.
         I dropped a slice of pizza on the floor of a pizza place the other evening. My
friend Don, who was less hungry than I was, immediately sympathized, saying, "Too
bad I didn't drop one of my pieces-or that you didn't drop one of mine instead of one
of yours." Sounds sensible. But why didn't he say, "Too bad the pizza isn't larger"?
His choice revealed that to his unconscious mind, it seemed sensible to switch the
role-filler in a given event, as if to imply that a pizza-slice-droppage had been in the
cards for that evening, that God had flipped a coin and, unluckily for me, it had come
out with me as the dropper instead of Don-but that it might have come out the other
way around.
         Some hypothetical replacement scenarios-I like to call them "subjunctive
instant replays"-are compelling, and come to mind by reflex. They are not idle
musings but very natural human emotional responses to a common type of
occurrence. Other subjunctive instant replays have little intuitive appeal and seem far-
fetched, although it is hard to say just why. Consider the following list:

       Too bad they didn't give us a replacement piece.
       Lucky we weren't in a really fancy restaurant.
       Too bad gravity isn't weaker, so that you could have caught it before it hit the
       ground.
       Lucky it wasn't a beaker filled with poison. Too bad it wasn't a fork. Lucky it
       wasn't a piece of good china. Too bad eating off floors isn't hygienic. Lucky
       you didn't drop the whole pizza.
       Too bad it wasn't the people at the next table who dropped their pizza. Lucky
       there was no carpet in here. Too bad you were the hungry one, rather than me.




Variations on a Theme as the Crux of Creativity                                     257

I'll leave it to you to generate other subjunctive instant replays that he might have
come up with. There is a rough rank ordering to them, in terms of plausibility of
springing to mind. It's the rhyme and reason behind that ordering that fascinates me.
Why do people find it not only plausible but even compelling to make remarks like
the following?

       If Jesse Jackson were a white man, he'd be elected President.
       If Jesse Jackson were a white man, he'd be running for dogcatcher.

These two sentences came from random voters, as quoted in Newsweek. I wonder
what slips in people's minds when they imagine a white Jesse Jackson. Do they
envision a preacher in a Baptist church? Is this person an ardent fighter for civil
rights? Or, conversely, an ardent fighter against the quota system? Similarly, what
does a high-school boy mean when he says, "If I were my father, I wouldn't lend me
the car"? Does he ever notice that if he were his father, he would ipso facto be his
own son? Or need that be so? Would the two have exchanged roles? The point is,
there are a host of questions left completely open here, yet no one balks for a second
at such counterfactuals. In fact, they are common currency, they are daily bread, they
are the meat and potatoes of communication, But some types of counterfactuals never
(or hardly ever) come up, while others, equally realityviolating, are a dime a dozen.
Daniel Kahneman and Amos Tversky, cognitive psychologists, have made studies of
how much emotion people generate upon reading stories of just-missed airplanes or
just-caught airplanes-especially ones that crash. These kinds of near misses, whether
fortunate or unfortunate, tug at our hearts and do so in nearly universal ways.
Something about these slippability examples is truly at the core of what it is to be
human and to experience the world through the filter of the human mind.
Philosophers and artificial-intelligence researchers by and large have not paid much
attention to the "catchiness" of a given counterfactual. Logicians have devoted a lot of
time and effort to trying to figure out what it would mean for a given counterfactual to
be true, but to my mind, that's not nearly as interesting-or even as meaningful-a
question as these more psychological questions:

       Which counterfactuals are likely to be triggered in a human mind by various
       types of events in the world?

       Why are some events perceived to be "near misses", while others are not?

       Why are some deaths of innocent people viewed as more tragic than other
       deaths of innocent people?




Variations on a Theme as the Crux of Creativity                                     258

At such points where deep human emotion, identification with other beings, and
perception of reality meet lies the crux of creativity-and also the crux of the most
mundane thoughts. Spinning out variations is what comes naturally to the human
mind, and is it ever fertile!




Variations on a Theme as the Crux of Creativity                                 259


                    Metafont, Metamathematics,
                   and Metaphysics: Comments on
                      .Donald Knuth's Article
                   "The Concept of a Meta-Font".
                                       August, 1982


        The Mathematization of Categories, and Metamathematics

DONALD Knuth has spent the past several years working of his system allowing him
to control many aspects of the design forthcoming books-from the typesetting and layout
down to the very shapes of the letters! Seldom has an author had anything remotely like
this power to control the final appearance of his or her work. Knuth's TEX typesetting
system has become well-known and available in many countries around the world. By
contrast, his METAFONT system for designing families of typefaces has not become as
well known or as available.
         In his article "The Concept of a Meta-Font", Knuth sets forth for the first time the
underlying philosophy of METAFONT, as well as some of its products. Not only is the
concept exciting and clearly well executed, but in my opinion the article is charmingly
written as well. However, despite my overall enthusiasm for Knuth's idea and article,
there are some points in it that I feel might be taken wrongly by many readers, and since
they are points that touch close to my deepest interests in artificial intelligence and
esthetic theory, I felt compelled to make some comments to clarify certain important
issues raised by "The Concept of a Meta-Font".
         Although his article is primarily about letterforms, not philosophy, Knuth holds
out in it a philosophically tantalizing prospect for us: that with the




Metafont, Metamathematics and Metaphysics                                                260

arrival of computers, we can now approach the vision of a unification of all typefaces.
This can be broken down into two ideas:

   (1) That underneath all 'A's there is just one grand, ultimate abstraction that can be
       captured in a finitely parametrizable computational structure-a "software
       machine" with a finite number of "tunable knobs" (we could say "degrees of
       freedom" or "parameters", if we wished to be more dignified);
   (2) That every conceivable particular `A' is just a product of this machine .with its
       knobs set at specific values.

    Beyond the world of letterforms, Knuth's vision extends to what I shall call the
mathematization of categories: the idea that any abstraction or Platonic concept can be so
captured-that is, as a software machine with a finite number of knobs. Knuth gives only a
couple of examples-those of the "meta-waltz" and the "meta-shoe"-but by implication one
can imagine a "meta-chair", a "meta-person", and so forth.
    This is perhaps carrying Knuth's vision further' than he ever intended. Indeed, I
suspect so; I doubt that Knuth believes in the feasibility of such a "mathematization of
categories" opened up by computers. Yet any imaginative reader would be likely to draw
hints of such a notion out of Knuth's article, whether Knuth intended it that way or not. It
is my purpose in this article to argue that such a vision is exceedingly unlikely to come
about, and that such' intriguingly flexible tools as metashoes, meta-fonts, modern
electronic organs (with their "oom-pah-pah" and "cha-cha-cha" rhythms and their canned
harmonic patterns), and other many-knobbed devices will only help us see more clearly
why this is so. The essential reason for this I can state in a very short way: I feel that to
fill out the full- "space" defined by a category such as "chair" or "waltz" or "face" or `A'
(see Figures 12-2, 12-3, and 12-4) is an act of infinite creativity, and that no finite entity
(inanimate mechanism or animate organism) will ever be capable of producing all
possible `A's and nothing but 'A's (the same could be said for chairs, waltzes, etc.).
    I am not making the trivial claim that, because life is finite, nobody can make an
infinite number of creations; I am making the nontrivial claim that nobody can possess
the "secret recipe" from which all the (infinitely many-) members of a category such as
`A' can in theory be generated. In fact, my Maim is that no such recipe exists. Another
way of saying this is that even if you were granted an infinite lifetime in which to draw
all the `A's you could think up, thus realizing the full potential of any recipe you had, no
matter how great it might be, you would still miss vast portions of the space of `A's.
    In metamathematical terms, this amounts to positing that any conceptual or semantic)
category is a productive set, a precise notion whose characterization is a formal
counterpart to the description in the previous paragraphs namely, a set whose elements
cannot be totally enumerated by any effective procedure without overstepping the bounds
of that set, but which can be




Metafont, Metamathematics and Metaphysics                                                 261

approximated more and. more fully by a sequence of increasingly complex effective
procedures). The existence and properties of such sets first became known as a result of
Gödel’s Incompleteness Theorem of 1931. It is certainly not my purpose here to explain
this famous result, but a short synopsis might be of help. (Some useful references are:
Chaitin, DeLong, Nagel and Newman, Rucker, and my book Gödel, Escher, Bach. )

                   An Intuitive Picture of Gödel’s Theorem
   Gödel was investigating the properties of purely formal deductive systems in the
sphere of mathematics, and he discovered that such systems-even if their ostensible
domain of discourse was limited to one topic-could be viewed as talking "in code" about
themselves. Thus a deductive system could express, in its own formal language,
statements about its own capabilities and weaknesses. In particular, System X could say
of itself through the Gödelian code:

        System X is not powerful enough to demonstrate the truth of Sentence S.

It sounds a little bit like a science-fiction robot called "ROBOT R-15" droning (of course
in a telegraphic monotone):

       ROBOT R-15-U.NFORTUNA TELY UNABLE TO COMPLETE TASK T-12-VERY
       SORRY.

Now what happens if TASK T-12 happens, by some crazy coincidence, to be not the
assembly of some strange cosmic device but merely the act• of uttering the preceding
telegraphic monotone? (I say "merely" but of course that is a bit ironic.) Then ROBOT R-
15 could get only partway through the sentence before choking: ROBOT R-15
UNFORTUNATELY UNABLE TO COMPL-.
   Now in the case of a formal system, System X, talking about its powers, suppose that
Sentence G, by an equally crazy coincidence, is the one that says,

  System X is regrettably not powerful enough to demonstrate the truth of Sentence G.

In such a case, Sentence G is seen to be an assertion of its own unprovability within
System X. In fact we do not have to rely on crazy coincidences, for Gödel showed that
given any reasonable formal system, a G-type sentence for that system actually exists.
(The only exaggeration in my English-language version of G is that in formal systems
there is no way to say "regrettably".) In formal deductive systems, this foldback takes
place of necessity by means of a Gödelian code, but in English no Gödelian code is
needed and the peculiar quality of such a loop is immediately visible.




Metafont, Metamathematics and Metaphysics                                             262

   If you think carefully about Sentence G, you will discover some amazing things.
Could Sentence G be provable in System X? If it were, then System X would contain a
proof for Sentence G, which asserts that System X contains no proof for Sentence G.
Only if System X is blatantly self-contradictory could this happen-and a formal reasoning
system that is self-contradictory is no more useful than a submarine with screen doors.
So, provided we are dealing with a consistent formal system (one with no self-
contradictions), then Sentence G is not provable inside System X. And since this is
precisely the claim of Sentence G itself, we conclude that Sentence G is true-true but
unprovable inside System. X.
   One last way to understand this curious state of affairs is afforded the reader by this
small puzzle. Choose the more accurate of the following pair of sentences:

         (1) Sentence G is true despite being unprovable.
         (2) Sentence G is. true because it is unprovable.

    You'll know you've really caught on to "Gödelism" when both versions ring equally
true to your ears, when you, flip back and forth between them, savoring that exceedingly
close approach to paradox that G affords. That's how twisted back on itself Sentence G
is!
    The main consequence of G's existence within each System X is that there are truths
unattainable within System X, no matter how powerful and flexible System X is, as long
as System X is not self-contradictory. Thus, if we look at truths as objects of desire, no
formal system can have them all; in fact, given any formal system we can produce on
demand a truth that it cannot have, and flaunt that truth in front of it with taunting cries of
"Nyah, nyah!" The set of truths has this peculiar and infuriating quality of being
uncapturable by any finite system, and worse, given any candidate system, we can use
what we know about that system to come up with a specific Gödelian truth that eludes
provability inside that system.
    By adding that truth to the given system, we come up with an enlarged and slightly
more powerful system-yet this system will be no less vulnerable to the Gödelian devilry
than its predecessor was. Imagine a dike that springs a new leak each time the proverbial
Dutch boy plugs up a hole with his finger. Even if he had an infinite number of fingers,
that leaky dike would find a spot he hadn't covered. A system that contains at least one
unprovable truth is said to be incomplete, and a system that not only contains such truths
but that cannot be rescued in any way from the fate of incompleteness is said to be
essentially incomplete. Another name for sets with this wonderfully perverse property is
productive. (For detailed coverage of the metamathematical ideas in this article, see the
book by Rogers.)
    My claim-that semantic categories are productive sets-is, to be sure, not a
mathematically provable fact, but a metaphor. This metaphor has been used by others
before me-notably, the logicians Emil Post and John Myhill-and I have written of it
myself before (see Chapter 23).




Metafont, Metamathematics and Metaphysics                                                 263

                           Completeness and Consistency
    Note that it is important to have the potential to fill out the full (infinite) space, and -
equally important not to overstep it. However, merely having infinite potential is not by
any means equivalent to filling out the full space. After all, any existing METAFONT
`A'-schema--even one having just one degree of freedom!-will obviously give us
infinitely many distinct 'A's as we sweep its knob (or knobs) from one end of the
spectrum to the other. Thus to have an 'A'-making machine with infinite variety of
potential output is not in itself difficult; the trick is to achieve completeness: to fill the
space.
    And yet, isn't it easy to fill the space? Can't one easily make a program that will
produce all possible 'A's? After all, any 'A' can be represented as a pattern of pixels (dots
that are either off or on) in an m X n matrix-hence a program that merely prints out all
possible combinations of pixels in matrices of all sizes (starting with I X 1 and moving
upwards to 2 X 1, 1 X 2, 3 X 1, 2 X 2, 1 X 3, etc., as in Georg Cantor's famous
enumeration of the rational numbers) will certainly cover any given `A' eventually. This
is quite true. So what's the catch?
    Well, unfortunately, it is hard-very hard-to write a screening program that will retain
all the `A's in the output of this pixel-pattern program, and at the same time will reject all
'K's, pictures of frogs, octopi, grandmothers, trolley cars, and precognitive photographs of
traffic accidents in the twenty-fifth century (to mention just a few of the potential outputs
of the generation program). The requirement that one must stay within the bounds of a
conceptual category could be called consistency-a constraint complementary to that of
completeness.
    In summary, what might seem desirable from a knobbed category machine is the joint
attainment of two properties-namely:

   (1) Completeness: that all true members of a category (such as the category of 'A's
      or the category of human faces) should be potentially producible eventually as
      output;
   (2) Consistency: that no false members of the category ("impostors") should ever be
      potentially producible (in short, that the set of outputs of the machine should
      coincide exactly with the set of members of the intuitive category).

   The twin requirements of consistency and completeness are metaphorical equivalents
of' well-known notions by the same names in metamathematics, denoting desirable
properties of formal systems (theorem-producing machines)-namely:

   (1) Completeness: that all true statements of a theory (such as the theory of
      numbers or the theory of sets) should be potentially producible eventually as
      theorems;




Metafont, Metamathematics and Metaphysics                                                  264

   (2) Consistency: that no false statements of the theory should ever be potentially
      producible (in short, that the set of theorems of the formal system should
      coincide exactly with the set of truths-of the informal theory).

   The import of Gödel’s Incompleteness Theorem is that these two idealized goals are
unreachable simultaneously for any "interesting" theory (where "interesting" really means
"sufficiently complex"); nonetheless, one can approach the set of truths by stages, using
increasingly powerful formal systems to make increasingly accurate approximations. The
goal of total and pure truth is, however, as unreachable by formal methods as is the speed
of light by any material object. I suggest that a parallel statement holds for any
"interesting" category (where again, "interesting" means something like "sufficiently
complex", although it is a little harder to pin down): namely, one can do no better than
approach the set of its members by stages, using increasingly powerful knobbed
machines to make increasingly accurate approximations.
   Intuition at first suggests that there is a crucial difference between the
(metamathematical) result about the nonformalizability of truth and the (metaphorical)
claim about the nonmechanizability of semantic categories; this difference would be that
the set of all truths in a mathematical domain such as set theory or number theory is
objective and eternal, whereas the set of all `A's is subjective and ephemeral. However,
on closer examination, this distinction begins to blur quite a bit. The very fact of Gödel’s
proven nonformalizability of mathematical truth casts serious doubt on the objective
nature of such truth. Just as one can find all sorts of borderline examples of 'A'-ness,
examples that make one sense the hopelessness of trying to draw the concept's exact
boundaries, so one can find all sorts of borderline mathematical statements that are
formally undecidable in standard systems and that, even to a keen mathematical intuition,
hover between truth and falsity. And it is a well-known fact that different mathematicians
hold different opinions about the truth or falsity of various famous formally undecidable
propositions (the axiom of choice in set theory is a classic example). Thus, somewhat
counterintuitively, it turns out that mathematical truth has no fixed and eternal
boundaries, either. And this suggests that perhaps my metaphor is not so much off the
mark.

                      A Misleading Claim for METAFONT
   Whatever the validity and usefulness of this metaphor, I shall now try to show some
evidence for the viewpoint that leads to it, using METAFONT as a prime example of a
"knobbed category machine". In his article, Knuth comes perilously close, in one
throwaway sentence, to suggesting that he sees METAFONT as providing us with a
mathematization of categories. I doubt he suspected that anyone would focus in on that
sentence as if it were the key sentence of the article-but as he did write it, it's fair game!
That sentence ran:




Metafont, Metamathematics and Metaphysics                                                 265

The ability to manipulate lots of parameters may be interesting and fun, but does anybody
really need a 6 1/7-point font that is one fourth of the way between Baskerville and
Helvetica?

   This rhetorical question is fraught with unspoken implications. It suggests that
METAFONT as it now stands (or in some soon-available or slightly modified version) is
ready to carry out, on demand, for any user, such an interpolation between two given
typefaces. There is something very tricky about this proposition that I suspect most
readers will not notice: it is the idea that jointly parametrizing two typefaces is no harder,
no different in principle, from just parametrizing one typeface in isolation.
   Indeed, to many readers, it would appear that Knuth already has carried out such a
joint parametrization. After all, in printing Psalm 23 (Figure 12-1) didn't he move from
an old-fashioned, compact, serifed face with relatively tall ascenders and descenders and
small x- height all the way to the other end of the spectrum: a modern-looking, extended,
sans-serif face with relatively short ascenders and descenders and large x- height? Yes, of
course-but the critical omitted point here is that these two ends of the spectrum were not
pre-existing, prespecified targets; they just happened to emerge as the extreme products
of a knobbed machine designed so that one more or less intermediate setting of its knobs
would yield a particular target typeface (Monotype Modern Extended 8A, in case you're
interested).
   In other words, this particular set of knobs was inspired solely and directly by an
attempt to parametrize one typeface (Monotype Modern). The two extremes shown in the
psalm are both variations on that single theme; the same can be said of every intermediate
stage as well. There is only one underlying theme (Monotype Modern) here, and a cluster
of several hundred variants of it, each one of which is represented by a single character.
The psalm does not represent the marriage of two unrelated families, but simply exhibits
many members of one large family.




   You can envision all the variants of Monotype Modern produced by twiddling the
knobs on this particular machine as constituting an "electron cloud" surrounding a single
"nucleus" (see Figure 12-5a). Now by contrast, joint parametrization of two pre-existent,
known typefaces (say, Baskerville and Helvetica, as Knuth suggests (see Figure 13-1)
would be like a cloud of electrons swarming around two nuclei, like a chemical bond (see
Figure 12-5c).
   In order to jointly parametrize two typefaces in METAFONT, you would need to find,
for each pair of corresponding letters (say Baskerville 'a' and Helvetica `a') a set of
discrete geometric features (line segments, serifs, extremal points, points of curvature
shift, etc.) that they share and that totally characterize them. Each such feature must be
equated with one or




Metafont, Metamathematics and Metaphysics                                                 266

FIGURE 13-1. Two typefaces of great beauty and subtlety. In (a), Baskerville; in (b),
Helvetica Light.

more parameters (knobs), so that the two letterforms are seen as produced by specific
settings of their shared set of knobs. Moreover, all intermediate settings must also yield
valid instances of the letter `a'. That is the very essence of the notion of a knobbed,
machine, and it is also the gist of the quote, of course: that we should now (or soon) be
able to interpolate between any familiar typefaces merely by knob-twiddling.
    Now I will admit that I think it is perhaps feasible-though much more difficult than
parametrizing a single typeface-to jointly parametrize two typefaces that are not radically
different. It is not trivial, to cite just one sample difficulty, to move between Baskerville's
round dot over the `i' to Helvetica's square dot-but it is certainly not inconceivable.
Conversely, it is not inconceivable to move between the elegant swash tail of the
Baskerville 'Q; and the stubby straight tail of the Helvetica `Q;-but it is certainly not
trivial.
    Moving from letter to letter and comparing them will reveal that each of these two
typefaces has features that the other totally lacks. (Incidentally, you should disregard
lowercase `g', since the `g's of our two typefaces are as different from each other as
Baskerville 'B' is from Helvetica `H'; in both cases, the two letterforms being compared
derive from entirely different underlying "Platonic essences". It is METAFONT's
purpose to mediate between' different stylistic renditions of a single "Platonic essence",
not between distinct "Platonic essences".) Presumably, in a case where one typeface
possesses some distinct feature that the other totally lacks, there is a way to fiddle with
the knobs that will make the feature nonexistent in one but present in the other. For
instance, a knob setting of zero might make some feature totally vanish. Sometimes it
will be harder to make features disappear-it might require several knobs to have
coordinated settings. Nonetheless, despite all the complex ways that Baskerville and
Helvetica




Metafont, Metamathematics and Metaphysics                                                 267

Differ. I repeat, it is conceivable that somebody with great patience and ingenuity could
jointly parametrize Helvetica and Baskerville. But the real question is this: Would such a
joint parametrization easily emerge out of two separate, independently carried-out
parametrizations of these typefaces?
   Hardly! The Baskerville knobs do not contain in them even a hint of the Helvetica
qualities-or the reverse. How can I convince you of this? Well, just imagine how great the
genius ofJohn Baskerville, an eighteenth-century Briton, would have had to be for his
design to have implicitly defined another typeface-and a typeface only discovered (or
invented) two centuries later, by Max Miedinger from Switzerland! To see this more
concretely, imagine that someone who had never seen Helvetica naively created a
METAFONT rendition of Baskerville (that is, a meta-font centered on Baskerville in the
same sense as Knuth's sample meta-font is centered on Monotype Modern). Now imagine
that someone else who does know Helvetica comes along, twiddles the knobs of this
Baskerville meta-font, and actually produces a perfect Helvetica! It would be nearly as
strange as having a marvelous music-composing program based exclusively on the style
of Dr. William Boyce (who composed in England in a baroque, elegant eighteenth-
century style) that was later discovered, totally unexpectedly, to produce many pieces
indistinguishable in style from the music of Arthur Honegger (who composed in
Switzerland in a sparse, crisp twentiethcentury style) when various melodic, harmonic,
and rhythmic parameters were twiddled. To me, this is simply inconceivable; eighteenth-
century style did not contain within it, no matter how implicitly, twentieth-century style -
whether in music or in visual arts.

           Interpolating Between an Arbitrary Pair of Typefaces
   The worst is yet to come, however. Presumably Knuth did not wish us to take his
rhetorical question in such a limited way as to imply that the numbers 6 1/7 and 1/4 were
important. Pretty obviously, they were just examples of arbitrary parameter settings.
Presumably, if METAFONT could easily give you a 6 1/7-point font that is 1/4 of the
way between Baskerville and Helvetica, it could as easily give you an 11 2/3-point font
that is 5/17 of the way between Baskerville and Helvetica-and so on. And why need it be
restricted to Baskerville and Helvetica? Surely those numbers weren't the only "soft"
parts of the rhetorical question! Common sense tells us that Helvetica and Baskerville
were also merely arbitrary choices of typeface. Thus the hidden implication is that, as
easily as one can twiddle a dial to change point size, so one can twiddle another dial (or
set of dials) and arrive at any desired typeface, be it Helvetica, Baskerville, or whatever.
Knuth might just as easily have put it this way:

   The ability to manipulate lots of parameters may be interesting and fun, but does
   anybody really need an n-point font that is x percent of the way between typeface
   TI and typeface T2 ?




Metafont, Metamathematics and Metaphysics                                               268

For instance, we might have set the four knobs to the following settings:

                        n: 36
                        x: 50 percent
                        TI: Magnificat
                        T2: Stop

   Each of these two typefaces (see Figure 13-2) is ingenious, idiosyncratic, and visually
intriguing. I challenge any reader to even imagine a blend halfway between them, let
alone draw it! And to emphasize the flexibility implied by the question, how about trying
to imagine a typeface that is (say) one third of the way between Cirkulus and Block Up?
Or one that is somewhere between Explosion and Shatter? (For these typefaces, see
Figure 13-2.)

              A Posteriori Knobs and the Frame Problem of Al
    Shatter, incidentally, provides an excellent example of the trouble with viewing
everything as coming from parameter settings. If you look carefully, . you will see that
Shatter is indeed a "variation on a theme", the theme being Helvetica Medium Italic (see
Figure 13-2). But does that imply that any meticulous parametrization of Helvetica would
automatically yield Shatter as one of its knob-settings? Of course not. That is absurd. No
one in their right mind would anticipate such a variation while parametrizing Helvetica,
just as no one in their right mind when delivering their Nobel Lecture would say, "Thank
you for awarding me my first Nobel Prize." When someone wins a Nobel Prize, they do
not immediately begin counting how many they have won. Of course, if they win two,
then a knob will spontaneously appear in most people's minds, and friends will very
likely make jokes about the next few Nobel Prizes. Before the second prize, however, the
"just-one" quality would have been an unperceived fact.
    This is closely related to a famous problem in cognitive science (the study of formal
models of mental processes, especially computer models) called the frame problem. This
knotty problem can be epitomized as follows: How do I know, when telling you I'll meet
you at 7 at the train station, that it makes no sense to tack on the proviso, "as long as no
volcano erupts along the way, burying me and my car on the way to the station", but that
it does make reasonable sense to tack on the proviso, "as long as no traffic jam holds me
up"? And of course, there are many intermediate cases between these two. The frame
problem is about the question: What variables (knobs) is it within the bounds of normalcy
to perceive? Clearly, no one can conceivably anticipate all the factors that might
somehow be relevant to a given situation; one simply blindly hopes that the species'
evolution and the individual's life experiences have added up to a suitably rich
combination to make for satisfactory behavior most of the time. There are too many
contingencies, however, to try to anticipate them all, even given the most powerful
computer. One reason for the extreme difficulty in trying to make




Metafont, Metamathematics and Metaphysics                                               269

FIGURE 13-2. A series of diverse typefaces: (a) Magn f cat; (b) Stop; (c) Cirkulus; (d)
Block




Metafont, Metamathematics and Metaphysics                                          270

machines able to learn is that we find it very hard to articulate a set of rules defining
when it makes sense and when it makes no sense to perceive a knob. It is a fascinating
task to work on making a machine capable of coaxing shy knobs out of the woodwork.
This brings us back to Shatter, seen as a variation on Helvetica. Obviously, once you've
seen such a variation, you can add a knob (or a few) to your METAFONT "Helvetica
machine", enabling Shatter to come out. (Indeed, you could add similar "Shatterizing"
knobs to your "Baskerville machine", for that matter!) But this would all be a posteriori:
after the fact. The most telling proof of the artificiality of such a scheme is, of course, that
no matter how many variations have been made on (say) Helvetica, people can still come
up with many new and unanticipated varieties, such as: Helvetica Rounded, Helvetica
Rounded Deco, Helvetican Flair, and so on (see Figure 13-3).
No matter how many new knobs-or even new families of knobs-you add to your
Helvetica machine, you will have left out some possibilities. People will forever be able
to invent novel variations on Helvetica that haven't been foreseen by a finite
parametrization, just as musicians will forever be able to devise novel ways of playing
"Begin the Beguine" that the electronic

FIGURE 13-3. Three "simple" offshoots of Helvetica: (a) Helvetica Rounded; (b)
Helvetica Rounded Deco; (c) Helvetican Flair.




Metafont, Metamathematics and Metaphysics                                                  271

organ builders haven't yet built into their elaborate repertoire of canned rhythms,
harmonies, and so forth. To be sure, the organ builders can always build in extra
possibilities after they have been revealed, but by then a creative musician will have long
since moved on to other styles. One can imagine Helvetica modified in many novel ways
inspired by various extant typefaces. I leave it to readers to try to imagine such variants.

                      A Total Unification of All Typefaces?
        The worst is still yet to come! Knuth's throwaway sentence unspokenly implies
that we should be able to interpolate any fraction of the way between any two arbitrary
typefaces. For this to be possible, any pair of typefaces would have to share the exact
same set of knobs (otherwise, how could you set each knob to an intermediate setting?).
And since all pairs of typefaces have the same set of knobs, transitivity implies that all
typefaces would have to share a single, grand, universal, all-inclusive, ultimate set of
knobs. (The argument is parallel to the following one: If any two people have the same
number of legs as each other, then leg-number is a universal constant for all people.)
        Thus we realize that Knuth's sentence casually implies the existence of a
"universal 'A'-machine"-a single METAFONT program with a finite set of parameters,
such that any combination of settings of them will yield a valid `A', and conversely, such
that any valid `A' will be yielded by some combination of settings of them. Now how can
you possibly incorporate all of the previously shown typefaces into one universal
schema?
        Or look again at the 56 capital 'A's of Figure 12-3. Can you find in them a set of
specific, quantifiable features? (For a comparable collection for each letter of the
alphabet, see the marvelous collection of alphabetical logos compiled by Kuwayama.)
Imagine trying to pinpoint a few dozen discrete features of the Magnificat 'A' (A7) and
simultaneously finding their "counterparts" in the Univers `A' (D3). Suppose you have
found enough to characterize both completely. Now remember that every intermediate
setting also must yield an `A'. This means we will have every shade of "cross" between
the two typefaces.
        This intuitive sense of a "cross" between two typefaces is common and natural,
and occurs often to typeface lovers when they encounter an unfamiliar typeface. They
may characterize the new face as a cross between two familiar typefaces ("Vivaldi is a
cross between Magnificat and Palatino Italic Swash") or else they may see it as an
exaggerated rendition of a familiar typeface ("Magnificat is Vivaldi squared") (see Figure
13-4). What degree of truth is there to such a statement? All one can really say is that
each Magnificat letter looks "sort of like" its Vivaldi counterpart, only about "twice as
fancy" or "twice as curly" or something vague along those lines. But how could a single
“curliness knob account for the mysteriously beautiful meanderings, organic and
capricious in each Magnificat Letter?




Metafont, Metamathematics and Metaphysics                                               272

FIGURE 13-4. A transition from curved to whirly to superswirly: (a) Palatino Italic
Swash caps; (b) Vivaldi caps; (c) Magnzcat caps. It is provocative to compare this figure
with Figure 16-7.

Can you imagine twisting one knob and watching thin, slithery tentacles begin to grow
out of the Palatino Italic `A', snaking outwards eventually to form the Vivaldi `A', then
continuing to twist and undulate into ever more sinuous forms, yielding the Magnificat
`A' in the end? And-who says that that is the ultimate destination? If Magnificat is
Vivaldi squared, then what is Magnificat squared?
         Specialists in computer animation have had to deal with the problem of
interpolation of different forms. For example, in a television series about evolution, there
was a sequence showing the outline of one animal form slowly transforming into another
one. But one cannot simply tell the computer, "Interpolate between this shape and that
one!" To each point in one there must be explicitly specified a corresponding point in the
other. Then one lets the computer draw some intermediate positions on one's screen, to
see if the choice works. A lot of careful "tuning" of the correspondences between figures
must be done before the interpolation looks good. There is no recipe that works in general
for interpolation. The task is deeply semantic, not cheaply syntactic.

        For a wonderful demonstration of the truth of this, look at the little book Double
Takes, in which artist Tom Hachtman has a lot of fun taking unlikely pairs of people and
combining their caricatures. His only prerequisite is that their names should splice
together amusingly. Thus he did "Bing Cosby (Bing Crosby and Bill Cosby), "Farafat"
(Farrah Fawcett-Majors and Yasir Arafat), "Marlon Monroe" (Marlon Brando and
Marilyn Monroe), and many others. The trick is to discern which features of each person
are the most characteristic and modular, and to be able to construct a new person having




Metafont, Metamathematics and Metaphysics                                               273

a subtle blend of those features, clearly enough that both contributoi1 be recognized. For
a viewer, it's almost like trying to recognize t~ parents in a baby's face.

                        The Essence of `A'-ness Is Not Geometrical

        Despite all the difficulties described above, some people, eve scrutinizing the
wide diversity of realizations of the abstract 'A'-concept maintain that they all do share a
common geometric quality. They sometimes verbalize it by saying that all `A's have "the
same shape' or are "produced from one template". Some mathematicians are inclined to
search for a topological or group-theoretical invariant. A typical suggestion might be:
"All instances of ´A' are open at the bottom and closed at the top” Well, in Figure 12-3,
sample A8 (Stop) seems to violate both of these criteria. And many others of the sample
letters violate at least one of them. In several examples, such concepts as "open" or
"closed" or "top" or "bottom" apply only with difficulty. For instance, is G7 (Sinaloa)
open at the bottom? Is F4 (Calypso) closed at the top? What about A4 (Astra)?
        The problem with the METAFONT "knobs" approach to the `A' category is that
each knob stands for the presence or absence (or size or angle, etc.) of some specifically
geometric feature of a letter: the width of its serifs, the height of its crossbar, the lowest
point on its left arm, the highest point along some extravagant curlicue, the amount of
broadening of a pen, the average slope of the ascenders, and so forth and so on. But in
many `A's, such notions are not even applicable. There may be no crossbar, or there may
be two or three or more. There may be no curlicue, or there may be a few curlicues.
        A METAFONT joint parametrization of two 'A's presumes that they share the
same features, or what might be called "loci of variability". It is a bold (and, I maintain,
absurd) assumption that one could get any `A' by filling out an eternal and fixed
questionnaire: "How wide is its crossbar? What angle do the two arms make with the
vertical? How wide are its serifs?" (and so forth). There may be no identifiable part that
plays the crossbar role, or the left-arm role; or some role may be split among two or more
parts. You can easily find examples of these phenomena among the 56 `A's in Figure 12-
3. Some other examples of what I call role splitting, role combining, role transferral, role
redundancy, role addition, and role elimination are shown in Figure 13-5. These terms
describe the ways that conceptual roles are apportioned among various geometric entities,
which are readily recognized by their connectedness and gentle curvatures.
        For a remarkable demonstration of ways to exploit these various role-
manipulations, see Scott Kim's book Inversions, in which a single written specimen, or
"gram", has more than one reading, depending on the observer's point of view. Often the
"grams" are symmetric and read the same both ways, but this is not essential: some have
two totally different




Metafont, Metamathematics and Metaphysics                                                 274

FIGURE 13-5. Examples off (a) role splitting; (b) role merging; (c) role transferal; (d)
role redundancy; (e) role addition; and (f) role elimination. The idea in all these
examples is that one smooth sweep of the pen need not fill exactly one coherent
conceptual role. It may fill two or more roles (or parts of two or more); it may fill less
than one, in which case several strokes combine to make one role; and so on. Sometimes
roles can be added or deleted without serious harm to the recognizability of the letter.
Angles, cusps, intersections, endpoints, extrema, blank areas, and separations often play
roles no less vital than those played by strokes. '

readings. The essence is imbuing a single written form with ambiguity. Both Scott and I
have for years done such drawings-dubbed "ambigrams" by a friend of mine-and a few of
my own are presented in Figure 13-6, as well as the one on the half-title page. The
strange fluidity of letterforms is brought out in a most vivid way by ambigrammatic art.
       Incidentally, it is most important that I make it clear that although I find it easier
to make my points with somewhat extreme or exotic versions of letters (as in ambigrams
or unusual typefaces), these points hold just as strongly for more conservative letters.
One simply has to look at a finer grain size, and all the same kinds of issues reappear.

                     Chauvinism versus Open-Mindedness:
                        Fixed Questionnaires versus Fluid Roles

        When I was twelve, my family was about to leave for Geneva, Switzerland for a
year, so I tried to anticipate what my school would be like. The furthest my imagination
could stretch was to envision a school that looked exactly like my one-story Californian
stucco junior high school, only with classes in French (twiddling the "language" knob)
and with the schoolbus that would pick me up each morning perhaps pink instead of
yellow (twiddling the "schoolbus color" knob). I was utterly incapable of anticipating the
vast




Metafont, Metamathematics and Metaphysics                                                275

Metafont, Metamathematics and Metaphysics   276

Metafont, Metamathematics and Metaphysics   277

difference that there actually turned out to be between the Geneva school and my
California school.
         Likewise, there are many "exobiologists" who have tried to anticipate the features
of extraterrestrial life, if it is ever detected. Many of them have made assumptions that to
others appear strikingly naive. Such assumptions have been aptly dubbed chauvinisms by
Carl Sagan. There is, for instance, "liquid chauvinism", which refers to the phase of the
medium in which the chemistry of life is presumed to take place. There is "temperature
chauvinism", which assumes that life is restricted to a temperature range not too different
from that here on the planet earth. In fact, there is planetary chauvinism-the idea that all
life must exist on the surface of a planet orbiting a certain type of star. There is carbon
chauvinism, assuming that carbon must form the keystone of the chemistry of any sort of
life. There is even speed chauvinism, assuming that there is only one "reasonable" rate
for life to proceed at. And so it goes.
         If a Londoner arrived in New York, we might find it quaint (or perhaps pathetic)
if he or she asked "Where is your Big Ben? Where are your Houses of Parliament?
Where does your Queen live? When is your teatime?" The idea that the biggest city in the
land need not be the capital, need not have a famous bell tower in it, and so on, seem
totally obvious after the fact, but to the naive tourist it can come as a surprise. (See
Chapter 24 for more on strange mappings between Great Britain and the United States.)
The point here is that when it comes to fluid semantic categories such as `A', it is equally
naive to presume that it makes sense to refer to "the crossbar" or "the top" or to any
constant feature. It is quite like expecting to find "the same spot" in any two pieces of
music by the same composer. The problem, I have found, is that most people continue to
insist that any two instances of `A' have "the same shape", even when confronted with
such pictures as Figure 12-3. Figure 12-4 helps, however, to dispel that sort of notion (as
does Figure 24-13).
         The analogy between Britain and the United States is a useful one to continue for
a moment. The role that London plays in England is certainly multifaceted, but two of its
main facets are "chief commercial city" and "capital". These two roles are played by
different cities in the U.S. On the other hand, the role that the American President plays
in the U.S. is split into pieces in Britain, part being carried by the Queen (or King), and
part by the Prime Minister. Then there is a subsidiary role played by the President's wife-
the "First Lady". Her counterpart in Britain is also split, and moreover, these days, "wife"
has to be replaced by "husband", no matter whether one considers that the "President of
England" is the Queen or the Prime Minister. (Again, see Chapter\_ 24 for much more
detail on this kind of analogy problem.)
         To think one can anticipate the complete structure of one country or language
purely on the basis of being intimately familiar with another one is presumptuous and, in
the end, preposterous. Even if you have seen




Metafont, Metamathematics and Metaphysics                                               278

dozens, you have not exhausted the potential richness and novelty in such domains. In
fact, the more instances you have seen, the more circumspect you are about making
unwarranted presumptions about unseen instances, although-a bit paradoxically-your
ability to anticipate the unanticipated (or unanticipable) certainly improves! The same
holds for instances of any letter of the alphabet or other semantic category.

                                       The `A' Spirit

        Clearly there is much more going on in typefaces than meets the eyeliterally. The
shape of a letterform is a surface manifestation of deep mental abstractions. It is
determined by conceptual considerations and balances that no finite set of merely
geometric knobs could capture. Underneath or behind each instance of `A' there lurks a
concept, a Platonic entity, a spirit. This Platonic entity is not an elegant shape such as the
Univers `A' (D3), not a template with a finite number of knobs, not a topological or
grouptheoretical invariant in some mathematical heaven, but a mental abstraction -a
different sort of beast. Each instance of the 'A' spirit reveals something new about the
spirit without ever exhausting it. The mathematization of such a spirit would be a
machine with a specific set of knobs on it, defining all its "loci of variability" for once
and for all. I have tried to show that to expect this is simply not reasonable. In fact, I
made the following claim, above:

   No matter how many new knobs-or even new families of knobs-you add to your....
   machine, you will have left out some possibilities. People will forever be able to
   invent novel variations .... that haven't been foreseen by a finite parametrization... .

        Of what, then, is such an abstract "spirit" composed? Or is it simply a mystically
elusive, noncapturable essence that defies the computationalindeed, the scientific-
approach totally? Not at all, in my opinion. I simply think that a key idea is missing in
what I have described so far. And what is this key idea? I shall first describe the key
misconception. It is to try to capture the essence of each separate concept in a separate
"knobbed machine"-that is, to isolate the various Platonic spirits. The key insight is that
those spirits overlap and mingle in a subtle way.

                Happy Roles, Unhappy Roles, and Quirk-Notes
        The way I see it, the Platonic essence lurking behind any concrete letterform is
composed of conceptual roles rather than geometric parts. (A related though not identical
notion called "functional attributes" was discussed by Barry Blesser and co-workers in
Visible Language as early as 1973.) A role, in my sense of the term, does not have a fixed
set of parameters defining the extent of its variability, but it has instead a set of




Metafont, Metamathematics and Metaphysics                                                 279

tests or criteria to be applied to candidates that might be instances of it. For a candidate to
be accepted as an instance of the role, not all the tests have to be passed; not all the
criteria have to be present. Instead, the candidate receives a score computed from the tests
and criteria, and there is a threshold point above which the role is "happy" and below
which it is "unhappy". Then below that, there is a cutoff point below which the role is
totally dissatisfied, and rejects the candidate outright.
         An example of such a role is that of "crossbar". Note that I am not saying
"crossbar in capital 'A"', but merely "crossbar". Roles are modular: they jump across
letter boundaries. The same role can exist in many different letters. This is, of course,
reminiscent of the fact that in METAFONT, a serif (or generally, any geometric feature
shared by several letters) can be covered by a single set of parameters for all letters, so
that all the letters of the typeface will alter consistently as a single knob is turned. One
difference is that my notion of "role" doesn't have the generative power that a set of
specific knobs does. From the fact that a given role is "happy" with a specific geometric
filler, one cannot deduce exactly how that filler looks. There is, of course, more to a role's
"feelings" about its filler than simply happiness or unhappiness; there are a number of
expectations about how the role should be filled, and the fulfillment (or lack thereof) can
be described in quirk-notes. Thus, quirk-notes can describe the unusual slant of a crossbar
(see Arnold Bocklin-E1 in Figure 12-3), the fact that it is filled by two strokes rather than
one (Airkraft-E3), the fact that it fails to meet (or has an unusual way of meeting) its
vertical mate (Eckmann Schrift-A2; Le Golf-F5), and many other quirks.
         These quirk-notes are characterizations of stylistic traits of a perceived letterform.
They do not contain enough information, however, to allow a full reconstruction of that
letterform, whereas a METAFONT program does contain enough information for that.
However, they do contain enough information to guide the creation of many specific
letterforms that have the given stylistic traits. All of them would be, in some sense, "in
the same style".

                                  Modularity of Roles
         The important thing is that this modularity of roles allows them to be exported to
other letters, so that a quirk-note attached to a particular role in 'A' could have relevance
to 'E', 'L', or 'T'. Thus stylistic consistency among different letters is a by-product of the
modularity of roles, just as the notion of letter-spanning parameters in METAFONT
gives rise to internal consistency of any typeface it might generate.
         Furthermore, there are connections among roles so that, for instance, the way in
which the "crossbar" role is filled in one letter could influence the way that the "post" or
"bowl" or "tail" role is filled in other letters. This is to avoid the problem of overly
simplistic mappings of one letter onto another, analogous to the Londoner asking an
American where the




Metafont, Metamathematics and Metaphysics                                                  280

American Houses of Parliament are. Just as one must interpret "Houses of Parliament"
liberally rather than literally when "translating" from England to' the U.S., so one may
have to convert "crossbar" into some other role when looking for something analogous in
the structure of a letter other than 'A', such as 'N'. In certain typefaces, the diagonal stroke
in 'N' could well be the counterpart of the crossbar in 'A'. But it is important to emphasize
that no fixed (i.e., typeface-independent) mapping of roles in 'A' onto roles in 'N' will
work; only the specific letterforms themselves (via their quirk-notes) can determine what
roles (if any) should be mapped onto each other. Such cross-letter mappings must be
mediated by a considerable degree of understanding of what functions are fulfilled by all
the roles in the, two particular letters concerned.

                   Typographical Niches and Rival Categories
         So far I have sketched very quickly a theory of "Platonic essences" or "letter
spirits" involving modular roles-roles shared among several letters. This sharing of roles
is one aspect of the overlapping and mingling that I spoke of above. There is a second
aspect, which is suggested by the phrase typographical niche. The notion is analogous to
that of "ecological niche". When, in the course of perception of a letterform, a group of
roles have been activated and have decided that they are present (whether happily or
unhappily), their joint presence constitutes evidence that one of a set of possible letters is
present. (Remember that since a role is not the property of any specific letter, its presence
does not signal that any specific letter is in view.)
         For instance, the presence of a "post" role and a "bowl" role in certain relative
positions would suggest very strongly that there is a 'b' present. Sometimes there may be
evidence for more than one letter. The eye-mind combination is not happy with any such
unstable state for long, and strains to make a decision. It is as if there is a very steep and
slippery ridge between valleys, and a ball dropped from above is very unlikely to come to
settle on top of the ridge. It will tumble to one side or the other. The valleys are the
typographical niches.
         Now, the overlapping of letters comes about because each letter is aware of its
typographical rivals, its next-door neighbors, just over the various ridges that surround its
space. The letter 'h', for instance, is acutely sensitive to the fact that it has a close rival in
'k', and vice versa (see Figure 13-7). The letter 'T' is very touchy about having its crossbar
penetrated by the post below, since even the slightest penetration is enough too destroy
its 'T'-ness and to slip it over into 'T's arch-rival niche, 't'. It's a low ridge, and for that
reason, 'T' guards it extra-carefully.




Metafont, Metamathematics and Metaphysics                                                    281

FIGURE 13-7. Have we "hen" or "ken" here? In each case, two niches in the Platonic
alphabet compete for possession of a single physical specimen. Again, the fluid way in
which minds are willing to let roles and fillers align is the source of all the trouble.

                       The Intermingling of Platonic Essences
        This image is, I hope, sufficiently strong to convey the second sense of
overlapping and intermingling of Platonic essences. "No letter is an island", one might
say. There has to be much mutual knowledge spread about among all the letters. Letters
mutually define each others' essences, and this is why an isolated structure supposedly
representing a single letter in all its glory is doomed to failure.
        A letterform-designing computer program based on the above-sketched notions of
typographical roles and niches would look very different from one that tried to be a full
"mathematization of categories". It would involve an integration of perception with
generation, and moreover an ability to generalize from a few letterforms (possibly as few
as one) to an entire typeface in the style of the first few. It would not do so infallibly; but
of course it is not reasonable to expect "infallible" performance, since stylistic
consistency is not an objectively specifiable quality.
        In other words, a computer program to design typefaces (or anything else with an
esthetic or subjective dimension) is not a conceptual impossibility;




Metafont, Metamathematics and Metaphysics                                                 282

but one should realize that, no less than a human, any such program will necessarily have
a "personal" taste-and it will almost certainly not be the same as its designer's (or
designers') taste. In fact, to the contrary, the program's taste will quite likely be full of
unanticipated surprises to its programmers (as well as to everyone else), since that taste
will emerge as an implicit and remote consequence of the interaction of a myriad features
and factors in the architecture of the program. Taste itself is not directly programmable.
Thus, although any esthetically programmed computer will be "merely doing what it was
programmed to do", its behavior will nonetheless often appear idiosyncratic and even
inscrutable to its programmers, reflecting the fact-well known to programmers-that often
one has no clear idea (and sometimes no idea at all) just what it is that one has
programmed the machine to do!

                   The Vertical and Horizontal Problems
                Two Equally Important Facets of One Problem
       I have made a broad kind of claim: that true understanding of letterforms depends
on more than understanding something about each Platonic letter in isolation; it depends
equally much on taking into account the ways that letters and their pieces are interrelated,
on the ways that letters depend on each other to define a total style. In other words, any
approach to the impossible dream of the "secret recipe" for 'A'-ness requires a
simultaneous solution to two problems, which I call the vertical and the horizontal
problems (see Figures 13-8 and 24-14).

       Vertical- What do all the items in any column have in common?
       Horizontal: What do all the items in any row have in common?

FIGURE 13-8. The vertical and horizontal problems. What do all the items in any column
have in common? What do all the items in any row have in common? Answers: Letter;
Spirit. (Compare this figure with Figure 24-14.)




Metafont, Metamathematics and Metaphysics                                                283

FIGURE 13-9. Six elegant faces created by the contemporary designer Hermann Zapf. In
(a), Optima; in (b), Palatino; in (c), Melior; in (d), Zapf Book; in (e), Zapf International;
and in (f), Zapf Chancery.

Actually, there is no reason to stop with two dimensions; the problem seems to exist at
higher degrees of abstraction. We could lay out our table of comparative typefaces more
carefully; in particular, we could make it consist of many layers stacked on top of each
other, as in a cake. On each layer would be aligned many typefaces made by a single
designer. This idea is illustrated in Figure 13-9, showing a few faces designed by
Hermann Zapf (Optima, Palatino, Melior, Zapf Book, Zapf International, and Zapf
Chancery). Along with the Zapf layer, one can imagine a Frutiger layer, a Lubalin layer,
a Goudy layer, and so on. One could try to arrange the typefaces in each layer in such a
way that "corresponding" typefaces by various designers are aligned in "shafts".
Now in this three-dimensional cake, the two earlier one-dimensional questions still apply,
but there is also a new two-dimensional question: What do all the items in a given layer
have in common? The third dimension can be explored as one moves from one layer to
another, asking what all the




Metafont, Metamathematics and Metaphysics                                                284

typefaces in a given "shaft" have in common. Moreover, a fourth dimenst can be added if
you imagine many such "layer cakes", one for each distinguishable period of
typographical design. Thus our fourth dimension, like Einstein's, corresponds to time.
Now one can ask about each layer cake: What do all the items herein have in common?
This is a three-dimensional question. Presumably, one could carry this exercise even
further.
         If we go back to the "simplest" of these questions, the original "vertical" question
applying to Figure 13-8, a naive answer to it could be stated in one word: Letter. And
likewise, a naive answer to the "horizontal" question of that figure is also statable in onl'
word: Spirit. In fact, the word "spirit" is applicable, in various senses of the term, to all
the higher-dimensional questions, such as "What do all the typefaces produced in the Art
Deco era have in common?" There is such a thing, ephemeral though it may be, as "Art
Deco spirit", just as there is undeniably such a thing as "French spirit" in music or
"impressionistic spirit" in art. (Marcia Loeb has recently designed a whole series of
typefaces in the Art Deco style, in case anyone doubts that the spirit of those times can be
captured. And then there is the book Zany Afternoons by Bruce McCall, in which the
entire spirit of several recent decades is wonderfully spoofed on all stylistic levels
simultaneously.)
         Stylistic moods permeate whole periods and cultures, and they indirectly
determine the kinds of creations-artistic, scientific, technological-that people in them
come up with. They exert gentle but definite "downward" pressures. As a consequence,
not only are the alphabets of a given period and area distinctive, but one can even
recognize "the same spirit" in such things as teapots, coffee cups, furniture, automobiles,
architecture, and so on, as Donald Bush clearly demonstrates in his book The Streamlined
Decade. One can be inspired by a given typeface to carry its ephemeral spirit over into
another alphabet, such as Greek, Hebrew, Cyrillic, or Japanese. In fact, this has been
done in many instances (see Figure 13-10). The problem I am most concerned with in my
research is whether (or rather, how) susceptibility to such a "spirit" can be implanted in a
computer program.

                                   Letter and Spirit
        These words "letter" and "spirit", of course, recall the contrast between the "letter
of the law" and the "spirit of the law", and the way in which our legal system is
constructed so that judges and juries will base their decisions on precedents. This means
that any case must be "mapped", in a remarkably fluid way, by members of a jury, onto
previous cases. It is up to the opposing lawyers, then, to be advocates of particular
mappings; to try to channel the jury members' perceptions so that one mapping dominates
over another. It is quite interesting that jury decisions are supposed to be unanimous, so
that in a metaphorical sense, a "phase transition" or "crystallization" of opinion must take
place. The decision must be solidly locked in, so that it reflects not simply a majority or
even a consensus, but a totality, a unanimity (which, etymologically, means "one-
souledness"). (For discussions of such "phase transitions of the mind", see Chapters 25
and 26, and for descriptions of




Metafont, Metamathematics and Metaphysics                                                285

FIGURE 13-10. Transalphabetic leaps by the ethereal "spirit" inherent in a given
typeface. In (a), we see the "Times" spirit jump across the gap between the Latin and
Cyrillic alphabets. In (b), the "Optima " spirit transplants itself to Greek soil. In (c), a
Hebrew spirit leaps out of the mirror and jumps into Latin clothes. Finally, in (d), a
gigantic trans-Pacific (or trans-Asiatic) leap in which a Kana spirit (Japanese syllabic
characters) jumps into Latin letters.


Metafont, Metamathematics and Metaphysics                                                  286

        In recent years there has been a spate of reported sightings of unidenl fled font-
like objects (UFO's). Many people who claim to have seen UFO 's insist that they come
from other planets. Some claim, for instance, to have seen Venusian written in the
Baskerville style, while others say they have seen Martian in the Helvetica style. There
are even claims of a complete Magn ycat-style Alphacentauribet! Often these claims are
contradictory. For instance, one witness will maintain that the bowl of the g' was cigar-
shaped, while another maintains equally vehemently that it resembled a saucer. Needless
to say, not a single such sighting has ever been scientifically validated.

computer models of perception in which a form of collective decisionmaking is carried
out, see the book by McClelland, Rumelhart, and Hinton, and my article on the Copycat
project.)
         In law, extant rules, statutes, and so on, are never enough to cover all possible
cases (reminding us once again of the fact that no fixed and rigid set of 'A'-defining rules
can anticipate all `A's). The legal system depends on the notion that people, whose
experience covers much more than the specific case and rules at hand, will bring to bear
their full range of experience not only with many categories but also with the whole
process of categorization and mapping. This allows them to transcend the specific, rigid,
limited rules, and to operate according to more fluid, imprecise, yet more powerful
principles. Or, to revert to the other vocabulary, this ability is what allows people to
transcend the letter of the law and to apply its spirit.
         It is this tension between rules and principles, tension between letter and spirit,
that
         is so admirably epitomized for us by the work of Donald Knuth and others
exploring the relationship between artistic design and mechanizability. We are entering a
very exciting and important phase of our attempts to realize the full potential of
computers, and Knuth's article points to many of the significant issues that must be
thought through, very carefully.
         In summary, then, the mathematization of categories is an elegant goal, a
wonderful beckoning mirage before us, and the computer is the obvious medium to
exploit to try to realize this goal. Donald Knuth, whether he has been pulled by a distant
mirage or by an attainable middle-range goal', has contributed immensely, in his work on
METAFONT, to our ability to deal with letterforms flexibly, and has cast the whole
problem of letters and fonts in a much clearer perspective than ever before. Readers,
however, should not pull a false message out of his article: they should not confuse the
chimera of the mathematization of categories with the quest after amore modest but still
fascinating goal. In my opinion, one of the best things METAFONT could do is to inspire
readers to chase after what Knuth has rightly termed the "intelligence" of a letter, making
use of the explicit medium of the computer to yield new insights into the elusive "spirits"
that flit about so tantalizingly, hidden just behind those lovely shapes we call "letters".




Metafont, Metamathematics and Metaphysics                                               287

Post Scriptum.
        Some months after this article appeared in Visible Language, the editor of that
journal published a most interesting commentary by Geoffrey Sampson, now a professor
in the Linguistics Department at the University of Leeds in England. Here are some
extracts from his article, giving the gist of it:

           I believe that Douglas Hofstadter is unfair in his critique of Donald Knuth's
  "Meta-font" article .... Human life involves both open-ended categories and closed
  categories, and in many cases it is very hard to say whether a given intuitively
  familiar category is open-ended or closed .... Hofstadter writes as if Knuth assumes
  an obviously open-ended category to be closed; but I cannot see that Hofstadter has
  demonstrated this .... Baskerville and Helvetica are both book faces, rather than
  faces designed exclusively for display. On the other hand, the 56 'A's of
  Hofstadter's figure [Figure 12-3] are all drawn from display faces. It is much less
  obvious that the class of book faces is open-ended than that the class of display
  faces is ....
           If we restrict the task to book faces (which are the only faces discussed by
  Knuth) then the open-endedness of the range really does become questionable.
  Hofstadter denies that this restriction affects his point: with `more conservative
  letters .... one simply has to look at a finer grain size, and all the same kinds of
  issues reappear'. Do they? ....
           The only argument Hofstadter gives for this is the difficulty of
  'parametrizing' the contrast between the round dots of Baskerville '1', j' and the
  square dots in Helvetica, and between the tails of `Q; in the two faces. But
  Hofstadter concedes that it is not `inconceivable' that these problems could be
  solved. Furthermore it seems to me that the number of such points, where two faces
  differ with respect to some property of an individual letter in a way that appears not
  to be predictable on the basis of more general differences between the faces, is
  fairly limited. The tail of `Q is an oddity in many faces; likewise the terminal of
  `G'; but on the other hand if you know what (say) `P' looks like in a given book face
  you will have a very good idea what `D' or `H' or `T' looks like.
           I would suggest that it is an entirely reasonable research programme to
  attempt to define a finite (no doubt large) set of variables (many of which would no
  doubt be very subtle) which generate all roman book faces, including faces not
  explicitly taken into consideration when formulating the variables, and excluding
  pathological letterforms .... If Hofstadter's view of typography is correct, the task
  proposed will prove to be impossible: every extra face considered will force the
  addition of yet more independent variables to the meta-font. However, I believe we
  have no adequate reason to reach this negative conclusion a priori.

       When I first read this letter, I must admit, I felt that it made sense; that I had
perhaps overstated my case. Sampson's point seemed reasonable. But then I started
wondering, `Just where are the boundary lines of




Metafont, Metamathematics and Metaphysics                                              288

'book-face-ness'?" This issue is beautifully exemplified by a tacit assumptio made by
Sampson. He calls Helvetica a book face, without any qualms. I doing so, he practically
kicks the ball between his own goal posts for m Helvetica is almost always thought of as
a display face, and is most ofte used in book titles and advertising displays. It is a sans-
serif face, lik Optima, Eras, and many others of a similar vintage. I wonder what Sampso
feels about serifed faces such as Goudy, Italia, Souvenir, Korinna, etc. (Se Figure 13-11.)
Which of these would count as, display faces, and which a book faces?
        Treacherous waters, these. The "problem" (actually not a problem at all but a
marvelous fact) is that the same typeface designers who design our favorite book faces
also design our favorite display faces. And the same sense of style and joyous creation is
called upon in both tasks. The way I

        FIGURE 13-11. Showing the futility of trying to draw a firm line between display
faces an book faces. From top to bottom, we have: Eras Demi, Romic Light, Goudy Extra
Bold, Itali Medium, Souvenir Light, and Korinna Extra Bold. It is easy to conceive of a
book being printed in any of these faces (in a light weight), yet none is a standard book
face.




Metafont, Metamathematics and Metaphysics                                               289

think of it is that each designer has a "wildness knob" with which to fiddle. When it's set
low, the complexities and trickeries "retreat" into the nooks and crannies of the
letterforms: how strokes terminate, swerve, change width, meet, and so on, and so the
resulting typeface appears reserved and dignified, conventional yet graceful and stylish,
still full of the designer's known characteristics. When wildness is set high, the desire for
unusual, exuberant effects is let out of the closet, and the resulting typeface is full of bold
flair and exciting, risky bravado: strokes are doubled, omitted, have extravagant shapes,
flourishes, and so on. It is quite naive to think that low wildness means "the same old
book-face knobs are twiddled" no matter who's doing it, whereas high wildness involves
an open-ended set of concepts.
         No creative designer with any pride would ever feel content creating within a pre-
set formula, a predetermined set of knobs. The joy of any kind of creation is in playing at
the boundaries of what has been done. Every perceptive observer has an intuitive sense of
the implicosphere centered on each standard letter and each role within it-a sense of just
how daring various deviations will seem and of just where they will begin veering off
into unacceptability. At the blurry boundaries of an implicosphere is exactly where an
artist most loves to play. With wildness set low, a designer will flirt with the boundaries
largely from within, making most decisions on the conservative side. With wildness set
high, many more risks will be taken, and the flirting will carry the designer noticeably
further from the implicosphere's center, like a satellite in a wider orbit. Norm violation is
the name of the game in creation, no matter where the "wildness" knob is set. High
wildness or low, it's still the same designer and the same creative forces expressing
themselves. It's just a question of how subtly, how subduedly, those influences will show
up.

                                          *   *   *

        Hermann Zapf is the designer of the famous sans-serif face Optima, a typeface
that some books have been printed in (see Figure 13-9a). Optima is deceptively simple-
looking. People tend to think that given one letter, they could determine all the rest easily.
Sampson says as much: "If you know what (say) 'P' looks like in a given typeface, you
will have a very good idea what 'D' or 'H' or 'T' looks like." But if that's the case, then
why did it take Zapf-one of the world's foremost type designers-seven years to design it?
All I can say is that there is rampant naiveté about the complexity of letters, even among
people who visually are otherwise very astute.
A wonderful exercise to prove this to yourself is to try to draw the Helvetica Medium `a'
by memory (see Exhibit 'a', that is, Figure 13-12a). Study it for as long as you like, and
then try to reproduce it. The better an eye you have, the more errors you will see you
have made. Try it a few times. I myself must have attempted that 'a' several dozen times,
and still I have never drawn it perfectly. This letter is one of my favorite letters of all
time,




Metafont, Metamathematics and Metaphysics                                                  290

FIGURE 13-12. Details of two classic typefaces: the 'a' of Helvetica Medium and the `g'
of Italia Book.

and I have probably spent more time admiring it than any other letter-yet for all that, I
still have not fathomed it entirely.
         The case of Helvetica is interesting. What is characteristic about it? It was one of
the first typefaces in which negative and positive spaces were given equal attention. It
employed very simple, nearly mathematical curves. Why was it designed only in 1958?
Why did it take so long for such obvious things to be done so elegantly? It's like asking
why the ancient Greeks, with their love of purity and elegance, didn't discover group
theory, the branch of mathematics dealing with abstract binary operations. Well, some
ideas are so abstract that even though they are glimpsed through a fog centuries earlier,
their full-scale arrival takes much longer. (Group theory waited patiently for 2,000 years
after the Greeks to be discovered! Isn't group theory patient with our species?) Thus it
was with the pristine qualities of Helvetica. And what seems remarkable, but is actually
to be expected, is that in the same year as Max Miedinger designed Helvetica, Adrian
Frutiger designed Univers, a lovely typeface, in many ways nearly indistinguishable from
Helvetica. Some ideas are just ripe at certain times.
         The ideas in Helvetica were not visible to anyone in the 1930's, even though
people had thousands of book faces and display faces to look at. Likewise, the ideas in
Snorple (a classic book face to be designed by Argli Snorple in 2027) are not visible to us
today, even if, in some sense, they are implicitly defined by what is all around us.
Cultural pressures, such as the development of computers and low-resolution digital
typefaces, have profound impacts on how letters are perceived. Here is a striking
example. When Hermann Zapf heard about the curve called a "super-ellipse"-an elegant
mathematical interpolation between a circle and a square (or, more generally, between an
ellipse and a rectangle), devised in the 1950's by the Danish scientist and author Piet
Hein-he decided to base a typeface on that shape. The result: Melior, a now-standard
book face whose "bowls" are super-ellipses (see Figure 13-9c). The point is, type
designers are as susceptible as anyone else is to the subtle ebb and flow of cultural waves
and evidence of those waves shows up in book faces no less than in display faces. You
just have to look more closely. Book faces pose problems no less knotty than do display
faces, Sampson notwithstanding.
         So on reconsideration, I stick with my point that all the same issues as




Metafont, Metamathematics and Metaphysics                                                291

apply to "wild" letterfdrms apply to "tame" ones-that one merely needs to look at a finer
grain size to see the same kinds of problems. As I said above, modern book faces play
with stroke tips in incredibly creative and surprising ways. Just look, for example, at
Exhibit `g'-that is, the 'g' of Italia (Figure 13-12b). Check out some of the other letters and
then see what you think of Sampson's claim.

                                          *   *    *

         People tend to think that only extreme versions of things pose deep problems.
That's why few people see modeling the creativity of, say, the trite television character of
Archie Bunker as a difficult task. It's strange and disorienting to realize that if we could
write a program that could compose Muzak or write trashy novels, we would be 99
percent of the way to mechanizing Mozart and Einstein. Even a program that could act
like a mentally retarded person would be a huge advance. The commonest mental
abilities-not the rarest ones-are still the central mystery.
         John McCarthy, one of the founders of the field of artificial intelligence, is fond
of talking of the day when we'll have "kitchen robots" to do chores for us, such as fixing
a lovely Shrimp Creole. Such a robot would, in his view, be exploitable like a slave
because it would not be conscious in the slightest. To me, this is incomprehensible.
Anything that could get along in the unpredictable kitchen world would be as worthy of
being considered conscious as would a robot that could survive for a week in the Rockies.
To me, both worlds are incredibly subtle and potentially surprise-filled. Yet I suspect that
McCarthy thinks of a kitchen as Sampson thinks of book faces: as some sort of simple
and "closed" world, in contrast to "open-ended" worlds, such as the Rockies. This is just
another example, in my opinion, of vastly underestimating the complexity of a world we
take for granted, and thus underestimating the complexity of the beings that could get
along in such a world.
         Ultimately, the only way to be convinced of these kinds of things is to try to write
a computer program to get along in a kitchen, or to generate book faces. That's when you
finally come face to face with the extremely limiting notion of what a knob really is.
People's notion of knobs has too much intuitive fluidity to it. It's hard to identify with a
computer and to see things utterly and foolishly rigidly-but that's where you have to
begin if you want to understand why knobbifying the alphabet is a task of vast
magnitude, and is a microcosm of the task of knobbifying all of human thought.

                                          *   *   *

       It is very tempting to think that a few degrees of freedom, when combined, can
cover any possible situation. After all, the number of possible states of a multi-knob
machine is the product of the numbers of settings of each of




Metafont, Metamathematics and Metaphysics                                                 292

its knobs, and multiplying a bunch of relatively small numbers together gets you rapidly
into large-number territory. A perfect illustration of this line of thought is given in an ad I
once clipped for a book called Director's and Officer's Complete Letter Book, informally
nicknamed The Ghost. Here is some of what that ad says:

           This is not a book on letter-writing technique: It is a collection of 133
   business letters already written and ready to use. They cover virtually every
   business situation you will ever meet. Just change a few words. They are arranged
   by subject, with 988 alternate phrases and sentences, keyed so that you can adapt
   the right letter to your purpose with almost no effort .... EditorJ. A. VanDuyn
   traveled for four years, collecting the finest examples of business letters written
   today. They're in crisp, direct, informal language, without cliches .... In 30 seconds
   you can look up the letter you need, by subject. You may need only to change the
   name, address, and half-a-dozen words. Or you may use one or more of the
   alternate phrases, sentences, or paragraphs on the facing page. In minutes, you've
   got your letter. With the personal touch you want. Perfectly suited to the sense you
   wish to convey .... `
           Some letters are especially hard. When you're stuck for the tactful approach,
   the just-right expression of concern, the graceful apology, you'll be thankful you
   have The Ghost. Look at, some of these subjects:

      Letters to Public Officials; Declining Appointive or Elective Positions;
      Letters of Condolence; Letters of Apology; Soliciting for Charitable
      Contributions; Adjustments-When the Answer is "No"; Letters to Creditors;
      Contacting Inactive Accounts; Collection Letters; Requests for References-11
      chapters in all.

          New subjects are thoroughly covered. You'll find letters on contracting for
   computer services, apologizing for computer errors, contracting for hardware and
   software. Virtually every letter a business executive could ever need is here in The
   Ghost-waiting for you.

I wonder if it contains letters that apologize for the mechanically written tone of recent
letters, or letters that apologize for the incorrectly selected letter sent last time-and so on.
The idea that anyone could think that every possible situation has been anticipated just
boggles the mind. How credulous does one have to be to buy this book? (By the way, if
you're interested, it costs only $49.95, and you can order it from Prentice-Hall, Inc.,
Englewood Cliffs, New Jersey 07632. But act now-it won't last long.)

                                           *   *   *

       In talking about knobs and creativity once with some architects, I encountered
some advocates of "shape grammars" used to design houses, gardens, tea rooms, and so
on. I was shown how a certain class of Frank Lloyd Wright houses known as his "prairie
houses" had been parametrized




Metafont, Metamathematics and Metaphysics                                                  293

and embedded in a- shaper grammar. An `article by H. Koning and J. Eizenberg presents
the grammar and shows a large number of external and internal designs of pseudo-Wright
houses. This kind of art by formula reminds me of the famous aleatoric waltz by Mozart,
in which one-measure fragments can be assembled in any order to make an acceptable, if
feeble, piece of music. Shape grammars recognize more levels of structure than Mozart
did, but then he was doing it only as a joke. It seems to me, after perusing several articles
on architectural shape grammars, that the designs they produce are respectable-in fact
they are very similar to the input designs. But for that very reason, they strike me as
rather dull and dry designs, given that they are all ex post facto. We are back at the issue
of pseudo-Mondrian versus genuine Mondrian (see Figure 10-14 and the accompanying
discussion), and the questionable artistic value in extracting features of a once-novel
creation and using them to allow a machine to mimic or perhaps even improve upon that
one creation, but always in a blatantly derivative way.
        Readers might be surprised to learn that one part of my research is not that distant
from either shape grammars or METAFONT: the Han Zi project, whose goal is to make a
program able to produce Chinese characters in a "twiddlable" style. All characters are
reduced to smaller units, which in turn are reduced to smaller units, and so on, until the
level of basic strokes is reached. Traditional Chinese calligraphers will tell you that there
are seven or eight such basic strokes, but that is only for humans, whose vision and
concepts are very fluid. For rigid machines, the number has to be increased. I have found
that somewhere around 40 will suffice to make just about any character, although for
most purposes 30 or 35 will do. The definition of each character is style-independent,
which means that if you change the basic strokes, all characters will change in
appearance. An example of this is shown in Figure 13-13, in which a short sentence is
printed out by Han Zi in two different styles (and in which the program says two different
things about its output).
        My co-worker David Leake and I do not harbor any illusions as to the generality
of this approach to style in Chinese. It is quite obviously subject to all the limitations of
any parameter-based approach to style: rigidity and non-creativity. Still, we find it an
exciting challenge to try to do the best we can within the obvious limitations of such a
system. It helps us see just how far these systems can be pushed, it teaches us more about
Chinese writing, and perhaps best of all, it entertains and intrigues the many. Chinese
students we know.

                                         *   *   *

        The creative, non-rut-stuck mind is always coming up with ideas that jump out of
preconceived categories. A lovely cover on Science News (January 8, 1983) shows four
new ideas for airplanes. One is a fuselage-less flying wing with six engines and with
vertical tails at both ends of the wing. Another is




Metafont, Metamathematics and Metaphysics                                                294

FIGURE 13-13. Self-descriptive Chinese sentences. The upper one, in a rather
calligraphic hand, says: "These Chinese characters I've written are really not bad. " The
lower one, in a rather robot-like hand, says: "These Chinese characters I've written are
really not good. " Both were written by the Han Zi program, with only about twenty basic
strokes changed. The basic strokes themselves are shown in the boxes. All 50, 000 (or so)
characters in the Chinese language can be built up by the Han Zi program from about 40
distinct basic strokes, so that one can switch the visual mood of any passage simply .by
switching 40 basic graphic objects. Still, we-David Leake and I-are nowhere near being
able to capture, in a few simple stroke-redefinition, the creative variety of Figure 12-4.
Our program does not see what it produces, and perception of what one has produced is
essential to good creativity.

a propeller-driven craft whose curvy propeller blades look more like flower petals than
like fan blades. The third is a' plane whose two wings bend up and over its fuselage,
meeting each other to form a complete circle (thus there is really only one wing, strictly
speaking). The fourth shows a kind of "Siamese twin" plane, with one giant wing being
shared by two parallel fuselages. Marvelous images of "Future Flight", as the caption
says. Try to put all possible future aircraft designs into a set of fixed knobs! Here is a
case where roles are split and merged with the greatest of ease. Visions of thefuture often
feature these kinds of exciting "twists" on present ideas, full of novelty and considerably
beyond trivial knob-twisting-yet even they usually fall far short of anticipating how the
future really turns out.
        An entertaining use of knobs is in the new - movie genre called "Choice-a-Rama".
The slogan says, "Where you decide what happens next!" Presumably, the audience votes
at predetermined choice points, and this selects one pathway out of a predetermined set of
possible continuations. It is like making dynamic choices at every possible turn while
driving through a city, and being surprised by where one winds up. But it must be very
expensive to have more than a few choice points, because the numbers multiply. If there
are ten binary choice points, that means 210, or,. 1,024, different pathways have to be
stored somewhere on film. It's an amazing, if decadent, symbol of our society.



Metafont, Metamathematics and Metaphysics                                              295

        In conclusion, let me mention an inspiring use of knobs: in tactical nuclear
weapons whose "yield" can be controlled. This is called, naturally enough, "dial-a-yield",
in the same spirit as "dial-a-pizza" or "dial-a-prayer" services. Depending on your need,
you can decide just how much of the enemy forces you wish to take out. A high setting
has the appealing advantage of making a bigger "kill" (although one shouldn't use crude
words like that) but the annoying disadvantage that it may trigger a similar or bigger
nuclear retaliation on the part of the enemy, thus triggering the rapid slide down a
slippery slope toward an all-out holocaust. Bother! All other things being equal, that's
undesirable, so one is encouraged to use lower settings unless one is particularly peeved
or impatient. After all, who wants to bring about Armageddon unnecessarily or
prematurely? By gosh, don't knobs have the darndest uses?




Metafont, Metamathematics and Metaphysics                                             296

